\documentclass[12pt,a4paper,oneside,openany]{book}
\usepackage[margin=25mm,headheight=15pt]{geometry}
\usepackage{fontspec}
\setmainfont{TeX Gyre Pagella}
\setsansfont{TeX Gyre Heros}
\setmonofont[Scale=0.84]{DejaVu Sans Mono}
\usepackage{amsmath,amssymb,mathtools,amsthm}
\usepackage{microtype,setspace}
\setstretch{1.06}
\usepackage{booktabs,longtable,tabularx,array}
\usepackage{graphicx,xcolor,tikz}
\usetikzlibrary{arrows.meta,positioning,calc,shapes.geometric}
\usepackage{enumitem,listings,titlesec}
\usepackage[round,authoryear]{natbib}
\usepackage{url}
\Urlmuskip=0mu plus 2mu
\definecolor{navy}{HTML}{173A53}
\definecolor{teal}{HTML}{186C75}
\definecolor{pale}{HTML}{F1F5F6}
\usepackage[unicode,colorlinks=true,linkcolor=navy,citecolor=navy,urlcolor=navy]{hyperref}
\hypersetup{pdftitle={Interpreting Regulation with Accountable Ensembles},pdfauthor={LegalMath research and engineering},pdfsubject={SFC circular interpretation, legal argumentation, ensemble assurance and verified Java delivery}}
\titleformat{\chapter}[display]{\sffamily\bfseries\color{navy}\raggedright\hyphenpenalty=10000}{\large Chapter \thechapter}{8pt}{\huge}
\titlespacing*{\chapter}{0pt}{0pt}{24pt}
\titleformat{\section}{\Large\sffamily\bfseries\color{navy}\raggedright}{\thesection}{0.7em}{}
\titleformat{\subsection}{\large\sffamily\bfseries}{\thesubsection}{0.65em}{}
\setlength{\parindent}{0pt}
\setlength{\parskip}{0.45em}
\setlength{\emergencystretch}{3em}
\setlist{itemsep=0.15em,topsep=0.35em,leftmargin=1.6em}
\setcounter{tocdepth}{1}
\makeatletter
\def\ps@legalmath{%
  \def\@oddhead{\small\sffamily\color{navy}LegalMath\enspace|\enspace Interpretation assurance\hfil\leftmark}%
  \let\@evenhead\@oddhead
  \def\@oddfoot{\hfil\small\thepage\hfil}%
  \let\@evenfoot\@oddfoot}
\makeatother
\renewcommand{\chaptermark}[1]{\markboth{Chapter \thechapter}{}}
\pagestyle{legalmath}
\lstset{basicstyle=\small\ttfamily,backgroundcolor=\color{pale},frame=single,rulecolor=\color{pale},breaklines=true,columns=fullflexible,keepspaces=true,showstringspaces=false,aboveskip=0.8em,belowskip=0.8em}
\newcommand{\True}{\mathsf{T}}
\newcommand{\False}{\mathsf{F}}
\newcommand{\Unknown}{\mathsf{U}}
\newcommand{\Conflict}{\mathsf{C}}
\newcommand{\OO}{\mathsf{OUT\_OF\_SCOPE}}
\newcommand{\IR}{\mathrm{IR}}
\newcommand{\CNL}{\mathrm{CNL}}
\newcommand{\eval}{\operatorname{eval}}
\newcommand{\srcpath}[1]{\path{#1}}
\newcolumntype{Y}{>{\raggedright\arraybackslash}X}
\newcolumntype{P}[1]{>{\raggedright\arraybackslash}p{#1}}
\newtheorem{proposition}{Proposition}[chapter]
\newtheorem{definition}{Definition}[chapter]

\begin{document}
\frontmatter
\hypersetup{pageanchor=false}
\begin{titlepage}
\vspace*{15mm}
{\sffamily\color{teal}\large RESEARCH MONOGRAPH AND IMPLEMENTATION DESIGN\par}
\vspace{12mm}
{\sffamily\bfseries\color{navy}\fontsize{32}{37}\selectfont Interpreting Regulation\\with Accountable\\Ensembles\par}
\vspace{12mm}
{\Large From SFC circulars to competing interpretations,\\bounded resolution and verified Java controls\par}
\vspace{12mm}
{\large LegalMath\par}
\vfill
\begin{minipage}{0.94\textwidth}
A technical account for compliance, legal, engineering and model-risk readers.
The running case is SFC circular 23EC46 and a retained, deliberately unapproved
implementation of its selected gift restriction. The proposed ensemble makes
agreement, disagreement, missing evidence and the limits of automatic resolution
visible before a control reaches a bank system.
\end{minipage}
\vspace{16mm}
\hrule
\vspace{5mm}
22 September 2026\hfill Complete author draft, version 1.0\par
\small Public-source and synthetic examples. Independent compliance, technical
and reader acceptance remain pending. Proposed extensions are distinguished
from the currently implemented workbench throughout.
\end{titlepage}
\hypersetup{pageanchor=true}
\chapter*{The decision this book supports}
\addcontentsline{toc}{chapter}{The decision this book supports}
A private bank cannot treat a software test result as a complete answer to a
regulatory question. A program may reproduce a decision table exactly while the
table overlooks an exception, applies to the wrong activity, or relies on a term
whose meaning has never been established. The practical question is how to make
those mistakes harder to introduce, easier to discover, and impossible to hide
inside a confident-looking automated answer.

This book develops an interpretation workbench that retains several candidate
readings, investigates the reasons for and against them, checks their consequences,
and keeps unresolved questions attached to the control. Multiple methods provide
different opportunities to discover error. A discrepancy opens a recorded
resolution process with finite limits. Reaching those limits produces a report
of the remaining uncertainty and prevents release of an affected control; it does
not cause the system to choose the most popular answer.

The intended reader need not know a formal legal language, an argumentation
semantics or a tree-search algorithm. Those methods are developed through a
single circular and deliberately contrasting client scenarios. Engineers should
be able to derive data structures, state transitions and acceptance tests from
the text. Compliance readers should be able to challenge the source mappings,
scope, classifications and review decisions without reading the software.

There are two distinct contributions. The first is an explanation of what the
literature and the existing LegalMath evidence establish. The second is a new
engineering design that combines those methods. The design is a proposal for
implementation and evaluation; published results about another task do not
establish its performance on SFC circulars. Finite search also cannot establish
that every defensible legal reading has been found. The design instead makes
coverage obligations, source dependencies, disagreement and stopping reasons
explicit, and assigns responsibility for the decisions automation cannot settle.

The ten chapters follow the order in which a bank would have to justify a
control. The initial circular exercise exposes the gap between execution and
interpretation. Formal meaning and relevant languages then make the alternatives
precise. Legal argumentation and search preserve competing readings. The central
ensemble chapter specifies their bounded investigation. Verification, Java
integration, service contracts, evaluation and continuing operation complete the
route from a source document to an attributable bank decision.

The companion machine-readable contracts describe a proposed extension, not a
silent alteration of the accepted runtime. The existing oracle and source
snapshots are preserved. A successful document build establishes that this
manuscript compiles; it does not supply the independent legal judgment, empirical
validation or reader feedback required to adopt the product.
\tableofcontents
\mainmatter
\chapter{A correct program can enforce the wrong rule}
\label{ch:problem}

\section{The decision before the decision}
\label{sec:problem-decision}
A distributor proposes a marketing offer involving an investment fund. The
offer includes an incentive. A compliance control must determine whether the
incentive falls within a restriction, whether an exception applies, and whether
the relevant evidence is sufficient to make that determination. A software team
can make these questions look simple: gather a few fields, evaluate a condition,
and return a result to the bank's Java system. The difficult decisions have
already occurred before the condition runs. Someone has decided what counts as
an incentive component, what the words in the circular cover, what evidence
establishes a promotional link, and which version of a referenced rule applies.

The distinction is consequential. Suppose a gift restriction covers both a
named investment product and a type of investment product. An implementation
that checks only the first route may work perfectly for a demonstration using
one named fund. Its unit tests may be comprehensive with respect to the code
that was written. Its developer may correctly explain every branch. The defect
is that a branch required by the adopted interpretation never reached the code.
Testing generated branches alone will not reveal a missing source obligation.

Now suppose that both promotional routes are represented, but an exception for
a fee discount is applied to an entire offer containing a fee waiver and a
separate voucher. The formal expression can be well typed, consistent and
executable. Its output can be reproduced by Python, Java and a solver. The
unresolved question is the object to which the exception attaches. A Boolean
variable named \texttt{discount} conceals that question unless its definition
identifies the assessed component and the evidence supporting the classification.

These are hypothetical failure scenarios derived from the retained 23EC46
exercise, not findings about a deployed bank system. They show how interpretation error can become a material operational risk. Regulatory action,
reputational damage and accountability for an erroneous control make the quality
of the decision process important. This book does not estimate the probability
or financial severity of an enforcement outcome. It specifies the evidence a
bank would need to understand, challenge and authorize its own interpretation.

The system we want to build therefore begins one step earlier than conventional
rule implementation. It asks what the bank has grounds to encode, which
alternative readings have been considered, what remains unknown, and who has
authority to resolve the uncertainty. Code generation follows that record. It
does not substitute for it.

The unit of work is a \emph{regulatory change assessment}. It includes source
material, a declared business scope, proposed interpretations, affected controls,
data requirements, implementation changes and review decisions. A single
circular can create several such assessments because different provisions may
concern different actors or activities. Conversely, one assessment may need
several circulars, Code provisions, FAQs and earlier decisions. Treating one
downloaded page as a self-contained legal specification would lose these
relationships.

\section{What the existing experiment actually established}
\label{sec:problem-evidence}
The present LegalMath workbench provides a useful starting point because it has
already crossed the engineering boundary from a formal rule to generated Java.
Its second-circular exercise concerns a deliberately narrow interpretation of
paragraph 10 of SFC circular 23EC46, dated 24 October 2023
\citep{sfcmarketing}. The example assesses one incentive component in a specified
distributor and SFC-authorised-fund profile. Whether the component is a gift,
whether it is a fee discount, and whether it has a relevant product link are
supplied classifications. The program does not establish those classifications
by reading an offer document.

The exercise preserved a packet of 22 named expected outcomes before constructing
the candidate rule. A separate truth-table implementation enumerated 729
combinations of true, false and unknown values across six inputs. The actual
generated Java class and the Python evaluator agreed on the declared results.
Four deliberately altered rules were compiled and detected by the named cases.
Those are substantial engineering observations: they show that the exercised
implementation computes the proposed rule in the checked domain, and that the
cases can expose several particular errors.

They do not provide independent evidence that the proposed rule is the correct
reading of the circular. The same implementing agent authored the interpretation
and the expected outcomes. Writing one before the other makes later alteration
detectable; it does not create a second legal judgment. The separate truth-table
oracle removes a direct dependency on the candidate's abstract syntax tree. It
still expresses the same chosen formula. If that choice omitted a legal
condition, all three implementations could agree on the omission.

The candidate also retained five unresolved review issues and two unresolved
source dependencies. Its attempt to enter review was rejected. It remained a
draft and no release was created. That refusal is part of the useful result.
The workbench did not mistake successful compilation for authority to deploy a
regulatory control. The monograph preserves this example as a record of an
unapproved candidate; resolving the cited dependencies here by assumption would
destroy the very distinction the example teaches.

The broader application acceptance contains 154 passing tests, but its drafting
provider returns supplied fixtures. It has not interpreted arbitrary circulars
in a measured live-model evaluation. The acceptance count is evidence about
source intake, data contracts, execution, review workflow and delivery under the
tested conditions. It is not an estimate of translation accuracy. A dashboard
that presents the count without the target of the tests would invite a stronger
conclusion than the evidence supports.

This distinction suggests the question the new product must answer. Given a
source packet and a proposed operational scope, can it expose a material
misinterpretation before that misinterpretation becomes an approved executable
control? Success requires more than generating valid syntax. It requires finding
missing alternatives, retrieving the context that changes a reading, and making
disagreements concrete enough for a competent reviewer to resolve.

\section{Five places where meaning can be lost}
\label{sec:problem-boundaries}
It is tempting to describe the workflow as one translation from English to
Java. That description hides several different kinds of error. The first is
source error: the reviewed text is incomplete, extracted incorrectly, superseded,
or wrongly treated as authoritative. An omitted footnote can remove a definition.
A copied paragraph can lose a heading that limits its scope. A source's
publication date can be mistaken for its effective date. None of these failures
requires a language model to reason incorrectly; the model may never receive
the material needed for the correct decision.

The second is interpretation error. The words are present, but their legal
function is misrepresented. A qualification becomes an unconditional rule. An
exception is attached to the wrong object. A statement about possible
classification becomes a finding of fact. A normative expectation is converted
into an operational block without recording the bank's choice. These are
decisions about meaning, not merely extraction or programming defects.

The third is specification error. The reviewer may have a defensible
interpretation, but the intermediate rule expresses something else. A conjunction
is replaced by a disjunction, a threshold loses equality, or an unknown condition
is interpreted as false. The specification can also fail through insufficient
expressiveness: a time-dependent obligation may be forced into a one-off Boolean
test. A software language that cannot represent a required distinction must
report the gap rather than approximate it silently.

The fourth is execution error. Java can evaluate a correct intermediate rule
incorrectly because of a compiler defect, numeric conversion, time-zone mismatch,
or inconsistent evidence handling. A data adapter may supply dollars where the
rule expects cents. A build may package an older runtime. An apparent match on a
few simple examples can miss a defect in unknown-value or conflict precedence.
These failures are often amenable to strong formal or exhaustive checks once
the specification and input domain are fixed.

The fifth is operational error. The correct Java result can be used for the
wrong client, product, source version or transaction time. Consent can change
between evaluation and order commitment. A system can treat the absence of one
prohibition as permission to proceed despite unresolved checks elsewhere. A
reviewed rule can remain active after its source is replaced. These are failures
of integration and change control, even when the individual rule evaluator is
correct.

Figure~\ref{fig:assurance-chain} gives each transition its own question. The
arrows do not imply that success flows automatically from one stage to the next. Every
arrow needs evidence appropriate to the claim it makes.

\begin{figure}[htbp]
\centering
\begin{tikzpicture}[node distance=7mm, every node/.style={font=\small},
box/.style={draw=navy,rounded corners,align=center,text width=125mm,minimum height=11mm},
arr/.style={-{Latex},draw=teal,thick}]
\node[box] (s) {Retained source: are these the complete, applicable words?};
\node[box,below=of s] (i) {Interpretation: what reading do the words and context support?};
\node[box,below=of i] (f) {Formal specification: does the rule express that reading?};
\node[box,below=of f] (j) {Java execution: does this binary preserve the rule?};
\node[box,below=of j] (a) {Bank action: is this result used with the right evidence and authority?};
\draw[arr] (s)--(i); \draw[arr] (i)--(f); \draw[arr] (f)--(j); \draw[arr] (j)--(a);
\end{tikzpicture}
\caption{A correct downstream step does not repair a mistaken upstream premise.
Each transition has a different assurance question.}
\label{fig:assurance-chain}
\end{figure}

A useful consequence follows immediately. Two checks count as meaningful
redundancy only to the extent that their opportunities for failure differ. A
second compiler can challenge code generation. It cannot detect a missing source
document if it receives the same incomplete specification. A second language
model may propose a different reading, but it cannot authenticate a source merely
by describing it confidently. Redundancy must be designed around the failure
being investigated.

\section{The denominator problem in completeness}
\label{sec:problem-denominator}
Suppose a drafting agent extracts eight requirements from a circular and produces
eight rules. A coverage report then shows that every extracted requirement has a
rule, every rule has a test, and every test passes. The result sounds complete.
It is complete only relative to the agent's own extraction. If the circular
contains a ninth operative requirement, the report has no row on which that
omission can appear. The system has used its answer to define the denominator
of its completeness measure.

The source inventory must therefore originate independently of the proposed
rule set. At the document level it records paragraphs, headings, annexes,
footnotes, tables and incorporated references. At the normative level it records
potential duties, prohibitions, permissions, definitions, exceptions and
qualifications. These two inventories serve different purposes. A paragraph
inventory can establish that a passage was considered, but one paragraph may
contain several operative propositions. A rule inventory can establish that an
identified proposition has a disposition, but it cannot prove that the
propositions were all identified.

For the retained circular, the existing inventory contains 18 numbered
paragraphs and four footnotes. Only one paragraph has an executable control in
the second-circular example. Other passages concern context, other controls,
qualitative assessment or external dependencies. It would be wrong to describe
the circular as implemented because its selected gift rule passes 729 input
combinations. It would also be wrong to implement every paragraph mechanically:
an observation about market practices is not necessarily a new obligation, and a
qualitative duty may need a reviewed procedure rather than a fabricated numeric
threshold.

A complete disposition process gives every inventoried proposition a reasoned
destination. It may be linked to an executable rule, to a reviewed manual
procedure, to a justified exclusion from the declared assessment, or to an open
question. An open question is visible incompleteness. An exclusion requires a
scope rationale and reviewer authority. A manual procedure needs an owner and
evidence of performance; attaching the word ``manual'' to an inconvenient clause
does not implement it.

The same denominator problem occurs in search. If five candidate interpretations
agree, the correct conclusion is agreement among those five candidates under
their recorded assumptions. The system has not measured the interpretations it
failed to generate. Breadth constraints and targeted challenges improve the
opportunity to find missing readings, but there is no finite list of prompts
that proves unrestricted legal interpretation has been exhausted. The product
must report which dimensions were explored and which remain outside the search.

This makes completeness a collection of bounded claims. We can check that every
retained paragraph has a disposition. We can check that every cited reference has
been resolved or explicitly left open. We can check that each represented
exception has positive, negative and uncertain cases. We cannot infer from those
facts that no relevant circular exists elsewhere or that every legally
defensible reading has been found. A bank needs a process for setting and
maintaining its regulatory perimeter as well as a translator inside that perimeter.

\section{Why agreement is useful and insufficient}
\label{sec:problem-agreement}
An ensemble is a group of methods whose results are examined together. In a
prediction task it is common to combine answers by majority vote or averaging.
For regulatory interpretation, the useful unit of combination is often a
reasoned claim: one method identifies a scope restriction, another locates an
exception, and another generates a case that exposes a missing promotional route.
Their outputs need not be identical to contribute to the same assessment.

Agreement can provide evidence that a representation is stable across some
transformations. A controlled-English sentence and its formal rule may produce
the same decisions. A back-translation may preserve the formal structure. A
second implementation may agree with the first. Each observation rules out some
failure explanations. None rules out every explanation. In particular, methods
that share the same omitted premise can agree for the same wrong reason.

Disagreement is therefore not merely a nuisance to suppress. It may be the first
observable sign of a missing definition, a conflict in evidence, or a genuine
interpretive choice. If one candidate treats a product-type promotion as covered
and another does not, the difference points to a specific source phrase. If two
candidates disagree only when a classification is unknown, the problem may be
their evaluation semantics. If two sources have different effective intervals,
the disagreement may be temporal rather than substantive. The controller should
classify these possibilities before attempting a repair.

The converse is equally important: disagreement alone does not make every
candidate legally respectable. A generated rule that deletes an express exception
may be a useful negative control but a poor interpretation. The system should
retain it long enough to explain the rejection and test the detectors, while
distinguishing it from a reading that remains genuinely unresolved. A shortlist
filled with contrived alternatives can increase reviewer workload without
improving assurance.

The proposed ensemble consequently has two obligations. It must preserve
material dissent until there is a recorded reason to resolve it, and it must
identify the evidentiary standing of that dissent. An objection supported by an
applicable source is different from an unsupported model suggestion. A solver
counterexample proves a difference in formal consequences; it does not by itself
give one interpretation greater legal authority. A human reviewer can accept a
reading for a declared scope, but that decision must not be mislabelled as a
machine proof of English meaning.

Chapter~\ref{ch:ensemble} will specify how these distinctions become a finite
automatic process. The core idea is that a discrepancy creates work with a
defined purpose: retrieve missing authority, correct an encoding, investigate an
alternative or clarify a fact. Repetition without a new diagnostic objective is
not additional assurance. When the permitted actions have been exhausted, the
system must make the remaining issue easier for a reviewer to understand.

\section{The object being classified is part of the rule}
\label{sec:problem-objects}
The incentive example reveals a general weakness of compact decision tables:
they often define a field without defining the entity to which it belongs. A
field called \texttt{is\_discount} could describe a payment, a fee line, a
marketing offer, a product or a client relationship. These objects are not
interchangeable. If a fee waiver and voucher are separate incentive components,
a property of the waiver cannot automatically be transferred to the voucher.
The rule's subject identity is therefore part of its meaning.

The source-to-rule specification must state the unit of assessment before
introducing its predicates. In this example it is one proposed component within
a documented offer. The specification must also state how components are
identified and when a purported separation is legally meaningful. A data
designer cannot settle that second question solely by assigning two database
identifiers. Operational decomposition needs to reflect the reviewed legal
analysis; otherwise a bank could accidentally change the meaning of an exception
by changing how it stores an offer.

The same problem appears in private-bank suitability workflows. A statement
about a client's status does not necessarily apply to every account, every
product category or every transaction channel. Consent may be linked to an
identified service and period. A relationship manager's assessment may have a
subject, scope and evidence date that differ from those of a wealth record.
The intermediate representation must retain these bindings, not merely a
collection of client-level Boolean fields.

A useful specification sentence therefore contains both a predicate and its
subject: ``component C in offer O is classified as a fee discount under
definition D, on the evidence recorded in assessment A.'' That sentence is longer
than ``discount = true,'' but the extra information explains what the Boolean
means. The eventual runtime can carry compact identifiers so that ordinary
transactions remain efficient. The review record retains the definitions and
evidence those identifiers name.

This is also why a language model's ability to extract entities is insufficient.
It may correctly identify the words ``fee waiver'' and ``voucher'' while failing
to determine the legal relationship between them. Extraction can propose
candidates for classification. The system must separately record the reasoning
that accepts their relationship and the uncertainty that survives it. That
separation lets a later source clarification update a definition without
pretending that the underlying marketing document changed.

\section{Time changes the question being answered}
\label{sec:problem-time}
An apparently straightforward rule can yield different legitimate assessments
at different times. The regulation may change, the facts may change, or the
bank may learn something new about an earlier event. These are three distinct
sources of variation. A source publication date identifies when a document was
issued. An effective interval identifies when a reviewed rule is intended to
apply. An evidence recording time identifies when the bank had the relevant
information. Substituting one date for another changes the assessment.

Consider a gift classification entered on Tuesday for an offer made on Monday.
A historical review can ask what the bank should conclude now about Monday's
offer, using Tuesday's evidence. It can also ask what the system could conclude
on Monday from the information then recorded. Both questions are useful. They
require different knowledge cutoffs. Replaying the earlier decision with all
subsequently available facts would conceal the information gap that existed at
the time.

The current runtime makes assessment time and knowledge cutoff explicit through
\texttt{valid\_at} and \texttt{known\_at}. Known evidence has a validity interval
and a recording timestamp. An expired classification becomes unknown for an
assessment outside its interval. Evidence recorded after the historical cutoff
cannot be used in that historical assessment. A superseding record names the
record it corrects; simple arrival order does not determine which assertion is
authoritative.

These are engineering conventions selected for a precise replay question. They
do not determine the legal effect of every late discovery or correction. A
corrected factual record may establish that an obligation was performed on time,
while the bank's earlier inability to demonstrate that fact remains part of the
audit history. A regulatory amendment may apply prospectively, retrospectively or
with transition rules that need interpretation. The translator must represent
the relevant temporal question before the runtime can answer it.

The ensemble must therefore challenge time assumptions as well as the visible
substantive condition. One candidate can be correct under a historical source
version and wrong under the current one. Another can use a current definition
to interpret an earlier transaction without justification. A disagreement report
that omits the source revision and assessment times can make these distinct
questions look like a conflict about the same rule. Immutable context identifiers
are necessary to know what is actually being compared.

\section{From an answer to an accountable decision}
\label{sec:problem-accountability}
A useful product result contains more than a sentence announcing compliance or
non-compliance. It identifies the precise condition evaluated, the reviewed scope,
the facts and definitions used, the source version, the proposed interpretation,
the evidence of correct execution, and any unresolved dependency. These details
serve different readers. A compliance reviewer needs to see the interpretation.
A data owner needs to see the evidence mapping. An engineer needs reproducible
inputs and an artifact identity. A release owner needs to know which approvals
bind the exact binary being deployed.

The design should not force every reader to inspect every detail on every case.
The review interface can begin with the decision and the material reason, then
allow inspection of assumptions and counterexamples. The underlying record must
still support that inspection. A concise explanation is valuable when it is a
faithful view of a fuller record; it is dangerous when it replaces a missing
record with persuasive prose.

Accountability also requires distinguishing an interpretation from a bank policy.
A bank may decide to prohibit a class of offers while a regulatory question
remains unresolved. That can be a sensible operational choice under its own
authority. The system should record the policy owner, scope and reason for the
restriction. It should not attribute the restriction to an SFC sentence unless
the interpretation supports that attribution. Otherwise later reviewers cannot
tell whether a change requires new legal analysis or a policy decision by the
bank.

Likewise, a manual review route must have a defined effect. ``Refer to compliance''
is incomplete if no owner, service expectation, required evidence or transaction
state is specified. An unresolved assessment may hold an automated approval,
route the case for review or invoke an already approved conservative policy. The
choice belongs to the bank's operating design. The translator should never
invent a transaction permission merely because its retries have reached a limit.

The product's ambition is thus substantial but specific. It should improve the
quality of the questions reaching a human reviewer, preserve alternatives that
would otherwise disappear, and verify execution of the interpretation that is
eventually adopted. The next chapter follows the second circular through that
entire route. Its apparent simplicity is useful: a small formula is sufficient
to expose why independent source judgment, explicit scope and uncertainty cannot
be replaced by a large number of successful program evaluations.

\chapter{The second circular, from source to an unreleased control}
\label{ch:circular}

\section{The source packet and the question it supports}
\label{sec:circular-packet}
The running example begins with an identifiable public document: SFC circular
23EC46, issued on 24 October 2023, concerning additional returns and other
services or arrangements offered when marketing SFC-authorised funds
\citep{sfcmarketing}. The retained experiment selected its paragraph-10 gift
restriction. The circular as a whole discusses more than that restriction. A
reader needs to know both what was selected and what remains outside the control
before a result from the control can be interpreted.

The experiment preserved the official source response, an extracted text
derivative and anchors connecting the selected rule to its source passages.
On 22 September 2026, a fresh retrieval matched the retained raw bytes. This
check establishes that the retained response matched what the retrieval route
returned at that time. It does not establish the complete current legal position
or resolve the effective versions of incorporated sources. A hash identifies
bytes; it does not by itself establish their authority, completeness or legal
effect.

The selected assessment concerns one proposed incentive component, offered by a
distributor in promoting a fund within the declared profile. It is a question
about a specific restriction. It is not a whole-offer compliance verdict. The
distinction is visible in the result vocabulary. A true result means that the
selected prohibition trigger is established. A false result means that this
trigger is not established as true on the supplied, resolved facts under the
specified semantics. It supplies no general permission to market the offer.

The retained packet explicitly leaves the applicable Code paragraph and cited
FAQ unresolved. This matters twice. Those sources may settle the meaning or
scope of a term; they may also establish that the selected profile is narrower
than the underlying requirement. Their absence is therefore a release-blocking
dependency, not a bibliographic inconvenience. The example is useful precisely
because it continues to expose that absence after its executable tests pass.

The book uses this historical, selected packet to teach implementation and
assurance. New commercial offers, cashback definitions and mixed-component
scenarios introduced below are synthetic investigations. They are not reported
SFC rulings. The unresolved references are not silently filled using a language
model's recollection of a Code provision. Doing so would convert a visible gap
into an untraceable premise.

\section{Reading the surrounding circular before isolating the rule}
\label{sec:circular-context}
A translation that begins by copying the most convenient sentence can miss why
that sentence is there. The circular's opening establishes its distribution
context. The next passages describe observed practices, including additional
features, standing investment arrangements and marketing. Those observations
help identify the activity being discussed, but they are not automatically
individual executable prohibitions. A source inventory should retain their
function without manufacturing a rule for every paragraph.

The discussion then turns to additional returns and their possible regulatory
consequences. Some questions concern the structure of an arrangement and public
promotion. Others concern impressions created about the underlying fund. The
gift discussion appears within that broader context. Consequently, a finding
that one separate classification is absent does not clear the gift issue. The
gift control is not conditional on first finding that every other regulatory
concern applies.

Later passages address redemption arrangements, administrative cutoffs,
advertising and the description of services. Some concerns require event or
workflow evidence; others require qualitative assessment of communication.
Those mechanisms differ from the paragraph-10 Boolean control. An adequate
implementation of the whole circular would need to preserve those differences.
It cannot simply accumulate a larger expression around the gift predicate and
call the circular complete.

The existing disposition inventory has 18 numbered paragraphs and four
footnotes. Its assignments are agent-authored and unreviewed. They include
context, separate controls, qualitative review and unresolved dependencies.
Paragraph 10 is the only numbered passage with an executable control in this
exercise. A future independent reviewer should challenge both the assignments
and the completeness of the inventory. Merely approving the final formula would
leave the earlier decision about what to implement unexamined.

This provides an immediate requirement for the ensemble. One path must inspect
the source independently of the candidates and ask what obligations, exceptions
and dependencies need a disposition. Another path may develop the selected
rule. Comparing those paths creates a place for an omitted requirement to
appear. If both start from the same already-filtered list of paragraphs, the
apparent redundancy begins too late.

\section{Naming the six inputs precisely}
\label{sec:circular-inputs}
The executable example has two scope inputs and four body inputs. Their names
are convenient for software, but their definitions do the substantive work.
The first scope input records whether the assessed actor is a distributor in
the declared profile. The second records whether the activity concerns promotion
of an SFC-authorised fund within that profile. Both are reviewed classifications.
The prototype does not infer them from the actor's name or the presence of the
word ``fund'' in a document.

The first body input records whether the proposed component is classified as a
gift. The second and third record a promotional link to a specific product and
to a particular product type. A link to funds from a particular fund house
requires an explicit classification where the relevant example is used; the
engine does not assume that every mention of a fund house creates such a link.
The fourth body input records whether the entire assessed component falls
within the fee-or-charge-discount exception.

The word ``entire'' in that final definition is an engineering expression of
the proposed component-level reading. It prevents an assessed voucher from
inheriting the classification of a separate fee waiver. It does not prove that
the decomposition is legally correct. A reviewer must establish the relevant
unit of assessment from the source and context. The definition and the review
issue must travel together until that question is resolved.

For a compact mathematical account, let $A$ denote established actor scope and
$S$ the fund/activity scope. Let $G$ denote the gift classification, $P$ the
specific-product link, $T$ the product-type link and $D$ the discount exception.
For now suppose these are known Boolean values. The body condition is
\begin{equation}
 B=G\land(P\lor T)\land\neg D.
 \label{eq:gift-body}
\end{equation}
It asks three things together: is the component a gift, does either promotional
route hold, and is the exception absent? The two alternatives occur inside one
parenthesized condition. Replacing that disjunction by a test of $P$ alone
changes the rule for offers linked only to a product category.

Scope is assessed separately. When both scope inputs are true, the body answers
the selected question. When a scope input is known false, the limited profile
returns \texttt{OUT\_OF\_SCOPE}. When scope is unresolved, it returns
\texttt{UNKNOWN}, even if a body condition might otherwise be false. The
distinction prevents an unresolved applicability question from being concealed
inside an apparently definitive in-scope result.

The current runtime also has a conservative conflict policy. It checks the
static dependency closure of the requested rule before evaluating branches.
Conflicting referenced evidence produces \texttt{CONFLICT}, including when a
later scope branch would otherwise make the body irrelevant. This policy is
part of the declared engineering semantics. It should not be attributed to the
wording of the SFC circular. Chapter~\ref{ch:proof} explains why such a policy
must be fixed before comparing implementations.

\section{The named cases are an argument about the rule}
\label{sec:circular-cases}
The 22 named cases are most useful when read as a sequence of challenges to an
interpretation. They do not merely fill rows in a test fixture. Each changes a
fact or evidentiary condition that a plausible implementation might mishandle.
The baseline case is a component classified as a gift, linked to one specific
fund, with no fee-discount exception, and known in-scope actor and activity.

\subsection{The alternative promotional route}
Case g01 establishes the baseline prohibition trigger. Under
equation~\eqref{eq:gift-body}, $G$ and $P$ are true and $D$ is false, so the
body is true. This is a necessary example but a weak test on its own. An
incorrect expression $G\land P$ would also return true, even though it omits
both an exception and an alternative route.

Case g02 changes the promotional facts: no individual fund is named, but the
component is linked to a fund category. Now $P$ is false and $T$ is true.
The disjunction remains true. A candidate that omitted the product-type route
would return false, making this case a concrete witness to the omission. The
source's inclusion of that route supplies the interpretive objection; the
different program outputs make its consequence inspectable.

Case g15 uses the fund-house example through a supplied reviewed product-type
classification. Its purpose is not to write a text classifier for fund-house
mentions. It checks that the chosen formal vocabulary has a place to represent
the classification when review establishes it. This distinction prevents a
source example from being silently generalized into a universal inference from
an entity name.

Case g09 establishes a specific-product link while leaving the product-type
link unknown. The known route is sufficient for the disjunction. Unknown
information about a second route does not erase the first. Case g14 reverses
the situation: the specific-product route is false and the remaining route is
unknown. Here the unresolved fact can change the outcome, so the correct result
under the declared semantics is unknown.

Taken together, these cases explain the meaning of ``either'' in an executable
rule with partial evidence. The system must neither require every route to be
known true nor convert every missing route to false. This is a semantic
distinction, not a preference about how verbose the explanation should be.

\subsection{The exception and the object it qualifies}
Case g03 changes the baseline component to one classified as solely a fee or
charge discount. The exception settles the selected body condition as false.
The result says that this particular gift restriction is not triggered under
the supplied classification. It does not resolve advertising, authorization or
any other requirement associated with the offer.

Case g04 returns to a voucher while another component of the same offer is a
fee waiver. The assessed component remains the voucher, so the waiver does not
set its $D$ value to true. This case challenges an implementation that stores a
single offer-level discount flag and reuses it for every component. If that flag
is the only available data, the adapter has insufficient information to apply
the proposed component-level rule. It should request a finer assessment rather
than derive a convenient value.

Case g07 leaves the discount classification unresolved for a component otherwise
meeting the trigger. If $D$ is false, the prohibition is triggered; if $D$ is
true, the exception applies. Both possibilities remain consistent with the
supplied evidence. Unknown is therefore the informative answer. A system that
turns this unknown into false would remove the exception by default; one that
turns it into true would invent the exception.

Case g16 provides known discount evidence while leaving the gift and promotional
classifications unknown. Within known scope, a true $D$ makes $\neg D$ false,
which settles this conjunction as false. The case demonstrates that every
missing fact need not prevent every narrow conclusion. Its safety depends on
the specific conclusion being reported. It would be unsafe to broaden the
result into a complete clearance of the marketing arrangement.

The contrast between g04 and g16 is particularly useful for reviewer training.
An established exception can settle a condition, but only for the object to
which it applies. A precise value attached to the wrong subject is not better
evidence than an unknown value attached to the right one.

\subsection{Classification is not supplied by the formula}
Case g08 concerns additional returns whose gift classification has not been
settled. The circular's qualified treatment of such returns motivates an
assessment; it does not establish $G$ for every arrangement. The proposed
candidate must retain unknown when the remaining facts would make that
classification decisive. Replacing the missing judgment by an unconditional
``additional return implies gift'' would create a new interpretive rule.

Cases g06 and g22 establish that no gift is offered. In g06 the discount
classification is also missing. The false gift predicate settles the body as
false under the declared Boolean semantics. Case g22 makes the point with a
known product link: a promotional link alone does not establish that a gift
exists. These cases catch a candidate that assumed $G$ from the context instead
of requiring evidence for it.

Case g05 establishes that neither product-link route holds. This also makes the
body false, while leaving other regulatory questions outside the result. It
does not establish that every incentive unrelated to an identified product is
permissible under all applicable requirements. The selected condition is only
one element of a broader assessment.

Case g17 separates regulatory classifications. A separate assessment finds no
structured product, but the component is a non-discount gift with a relevant
product link. The gift trigger remains true. The case challenges a workflow
that exits early after resolving one concern and never evaluates another.
Correct integration therefore requires an inventory of outstanding controls,
not a single reassuring flag named \texttt{regulatory\_check\_passed}.

\subsection{Applicability precedes an in-scope answer}
Cases g10 and g11 establish that, respectively, the actor and the product/activity
are outside the selected profile. The response is out of scope. The underlying
Code requirement may have a broader reach; the example cannot answer that
broader question. Treating out of scope as false would erase the distinction
between an in-scope condition that is not triggered and a question the selected
control does not purport to decide.

Case g12 leaves the fund/activity scope unknown. The system cannot establish
whether the profile applies. Case g13 leaves actor scope unknown while
establishing no gift. A Boolean programmer might be tempted to use the false
gift fact to return false immediately. The declared profile instead reports
unknown scope. This preserves the fact that the system has not established the
context in which its body answer would be meaningful.

This is a design choice that an independent implementation must reproduce.
Another language could define a different scope-combination policy, but it
would then implement a different specification. Differential testing only has
a clear target after the policy has been made explicit. A reviewer must also
decide how an unknown scope result affects the transaction workflow.

\subsection{Contradiction, expiry and later knowledge}
Case g18 supplies conflicting gift evidence. The result is conflict, not
unknown and not the value favored by the latest-arriving record. Unknown means
the evidence does not determine a value. Conflict means admissible records
assert incompatible values. These require different remediation: obtaining
missing evidence in the first case, and resolving or superseding contradictory
evidence in the second.

Case g19 combines an out-of-scope actor with conflicting discount evidence.
The current static-closure conflict policy returns conflict before scope
evaluation. This is deliberately conservative. The example must teach the
policy honestly because a reviewer might otherwise expect out of scope and
mistake the difference for a Java bug. The policy is recorded as engineering
semantics, subject to separate operational review.

Case g20 makes the gift evidence stale at assessment time. The rule does not
continue treating an expired classification as current. Case g21 places the
recording time after the historical knowledge cutoff. A historical assessment
cannot use that later knowledge. Both become unknown when the missing gift
classification remains decisive, but their reason codes and repair actions
differ. Refreshing an expired record does not justify rewriting what was known
at an earlier point in time.

All 22 cases therefore contribute a distinct aspect of the specified behavior.
Some concern the wording of the selected restriction. Others concern evidence,
scope and operational conventions needed to execute the reading. The case
packet should label those sources of expected behavior separately. Otherwise
an implementation convention can acquire the apparent authority of a legal
sentence simply because both are stored in the same JSON file.

\section{What the 729 combinations cover}
\label{sec:circular-enumeration}
Each of the six abstract inputs takes one of three values: true, false or
unknown. The number of assignments is therefore
\begin{equation}
 3^6=729.
 \label{eq:case-domain-size}
\end{equation}
This is a statement about a finite abstraction. It does not count legal
interpretations, marketing documents, evidence histories or real client offers.
The enumeration is exhaustive for that abstraction because it visits every
combination of the six declared values.

The independent truth oracle handles scope first. For the four body facts it
replaces each unknown by both possible Boolean values and evaluates the formula
under every completion. If every completion returns true, the oracle returns
true; if every completion returns false, it returns false; otherwise it returns
unknown. The code does not import RuleIR or the reference evaluator. This makes
it useful for detecting a defect in how the expression was encoded or evaluated.

For this particular body, each input occurs once and with a single polarity.
The completion calculation agrees with the current three-valued operations.
That agreement should not be generalized to every expression. For example, a
formula containing a variable and its negation can be true under every Boolean
completion while a compositional three-valued evaluator returns unknown. The
next chapter derives that difference. The oracle's representation is independent
of the candidate AST, but its semantic scope remains specific.

Conflict, stale evidence and later-recorded evidence are not extra members of
the three-value enumeration. Their normalization and precedence are exercised
by named cases and other application checks. The count 729 must not be used to
imply exhaustive testing of those mechanisms. Nor does it exhaust numerical
inputs, dates, arbitrary rule graphs or the Java runtime's entire implementation.
A result summary should always name the finite object that was exhausted.

The comparison checked full semantic responses and relevant traces, while
respecting each backend's own identity. It would be wrong to require the Python
and Java engine identifiers to be identical: those identifiers describe
different implementations. It would also be wrong to ignore all metadata in a
comparison, because source, evidence and rule identities are part of an
explainable decision. Chapter~\ref{ch:verification} specifies the distinction
between semantic agreement and implementation-specific provenance.

\section{Mutations ask whether the tests can notice an error}
\label{sec:circular-mutations}
Successful cases can be uninformative if they would also pass a wrong program.
Mutation testing addresses this by deliberately changing the rule and asking
whether the checks detect the change. A mutation is not a plausible legal
interpretation merely because it can be compiled. It is an instrument for
testing the sensitivity of the evidence.

The first retained mutation removed the fee-discount exception. Cases g03, g07,
g16 and g19 detected the change. Some of these detections concern the logical
exception, while g19 also exposes the effect of changing the static dependency
closure. Removing an expression can remove a referenced fact from conflict
analysis. A mutation report should therefore explain why a case detects a
change rather than treating every detection as the same kind of evidence.

The second mutation omitted the product-type route. Cases g02, g14 and g15
detected it. Their relationship to the source is direct: they isolate the
route that the mutation removed, including a case where that route is
unresolved. A baseline case with a named fund could not provide this evidence.

The third mutation assumed that a gift existed without requiring its evidence.
Cases g06, g08, g18, g20, g21 and g22 detected the resulting change. The set
includes known absence, missing classification, conflict, expiry and later
knowledge. This shows why evidence-handling cases are necessary even for a
small logical condition. A wrong classification default can alter outcomes
across several superficially different scenarios.

The fourth mutation ignored the selected actor/product scope. Cases g10 through
g13 detected it. The expected difference includes both out-of-scope and unknown
results. If the test harness compared only a final allow/deny bit after
flattening all other statuses, it could conceal these errors.

All four mutations were compiled into Java before comparison. The result thus
concerns executable altered candidates, not a textual inspection of a proposed
patch. Still, detecting four mutations does not establish detection of every
possible implementation or interpretation defect. Mutation operators are chosen
by people or software and can themselves omit an important failure mode. The
ensemble evaluation must maintain a documented fault model and extend it when
new failures are discovered.

\section{A new discrepancy: what counts as a discount?}
\label{sec:circular-cashback}
The existing tests assume that a reviewer can supply the discount classification.
To investigate the missing interpretation step, consider a synthetic offer in
which a customer pays a fee and later receives cash back. Two candidate
definitions might be proposed. One restricts the exception to a reduction of
the fee initially charged. Another includes a documented refund of a fee that
was actually paid. The source packet alone has not established that both
readings are legally tenable. Their purpose here is to expose a question whose
answer affects implementation.

Fix the actor, activity and promotional link, and assume for the comparison
that the component otherwise meets the gift condition. Let the fee paid be
HK\$100 and the later rebate also HK\$100. Under the first proposed definition,
the later transfer is outside the defined exception. Under the second, it may
fall within the exception if the required refund relationship is established.
Those conditions must be derived from authority or review; the example does
not introduce them as new SFC requirements.

The next action should be a targeted investigation of the applicable definition,
Code and FAQ, not a request for the same model to vote again on the word
``discount.'' A source that explicitly includes the relevant refund treatment
would support refinement of the second candidate. A source that excludes it
would defeat that candidate in the applicable scope. A source discussing
another product, actor or period may be informative without settling the
question. The retrieval result needs an applicability assessment before it can
resolve the discrepancy.

Now alter the cash amount to HK\$150. If a candidate definition treats only a
refund of an actual HK\$100 fee as falling within the exception, the extra
HK\$50 raises a further classification question. It does not follow that the
system can automatically split the transfer into a lawful and unlawful part.
The legal unit of assessment and treatment of mixed benefits must first be
settled. An arithmetic decomposition is easy; its legal significance is the
question under investigation.

A second variation changes the evidence. The marketing material describes the
transfer as a fee refund, but no record establishes that a fee was paid. The
candidate's definition may be clear while its decisive factual premise remains
unknown. This discrepancy belongs to fact classification or evidence
availability, not necessarily to interpretation. The resolution controller
should request the relevant evidence or return unresolved, rather than
changing the definition to make the missing record unnecessary.

A third variation combines a refund with a separate voucher. A whole-offer
exception would remove the need to assess the voucher independently. The
proposed component-level interpretation would retain it. This is a different
branch in the hypothesis space from the timing of a refund. Preserving separate
issue identifiers prevents a source that resolves the timing question from
being incorrectly treated as resolving the mixed-offer question as well.

These variations show how an ensemble can make review more precise. It should
produce a small set of materially different definitions, the source reasons
for each, and the smallest scenarios on which they diverge. The reviewer then
decides an identifiable question. A long narrative saying that an offer appears
reasonable would provide less actionable evidence.

\section{A complete dry run of the proposed workflow}
\label{sec:circular-dryrun}
The following dry run specifies proposed behavior. It is not a report that the
new ensemble has already run. Start by creating an assessment whose context
contains the retained circular revision, declared distributor/fund scope,
selected effective interval and an inventory of dependencies. The historical
engineering example uses a synthetic September 2026 interval; the workflow must
not infer a regulatory commencement date from that setting.

The source-review path inventories the circular and flags the incorporated Code
and FAQ as unresolved. The interpretation path proposes the component-level
gift rule. A second interpretation path omits the product-type route. A
definition-focused path asks whether cashback is a discount and whether a
combined offer changes the assessed unit. At this stage the system has a
candidate set and several issues, not a winner.

The comparator identifies a decision difference between the complete promotional
condition and the candidate omitting the product-type route. It constructs a
case with a gift linked only to a fund category and no exception. The
source-grounding check points to the relevant paragraph. The candidate author
can propose a repair adding the missing route. The system stores the repaired
candidate as a new revision, reruns the discrepancy case and the retained
regression cases, and preserves the original candidate and its rejection reason.

The cashback definition issue follows a different route. Retrieval attempts to
obtain the incorporated sources with the right identity and version. If the
sources are unavailable, another language-model explanation cannot close the
issue. The controller records the failed acquisition and the remaining
dependency. It may retry an authorized retrieval route within its budget, or
stop automatic work because no permitted action can supply the missing
authority. The issue remains unresolved either way.

The mixed-offer issue can produce a useful witness even before legal resolution.
The comparison holds the voucher facts fixed and toggles the presence of a
separate fee waiver. If a candidate's voucher decision changes solely because
of the other component's exception, the report explains the dependence. A
reviewer must then confirm the unit to which the exception applies. The witness
does not grant the controller authority to make that legal decision itself.

The accepted formal candidate can be compiled and checked in parallel with this
research. Candidate compilation is valuable because it exposes implementation
defects early. It does not change the issue status. The report may legitimately
say that the candidate passed all specified execution checks and remains
unavailable for release because source interpretation is unresolved. Keeping
those two observations together is more informative than suppressing either.

When automatic work reaches its finite limit, the controller freezes a report.
It identifies the source revision, surviving definitions, rejected encodings,
discriminating scenarios, completed actions, failed actions, unresolved
dependencies and the exact effect of each unresolved issue on the control.
There is no final vote to select a legally preferred answer. The state is ready
for targeted human review or blocked on a named dependency, according to the
reason for stopping.

A designated reviewer can later resolve an issue using newly obtained authority.
That creates a new context revision. Dependent candidate scores and comparisons
are invalidated where appropriate, and the formal rule and tests are updated.
The earlier report remains a faithful account of what was unresolved before the
new material arrived. This is essential for reproducibility: a current reviewer
must be able to see why the earlier system abstained.

\section{From the JSON oracle to independent adjudication}
\label{sec:circular-independent}
The filename \texttt{oracle.json} is an engineering term. It designates a source
of expected results for tests. It does not establish institutional authority or
legal independence. The file states that the same agent authored it and the
candidate. A future user interface should display that authorship rather than
letting the label ``oracle'' imply a gold standard.

Independent adjudication should begin from the retained source packet and
declared task, before exposing reviewers to the proposed formal expression.
One reviewer identifies requirements, definitions and difficult cases. Another
can independently develop or challenge the interpretation. Disagreements are
recorded and, where possible, resolved with source reasons. A third role may
adjudicate material disagreement. Independence here concerns information and
responsibility as well as the identity of a person or model.

After that source-based stage, reviewers can inspect the candidate rules and
their discriminating cases. This order reduces anchoring on the generated
answer. It does not make reviewers infallible. The reference packet must allow
an unresolved category, later correction and competing defensible readings.
If a reviewer cannot determine the treatment of a cash rebate without the
referenced FAQ, the correct reference outcome is not a guessed Boolean chosen
to make benchmark scoring convenient.

The original 22 cases should remain frozen as historical engineering evidence.
New adjudicated cases belong in a separately versioned packet with its own
authors, source context and review decisions. A changed expected outcome should
explain whether the earlier interpretation was rejected, the scope changed,
new evidence became available or the evaluation semantics were corrected.
Those causes imply different repairs and different consequences for existing
releases.

The chapter has now followed a source to a rule, tested the rule, exposed
unresolved definitions and specified how a future ensemble should handle them.
The remaining task is to state precisely what each formal check can prove.
Without that precision, a solver's success could once again be mistaken for
the missing judgment about meaning.

\chapter{What can be proved, and under which assumptions}
\label{ch:proof}

\section{Giving a correctness claim a precise target}
\label{sec:proof-target}
The second-circular exercise leaves us with two questions that sound similar but
have different evidentiary requirements. Does the program compute the proposed
gift rule correctly? Does the proposed gift rule express a justified reading of
the circular? The first question becomes mathematical once the rule and its
inputs have precise meanings. The second includes deciding which meanings the
words, context and authority support. A proof of the first cannot silently
include a proof of the second.

To see why, suppose a programmer writes a formula that checks gifts promoting a
specific product and omits the product-type route. A proof assistant can verify
that an implementation returns exactly that formula's value for every Boolean
assignment. The proof is correct about its target. The target is incomplete
relative to the proposed source interpretation. Making the proof longer, using a
more powerful solver or verifying another implementation will not add the
missing route. The source-to-formula mapping has to be challenged directly.

This is not a special weakness of legal language. Formal verification always
depends on the proposition that was formalized. The virtue of a formal method
is that it makes the conditional claim exact and inspectable. It does not remove
the need to establish that the proposition is the one the institution intended
to rely on. Catala's separation between legal authoring and compiler semantics
is useful precisely because it keeps these responsibilities distinct
\citep[core language and compilation sections]{catala2021}.

Let $E$ denote a versioned source packet and $S$ a controlled specification
adopted for a declared interpretation. Let $F$ be a RuleIR program and $J$ the
generated Java artifact. Let $A$ contain the explicit vocabulary definitions,
domain restrictions and operational assumptions. The input domain $D_A$ consists
of facts and histories allowed by those assumptions. We can now state a proposed
preservation obligation:
\begin{equation}
 \forall x\in D_A,\qquad
 \pi\bigl(\eval_{\CNL}(S,x)\bigr)
 =\pi\bigl(\eval_{\IR}(F,x)\bigr)
 =\pi\bigl(\eval_{J}(J,x)\bigr).
 \label{eq:preservation-chain}
\end{equation}
The projection $\pi$ selects the semantic result components that the contract
requires to agree. These include the decision category and the specified
explanation and evidence behavior. They exclude fields whose purpose is to
identify different backends. Chapter~\ref{ch:verification} makes that projection
concrete. Equality in equation~\eqref{eq:preservation-chain} is a proposed formal
target for a supported language fragment, not a theorem already proved about
the delivered application.

The claim that $S$ is justified by $E$ remains a separate interpretation claim.
It can be supported by source anchors, explicit alternatives, reviewed
definitions, adversarial cases and an attributable judgment. When some part of
$E$ has itself been translated into a formal language, a theorem can relate
those formal objects. The choice that made the formal source representation
faithful still needs justification. Repeating that translation boundary does not
eliminate it.

\section{Quantifiers change which entity carries the obligation}

The phrase ``for each'' introduces a choice that a Boolean expression can easily
hide. To see it without assuming knowledge of predicate logic, consider a
synthetic requirement: every relevant benefit in an offer must satisfy a specified
condition. Let $B(o)$ be the finite set of benefit components in offer $o$, and
let $P(b)$ mean that component $b$ has the property under investigation. The
universal reading is
\begin{equation}
 \forall b\in B(o),\ P(b).
 \label{eq:universal-benefit}
\end{equation}
It requires every component to satisfy the condition. An existential reading,
\begin{equation}
 \exists b\in B(o),\ P(b),
 \label{eq:existential-benefit}
\end{equation}
requires only one. These expressions are not two styles of writing the same rule.
For an offer with two components, one satisfying $P$ and one not, the universal
reading is false and the existential reading true.

The distinction becomes more consequential when an exception is involved. Suppose
$G(b)$ identifies a gift and $D(b)$ an excepted discount. A component-level
prohibition trigger might be
\begin{equation}
 \exists b\in B(o),\ G(b)\land\neg D(b).
 \label{eq:component-exception}
\end{equation}
An alternative whole-offer encoding might be
\begin{equation}
 \left(\exists b\in B(o),\ G(b)\right)
 \land\neg\left(\exists b\in B(o),\ D(b)\right).
 \label{eq:offer-exception}
\end{equation}
The first asks whether any component is a non-excepted gift. The second suppresses
the trigger whenever any component is a discount. A voucher combined with a fee
waiver distinguishes them immediately: the second formula can let the waiver
cancel the voucher's effect, while the first continues to assess the voucher.
These are proposed logical readings for analysis, not a legal determination about
mixed offers under 23EC46.

A source interpretation must decide what the variable ranges over before a proof
about either formula is useful. Is the subject a benefit, a contract, a campaign,
a transaction or a client relationship? Are components defined by economic
substance, contractual terms or an operational product record? The formal
notation exposes this question because $B(o)$ cannot remain undefined in an
implementation-ready specification.

The current RuleIR fragment has scalar facts rather than unrestricted quantified
collections. A reviewed adapter can create one snapshot per component and a host
orchestration can aggregate results under an explicit rule. That aggregation is
another specification obligation. It must define how true, false, unknown,
conflict and out-of-scope component results combine, and how the host knows the
component list is complete.

An empty list deserves attention. In classical logic, a universal statement over
an empty set is true and an existential statement is false. That convention is
mathematically coherent, but a missing component inventory must not be treated as
an established empty offer. The adapter needs an evidence status for collection
completeness. Otherwise a data omission can pass a universal check by vacuity.
The same pattern occurs when no argument extensions are returned or no test cases
are selected: an empty computation can look successful unless existence and
coverage are checked separately.

Unknown classification also affects aggregation. If one component is known to
trigger a prohibition and another is unknown, the selected logical aggregation
may already establish a trigger. Yet the unknown component may remain relevant
to another obligation or to the completeness of the offer review. A top-level
known result should not erase the record of missing facts. The distinction between
missing and blocking inputs in the runtime supports this more informative report.

\section{Necessary conditions, sufficient conditions and permissions}

Another frequent translation error changes the direction of implication. Consider
a synthetic statement that a process may use a streamlined route only if consent
has been obtained. Let $S$ mean use of the streamlined route and $C$ mean established
consent. The necessary-condition reading is $S\Rightarrow C$: using the route
requires consent. It does not follow that $C\Rightarrow S$ or that consent alone
authorizes the route. Other eligibility conditions may still apply.

A truth-table example makes the difference concrete. With consent true and an
unmet financial condition, the necessary-condition statement is not violated by
refusing the streamlined route. An implementation that treats consent as sufficient
would proceed. Both programs may mention the same two words, but their decisions
represent different logical relationships.

Likewise, an exception to one prohibition is not a general permission. If the
selected condition $R$ is false, we know only that this particular trigger is not
satisfied under the accepted facts and semantics. We do not infer that every
other prohibition is false or every positive obligation fulfilled. The host must
compose independently applicable controls without upgrading the absence of one
trigger into universal clearance.

In a closed, complete formal theory, one can define permission in terms of the
absence of a prohibition if that is the intended logic. Legal systems and practical
source packets often do not provide the completeness needed for that move.
Explicit permission, lack of evidence of prohibition and a bank's operational
choice should therefore have different representations. The normative-language
literature examined in Chapter~\ref{ch:languages} provides several treatments;
RuleIR's narrow condition results do not choose a universal deontic logic.

The words ``unless'' and ``except'' also require a structural decision. A simple
propositional rewriting may express an exception as a negated condition, while a
default system represents a defeasible conclusion that can be overridden. These
can agree on a selected complete fact pattern and differ under missing evidence,
conflicting exceptions or additional rules. The translator must identify which
behavior the approved specification requires.

For the circular example, using $\neg D$ inside the selected body is an explicit
modeling decision tied to the reviewed factual definition of discount. If $D$ is
unknown, the program should not act as though the exception is certainly absent.
The three-valued semantics makes that uncertainty visible. A closed-world database
query that treats a missing discount row as false would encode a different
assumption, which requires separate evidence and approval.

\section{A proof obligation should expose its assumptions}

A useful proof statement includes the source-independent mathematical property,
the accepted semantic definitions and the domain assumptions. For example:
for every well-typed Boolean expression in a specified grammar and every normalized
fact assignment without conflict, the reference and target evaluators produce
the same three-valued result. This statement is precise enough to prove by
structural induction. It does not include the claim that a particular English
paragraph was mapped to the right expression.

The interpretation assumption can be represented separately as a reviewed
proposition: under source packet $S$, profile $P$ and definitions $D$, expression
$e$ represents the selected requirement. A preservation proof then establishes
that target program $J(e)$ behaves like $e$. Composing these steps is useful, but
the first premise remains a source and review judgment unless a restricted
controlled language and its translation semantics are themselves the authoritative
input.

A proof assistant will faithfully carry an unjustified premise into a theorem.
If we assume that every cashback is a discount, it may prove the corresponding
exception applies. The theorem is conditional on that assumption. It is not a
discovery that the assumption is legally correct. The proof report must display
such premises in language the reviewer can understand.

The same principle applies to abstraction. A model may replace a complex offer
with six Boolean inputs. The proof can establish behavior for every assignment
of those inputs. To apply it to the real offer, the adapter must establish that
the six inputs faithfully represent the relevant facts and that the omitted
features cannot change the selected requirement within the declared scope.
That relationship is an abstraction argument, not an automatic consequence of
exhaustive truth-table testing.

The workbench should therefore attach an assumption register to every formal
result. Each assumption names its provenance, operational meaning and review
status. A proof that depends on an unresolved definition is retained as a valid
conditional engineering result, while release remains blocked. This preserves
the value of formal work without allowing its precision to conceal an uncertain
starting point.

\section{Expressions, types and meanings}
\label{sec:proof-language}
A formal language is a collection of permitted expressions with rules for
deciding what they mean. It need not look like a conventional programming
language. A decision table, a restricted English sentence or a state-transition
diagram can have formal semantics if each permitted construct has a precise
interpretation. An unrestricted paragraph usually has more interpretive freedom
than such a language allows.

The simplest fragment needed for the gift example contains Boolean literals,
references to named facts, conjunction, disjunction and negation. In a compact
grammar, an expression $e$ is generated by
\begin{equation}
 e ::= \mathsf{true}\mid\mathsf{false}\mid\mathsf{fact}(a)
       \mid\mathsf{not}(e)\mid\mathsf{all}(e_1,\ldots,e_k)
       \mid\mathsf{any}(e_1,\ldots,e_k),\quad k\geq 1.
 \label{eq:boolean-grammar}
\end{equation}
Here $a$ is the identifier of a declared Boolean fact. The grammar specifies
which trees are valid expressions; it does not decide whether the fact denotes
a legally appropriate classification. That meaning belongs to the declaration
and its source-linked interpretation.

Types prevent several avoidable errors. A Boolean condition is different from
a monetary amount. A date is different from a number of days. Money denominated
in HK cents is different from a foreign-currency amount before conversion.
The current RuleIR scalar types are Boolean, integer, HKD money and date. A
comparison accepts compatible ordered values. Addition preserves a numeric
type. No implicit conversion turns a missing Boolean into zero or a dollar
amount into cents.

Typing is valuable because it can reject an entire class of malformed
specifications before any client case is evaluated. It is also limited. Two
variables can both have the Boolean type while referring to different subjects
or definitions. A type checker cannot infer that an offer-level discount flag
was wrongly substituted for a component-level classification unless the type
and declaration system explicitly represent that distinction. The proposed
extension should therefore retain subject, scope and definition identities
alongside the scalar type.

An expression's semantics is its evaluation rule. For known Boolean inputs,
\texttt{all} is true exactly when every operand is true, and \texttt{any} is true
when at least one operand is true. The gift body is a tree whose root is an
\texttt{all}, with children for the gift fact, the two-route disjunction and
the negated exception. The tree makes parentheses explicit. A source sentence
can be attached to the root and narrower evidence anchors to its children.

The chosen language is deliberately finite and acyclic. A reference cycle such
as rule A depending on B while B depends on A is rejected in the current
profile. This avoids an undefined recursive evaluation. Richer legal formalisms
can assign meaning to recursive or mutually defeating rules, but doing so
requires a specified semantics. A compiler must not guess that an unsupported
recursive construct means the same thing as the current acyclic profile.

\section{Unknown is an information state}
\label{sec:proof-unknown}
The bank often lacks a decisive fact when a rule is evaluated. In the gift
example, the evidence may not establish whether an additional return is a gift.
The program should distinguish that absence from evidence that no gift exists.
We therefore need an unknown value, written $\Unknown$, alongside known true
$\True$ and known false $\False$.

The current Boolean operations follow strong-Kleene tables. The name identifies
a precise three-valued calculation, which can be taught through the two cases
already encountered. If a product-specific promotional link is true, then the
disjunction of that link with an unknown type link is true. If the specific
link is false and the type link is unknown, the disjunction is unknown. A true
route suffices; an absent route cannot settle the other one.

\begin{table}[htbp]
\centering
\begin{tabular}{cc|cc}
\toprule
$a$ & $b$ & $a\land b$ & $a\lor b$\\
\midrule
$\True$ & $\True$ & $\True$ & $\True$\\
$\True$ & $\False$ & $\False$ & $\True$\\
$\True$ & $\Unknown$ & $\Unknown$ & $\True$\\
$\False$ & $\False$ & $\False$ & $\False$\\
$\False$ & $\Unknown$ & $\False$ & $\Unknown$\\
$\Unknown$ & $\Unknown$ & $\Unknown$ & $\Unknown$\\
\bottomrule
\end{tabular}
\caption{The omitted rows follow by exchanging the operands. Negation swaps
true and false and leaves unknown unchanged. The table describes logical values;
evidence conflict and errors are handled separately.}
\label{tab:kleene}
\end{table}

An established exception can make the gift conjunction false even when another
body fact is unknown, because one false conjunct is enough. That is the logic
behind g16. It is not a general instruction to ignore missing information.
Unknown scope is handled before the body, and conflicting evidence has its own
precedence. The evaluation result must retain which question was answered and
why the remaining unknown body facts did not change that answer.

A convenient symbolic encoding represents a three-valued Boolean by two ordinary
Boolean indicators: one means ``known true'' and the other ``known false.'' The
unknown value has neither indicator. If $t_a$ and $f_a$ encode $a$, then
\begin{align}
 t_{a\land b}&=t_a\land t_b,& f_{a\land b}&=f_a\lor f_b,\label{eq:kleene-and}\\
 t_{a\lor b}&=t_a\lor t_b,& f_{a\lor b}&=f_a\land f_b,\label{eq:kleene-or}\\
 t_{\neg a}&=f_a,& f_{\neg a}&=t_a.\label{eq:kleene-not}
\end{align}
The domain excludes $(t_a,f_a)=(1,1)$. That pair would assert both truth and
falsity, which belongs to a different conflict mechanism in the current design.
The equations reproduce Table~\ref{tab:kleene}: for example, true and unknown
have true indicators 1 and 0, so their conjunction cannot be known true; neither
has a false indicator, so the result is unknown.

This encoding allows a solver to search disagreements while preserving missing
information. Replacing each unknown by an arbitrary ordinary Boolean inside
one execution would ask a different question. It could produce a definite answer
chosen from a possible completion without showing that other completions differ.
A comparison tool must name the semantics it encoded, not merely report that it
used a solver.

There is a subtler distinction between compositional three-valued evaluation
and evaluating every Boolean completion. Let $p$ be unknown. Strong-Kleene
evaluation gives
\begin{equation}
 p\lor\neg p=\Unknown\lor\Unknown=\Unknown.
 \label{eq:unknown-tautology}
\end{equation}
Under every consistent Boolean completion, however, the expression is true:
if $p$ is true the first term holds, and if $p$ is false the second holds. A
completion-based semantics can therefore return true. Neither calculation is
ambiguous once its rules are specified; they are different semantics.

For the read-once gift body, the retained truth oracle's completion method agrees
with the compositional calculation. That particular equivalence does not make
the completion oracle a general replacement for RuleIR semantics. A future
ensemble should treat an unexpected disagreement as a possible semantics
mismatch before assuming that one interpreter contains a bug. This example is
also a reason to test repeated variables when extending the comparison fragment.

\section{Conflict, scope and errors are not extra Boolean values}
\label{sec:proof-tags}
Unknown truth and conflicting evidence should not be combined casually into one
truth table. The current runtime first normalizes records by validity,
recording time and explicit supersession. Two admissible non-superseded records
with incompatible values produce conflict. That result records a problem in
the evidence supporting a fact. It does not ask the Boolean evaluator to choose
which record is right.

The static dependency closure contains every fact reachable from the requested
rule through its scope, body, referenced rules and exception branches. Conflict
in that closure prevents evaluation from treating a convenient branch as a way
to evade contradictory evidence. A conflict in an unrelated declared fact does
not automatically block the requested rule. This distinction makes the
conservative policy bounded and reproducible.

After successful structural validation and conflict handling, scope determines
whether the body is an appropriate question. A known-false scope produces out
of scope. An unknown scope produces unknown with a scope reason. Only a
known-true scope permits the in-scope body result. Thus an out-of-scope result
is a statement about the relation between the query and the selected rule,
not a false answer to the rule's substantive condition.

Malformed data and unsupported constructs introduce a further distinction.
An input with an invalid date or duplicated conflicting identifier may be an
error, not legal uncertainty. If the system cannot parse a request, it cannot
honestly report that the regulatory condition is unknown on a well-defined
request. The API should identify the boundary error and the location of the
malformation. Operationally both may prevent automatic action, but their causes
and repair routes differ.

We can therefore describe the result as a tagged value whose tag identifies
the kind of answer. The current contract includes true, false, unknown,
conflict, out of scope and error, together with diagnostics and provenance.
The tags are not interchangeable labels for a single confidence score. The
ensemble discrepancy taxonomy must preserve them so that a failed parser is
not processed as a legal disagreement and an unresolved definition is not
processed as a transient network failure.

The precedence is part of the specification. Reordering checks can change
both a result and its explanation. The g19 example deliberately returns
conflict even though an actor scope fact is false. A Java implementation that
short-circuits on scope before consulting the required conflict closure would
disagree with the current target. That disagreement is wrong relative to the
declared execution semantics, independently of whether the bank later chooses
to revise the policy.

\section{A readable specification with a fixed translation}
\label{sec:proof-cnl}
Controlled natural language keeps the familiarity of English while restricting
how sentences are formed. The restriction matters because the reader should
be able to determine the formal structure from the sentence without asking a
model to interpret it again. Logical English provides one developed example of
this approach through explicit templates, variable binding and restrictions on
ordinary linguistic ambiguity \citep[language and law examples]{logicalenglish2023}.

For our narrow gift example, a specification can state: ``For an assessed
component in the established distributor/fund scope, the selected gift
restriction is triggered if the component is a gift, at least one of the two
specified promotional links holds, and the component does not meet the defined
discount exception.'' The specification must then define each term and its
evidentiary state. This sentence is reviewer-facing exposition. A machine
translation needs an accompanying restricted structure in which the three
conjuncts, disjunction and negation are explicit.

An appropriate authoring interface might present one condition per line with
typed connectives and named definitions. The reviewer sees the English label
and can inspect the exact expression beneath it. Free text remains available
for reasoning and caveats, but it does not secretly alter the executable
condition. A change to a definition or connective creates a new specification
revision. Editing an explanatory paragraph must not bypass the formal diff.

The controlled specification also records the legal character of the source.
A statement framed as an obligation, a prohibition, a permission, a discretion
or an expectation should retain that category. Translating all of them into a
Boolean called \texttt{allowed} would destroy distinctions before verification
begins. The current RuleIR chiefly computes conditions. A reviewed operational
policy determines how those conditions affect a bank action. Richer normative
and temporal languages are introduced in Chapter~\ref{ch:languages}.

A fixed translation from restricted syntax to RuleIR is a valuable target for
proof. Every permitted construct has a known translation rule; unsupported
sentences are rejected or left as unresolved authoring work. If a reviewer
writes an unrestricted sentence outside that syntax, a model can propose a
restricted rendering, but that proposal returns to the interpretation process.
The system cannot claim deterministic controlled-language compilation merely
because an LLM usually produces valid JSON.

Roundtrip verification supplies an additional check. Translate a formal
candidate into readable language, formalize the rendering again, and compare
the two formal candidates. A mismatch can reveal drift. Agreement can show
stability across the tested transformations. The same omitted source condition
can survive both translations, so the roundtrip result remains one component
of the evidence, as acknowledged by \citet[method and limitations]{roundtrip2026}.
It cannot replace the independent source inventory established earlier.

\section{How a preservation argument is constructed}
\label{sec:proof-induction}
The phrase ``prove the compiler correct'' becomes more useful when the reader
can see the structure of such a proof. Consider only the finite Boolean
expression grammar in equation~\eqref{eq:boolean-grammar}. Let its source
evaluator use Table~\ref{tab:kleene}. Suppose a translation maps literals,
fact references and each connective to corresponding target operations.
Assume the target fact lookup returns the same normalized tagged value for
every declared fact and that its primitive operations implement the same
tables.

\begin{proposition}[Preservation for the Boolean expression fragment]
Under those assumptions, translating any well-formed finite expression in
equation~\eqref{eq:boolean-grammar} preserves its three-valued result for every
normalized fact assignment.
\label{prop:boolean-preservation}
\end{proposition}
\begin{proof}
Use induction on the construction of the expression. A true or false literal
has the same value by the literal translation. A fact reference has the same
value by the lookup assumption. For a negation, the induction hypothesis gives
the same child value in both evaluators, and the target negation table maps
that value to the same result. For a conjunction or disjunction, the induction
hypotheses give equality for every child. Applying the identical connective
table to the identical sequence of values gives the identical parent value.
These cases cover every constructor in the grammar. The finite syntax tree
ensures the induction reaches a base case on every path.
\end{proof}

This small proof is exact, but its scope is intentionally visible. It assumes
correct target primitives and lookup. It does not establish the current Java
implementation satisfies those assumptions. It does not cover scope, conflict
normalization, numeric operations, exceptions, errors or traces until separate
lemmas include them. And it says nothing about whether the expression is a
justified interpretation of paragraph 10.

The next proof obligations follow the same pattern. A normalization lemma must
show that both implementations select the same admissible evidence under the
same times and supersession relation. A scope lemma must preserve the tagged
result and the decision to skip the body. A conflict lemma must preserve the
static closure and conflict precedence. A trace lemma must connect each
evaluated or skipped node to the required source and evidence identifiers.
Only after these claims are assembled can the implementation claim be expanded
from a Boolean value to the full semantic result.

The actual generated Java class currently binds an immutable validated bundle
and invokes the Java semantic runtime. The proof boundary must follow that
architecture. Proving a small emitted wrapper correct while assuming an
unchecked runtime would leave most of the computation outside the theorem.
Conversely, a proof about a generic runtime still needs to show that the
generated class embeds the intended bundle and that the host invokes the
correct class. The theorem's trusted components should be named rather than
hidden under the word ``compiler.''

There is a useful discipline here for an implementation agent. Begin with a
small formal fragment whose semantics and generated behavior can be inspected.
Prove or exhaustively check its components, then extend the fragment with new
constructors and corresponding obligations. A failed lemma is information
about the specification or implementation. It is not a reason to weaken the
claimed target without explaining the change.

\section{Counterexamples and the domain they live in}
\label{sec:proof-counterexamples}
To compare two candidate interpretations, the system first needs a common
question and compatible inputs. Let $F_1$ and $F_2$ compute tagged results on
the same declared input domain. A solver can search for an assignment satisfying
\begin{equation}
 D(x)\land\left(\pi(\eval(F_1,x))\ne
                         \pi(\eval(F_2,x))\right).
 \label{eq:difference-query}
\end{equation}
Here $D$ states the permitted facts and relationships. In the product-type
example, a witness sets gift true, specific-product link false, type link true
and discount false, with scope established. The assignment has different
consequences under the complete and omitted-route rules.

The first solver check should establish that $D$ is satisfiable. If the domain
simultaneously requires a fact to be true and false, the difference query is
unsatisfiable for an unhelpful reason: there are no admissible inputs. The
correct response is \texttt{INCONSISTENT\_DOMAIN}, not equivalence. This
precheck is a concrete safeguard against a formally successful but meaningless
verification result.

If the domain is satisfiable and the difference query is satisfiable, the
assignment is a counterexample to equivalence in that domain. Before displaying
it to a reviewer, replay it through the ordinary interpreters. That replay
checks that the solver encoding and the runtime agree on what the witness
means. If replay disagrees with the solver prediction, the comparison harness
has failed and the witness cannot be used to adjudicate the candidates.

If the difference query is unsatisfiable, the supported conclusion is absence
of a counterexample under the encoded semantics and domain. The strength of
that result depends on the solver procedure and any checked certificate. A
machine-checked certificate provides a different assurance level from trusting
an unchecked solver response. A timeout, unsupported construct or unknown
solver result remains unresolved; it must never be converted to successful
equivalence because no witness was printed.

Domain assumptions can hide exactly the error we hoped to find. If a comparison
asserts that every product-type promotion also names a specific product, the
omitted-route mutation may become indistinguishable from the proposed rule.
That assumption would need a source or business justification. It cannot be
introduced merely because it makes the proof succeed. The report should show
which constraints exclude a witness and where those constraints came from.

Minimal witnesses are useful for human review, but ``minimal'' also requires a
definition. One option minimizes the number of facts changed from a familiar
baseline. Another minimizes monetary amounts within meaningful bounds. Neither
guarantees that the resulting scenario is legally or operationally feasible.
Source and domain validation remain necessary. The smallest formal assignment
may be a poor real-world question if it combines classifications that a reviewer
would never jointly accept.

\section{Proving a workflow property}
\label{sec:proof-workflow}
Some of the strongest useful guarantees concern the review process rather than
the legal reading. For example, the proposed ensemble should never mark an
assessment releasable while a material discrepancy remains unresolved. This is
a property of explicit records and transitions. It is amenable to a state-machine
argument once the terms ``material,'' ``unresolved'' and ``releasable'' have
precise stored meanings.

Let $Q$ be an assessment state, $M(Q)$ its set of unresolved material issues,
and $R(Q)$ the predicate that a release is authorized. A required invariant is
\begin{equation}
 R(Q)\ \Longrightarrow\ M(Q)=\varnothing.
 \label{eq:no-open-release}
\end{equation}
The invariant does not say that an empty issue set proves the interpretation
correct. It says that known unresolved material issues cannot be ignored by
the release transition. A faulty detector may fail to create an issue, which
is why independent coverage and detector evaluation remain necessary.

To establish the invariant, the initial state has no release. Every transition
that creates, reopens or materially changes an issue must invalidate dependent
release eligibility. The only transition that authorizes release must check
the issue set and all other required conditions in the same transaction that
records the authorization. If a concurrent worker can add an issue after the
check but before the release without a version check, the implementation does
not preserve the invariant. Database transaction boundaries are therefore part
of the proof assumptions, not an incidental engineering detail.

A second property concerns termination of automatic processing. Suppose a run
has a finite number of reserved actions, and every dispatched action consumes
one reservation even if the worker fails. If all automatic paths either
consume a reservation or enter a terminal state, then no run can dispatch more
actions than its budget. This argument requires crash and retry accounting;
counting only successful responses would permit unlimited failed calls.
Chapter~\ref{ch:ensemble} gives the exact controller and its termination
measure.

These guarantees directly support the user's requirement. They can establish
that the system reports remaining uncertainty and respects its retry limit.
They cannot establish that the detectors discovered every possible uncertainty.
The two responsibilities should be tested separately: detection quality against
independent cases, and workflow integrity against state-transition and failure
injection tests.

\section{Temporal duties require histories}
\label{sec:proof-temporal}
A Boolean assessment of an offer is not enough for every circular requirement.
Some duties become active after an event and must be performed before a
deadline. Others must remain satisfied throughout an interval. Consent can be
granted, withdrawn and granted again. A formal specification for these cases
needs histories and states, not just a snapshot of current fields.

Consider a hypothetical duty activated at time $a$ with an exclusive deadline
$d$. Under a declared non-preemptive achievement interpretation, a matching
performance at time $t$ satisfies the timing condition exactly when
\begin{equation}
 a\leq t<d.
 \label{eq:performance-window}
\end{equation}
The inequality states an engineering profile; a source may require a different
boundary or permit performance before activation. Those choices must be
reviewed and encoded explicitly. The event must also match the duty's actor,
action and subject. A timely action by the wrong party is not satisfaction.

Failure to find a performance record is not by itself evidence that the duty
was breached. The observation stream may be incomplete. The current event
runtime therefore requires a completeness assertion through the deadline before
inferring breach from absence. That assertion has provenance and a recording
time. A process clock cannot create it. This mirrors the source-coverage
problem: an empty retrieved set does not prove that the relevant object does
not exist.

Late performance and breach are also distinct. A duty can be performed after
its deadline without erasing the fact that timely performance was absent under
complete evidence. A later correction may establish that an earlier performance
was missed; that creates a new assessment with a different knowledge cutoff.
The original assessment remains available. Normative logic and contract-state
languages help express these distinctions, but their semantics must be selected
to match the intended duty \citep{violations2006,normative2016,stipula2021}.

The current profile supports a limited event model. It does not automatically
implement every waiver, sanction, compensatory duty or maintenance obligation
that a legal document could contain. Unsupported temporal constructs should
remain explicit implementation gaps. A translator that invents a deadline or
converts a continuing duty into a single check would create a new rule while
appearing to simplify the software.

\section{An assurance claim is a conjunction of different evidence}
\label{sec:proof-assurance}
We can now state the overall structure without pretending that one theorem
settles the entire process. Source retention establishes which material was
reviewed. Interpretation review establishes an attributable decision about that
material within a declared context. Formal checks establish consequences and
preservation properties for specified languages and domains. Integration checks
establish how a result reaches a bank action. Each supports a distinct premise
in the justification for using a control.

This structure prevents two opposite mistakes. The first is to dismiss formal
verification because it cannot prove the meaning of unrestricted English.
Preserving a reviewed specification across implementations is valuable, and
workflow invariants can rule out important classes of operational failure. The
second is to present formal verification as if it eliminated the interpretation
and evidence premises. A perfectly verified formula still needs the right
subject, source version and fact classifications.

For the bank's reviewers, the practical outcome is a set of answerable
questions. Which parts of the interpretation are accepted and by whom? Which
formal properties have actually been established? What lies outside the
supported fragment? Which data mappings remain unresolved? What would reopen
the conclusion? Those questions can be supported by exact records rather than
an undifferentiated confidence statement.

The next chapter examines existing languages and methods in that light. Each
is introduced by the problem it solves: readable authoring, exceptions,
alternative interpretations, events, or symbolic checking. Their combination
must preserve those distinctions instead of assuming that choosing one formal
language supplies the whole assurance argument.

\chapter{Languages that make interpretation inspectable}
\label{ch:languages}

\section{Choose the question before choosing the language}
\label{sec:languages-question}
The gift example needs a condition over classified facts. A consent example
needs a history of grants and withdrawals. A payment agreement may need parties,
assets and deadlines. A dispute about the meaning of an exception needs
alternative readings and their supporting reasons. These are related questions,
but a representation designed for one does not automatically solve the others.
The purpose of reviewing formal legal languages is to identify which structures
help make each question precise.

A formal legal language is a language in which statements about a legal or
normative domain have defined syntax and computational meaning. The word
``legal'' identifies the intended subject matter and useful abstractions. It
does not mean the language has legal authority, that a program written in it
is enforceable, or that translation into it establishes a correct reading.
Authority still comes from the relevant source and institutional decision.

For our product, a useful language must support at least three forms of
inspection. A reviewer needs to see how a condition corresponds to source
text. An engineer needs to know exactly what values and events the program
expects. A verification tool needs unambiguous semantics. A language that is
expressive enough to represent any program may make verification harder; a
language that permits only simple decision tables may fail to represent
important temporal or exception structure.

This chapter develops the relevant mechanisms through examples. It does not
choose a universal replacement for the bank's Java system. The recommended
architecture keeps a bounded executable representation, a richer interpretation
record and explicit adapters for specialized analyses. The relationship between
those representations must itself be checked. Adding more languages without
defining that relationship would create new translation boundaries rather than
redundancy.

\section{Rules do not all behave like ordinary implications}
\label{sec:languages-norms}
An ordinary material implication says that if a proposition is true, another
proposition is true. A legal norm often says that, under specified circumstances,
someone ought to do something. The existence of the duty does not establish
that the action occurred. A rule requiring a document to be supplied cannot
be used as evidence that the document was supplied. The distinction between
facts and normative outputs is central to input/output logics
\citep[secs.~2--4]{inputoutput2000}.

For a simple illustration, suppose an adopted specification says that a
distributor must provide a named warning when a condition holds. The condition
can create a duty record. Performance requires separate evidence of the
warning being delivered to the correct recipient in the required context.
If the program inserts the duty's content into its factual database as though
it had happened, every obligation can appear automatically satisfied.

Prohibitions and permissions also need care. Failure to derive a prohibition
may mean the act is permitted under one explicitly defined logic; in a bank
workflow it may instead mean the relevant sources or facts are missing. The
literature distinguishes strong and weak permissions under particular
defeasible-logical definitions \citep[secs.~2--4]{permissions2012}. Those
definitions are not a general license to treat an empty query result as
authorization. The selected RuleIR gift result remains a subcondition, with
permission supplied only by a separate approved decision policy.

Many legal rules are defeasible: they normally support a conclusion, but an
exception or stronger opposing rule can override that support. A defeasible
rule is not the same as an ordinary implication with an informal comment
attached. Its evaluation must account for applicable opposing rules and
priorities. Defeasible-logic systems formalize different ways to do this,
including rules that support conclusions and defeaters that block opposing
support \citep[sec.~2]{defeasible2000}.

The engineering implication is simple but demanding. ``Has exceptions'' is
not a complete semantics. The specification must say what happens when no
exception is known, one is established, two compete, or an exception remains
unknown. It must also say whether a priority is supplied by legal authority,
an explicit interpretation decision or a programming convention. File order
is not an adequate hidden source of legal priority.

The early logic-program treatment of the British Nationality Act demonstrated
how a selected body of legislation could be decomposed into rules with
inspectable reasoning \citep{sergot1986}. That work is valuable here as an
example of disciplined formalization, including attention to time and negative
information. It should not be read as a claim that ordinary legal prose can
always be converted automatically into a complete logic program.

\section{Controlled English: readable structure with restricted freedom}
\label{sec:languages-logicalenglish}
Ordinary English permits pronouns, elliptical references, implied subjects and
several possible attachments of a qualification. Those features make writing
compact, but they leave a translator with choices. Controlled English restricts
some of that freedom so a sentence can be mapped into logic by explicit rules.
Logical English develops this approach for law and education
\citep[language specification and examples]{logicalenglish2023}.

Consider the sentence ``If it is a gift and is linked to the product or its
type, it is prohibited unless discounted.'' The sentence is short, but ``it''
could refer to the offer or component, ``its'' needs a binding, and ``discounted''
does not identify the exception. A controlled rendering should repeat the
subject where necessary and name each definition. Repetition in the language
is useful if it removes a legally consequential ambiguity.

One possible authoring form for our own restricted profile is:
\begin{lstlisting}[caption={Illustrative LegalMath authoring notation. This is a proposed template, not claimed to be valid Logical English or Catala source.}]
subject: component C in offer O
scope:
  actor(O) satisfies distributor_profile
  activity(O) satisfies authorised_fund_promotion_profile
condition:
  C satisfies gift_definition
  AND (
    C has specific_product_link
    OR C has product_type_link
  )
  AND NOT (C satisfies fee_discount_definition)
consequence:
  selected_gift_restriction_triggered
\end{lstlisting}
The form makes four decisions inspectable: the subject, the scope, the
condition and the named consequence. Every definition identifier refers to
a versioned record. A reviewer can ask why a definition applies without
reverse-engineering the syntax tree. An engineer can translate the connectives
deterministically once the template is fixed.

Variable binding is a practical source of error. If one occurrence of $C$
refers to the voucher and another to the whole offer, the sentence has changed
its subject midway. The authoring language should enforce a declared variable
and type environment. A reference to an undeclared subject should be rejected.
An association such as ``component C belongs to offer O'' should be explicit
when the rule uses facts about both.

Negation requires similar precision. ``There is no evidence that C is a gift''
is not the same assertion as ``evidence establishes that C is not a gift.''
Logical-programming systems may distinguish explicit negation from failure to
prove a proposition. Our evidence model records known false, unknown and
conflict separately. A controlled sentence that asks about missing evidence
should use that information state explicitly; it should not rely on a model
to guess which sense of negation the author intended.

Controlled English improves the boundary between reviewed meaning and
implementation. It does not remove the earlier task of choosing the controlled
sentence from the original circular. A reviewer can still choose the wrong
definition or omit a phrase. That is why the product retains source anchors,
alternative templates and contrasting scenarios alongside the readable rule.
The language makes a chosen interpretation precise enough to challenge.

\section{Catala: defaults, exceptions and source-linked authoring}
\label{sec:languages-catala}
Catala was designed around a recurring pattern in legislation: a general rule
defines a value, and explicit exceptions may replace it. Its core calculus
gives that structure a formal meaning; its authoring approach keeps legal text
near the corresponding code \citep[secs.~2--4]{catala2021}. For a reviewer,
this can be more intelligible than a long sequence of ordinary conditional
statements whose priority is implicit in their order.

A hypothetical policy might normally require an explanation and waive a
particular step under a reviewed condition, while restoring it if the client
requests assistance. The important structure is the relationship between
the general rule and its exceptions. A nested exception can express a priority
that differs from two sibling exceptions. Moving a condition one level in
that structure can change the result even when every word remains present.

For a small teaching example, suppose a base value is false and one exception
returns true when $g$ holds. If $g$ is established, the exception supplies the
value. If $g$ is refuted, the base applies. If the specification supplies two
sibling exceptions and both are active, the system needs a conflict rule.
Catala's core treatment distinguishes an absent applicable value from a
conflict between exceptions; more than one active exception can conflict even
when the values agree. Agreement of the values does not prove that their
legal priority has been resolved.

LegalMath borrows explicit exception structure but adds its own treatment of
incomplete evidence. In the present RuleIR profile, every sibling guard is
evaluated. Two true guards conflict. One true guard with an unknown competitor
remains unknown because the competitor may become active. One true guard with
all competitors false selects its value. When all guards are false, the base
is selected. This is a separately specified design; it is not a claim that
RuleIR implements the whole Catala calculus.

The distinction matters for ensemble comparisons. A Catala model and a RuleIR
model can be useful independent representations, but they do not become
equivalent merely because both have an ``exception'' construct. The adapter
must define how missing evidence, no applicable value and competing exceptions
are represented. A comparison outside that common fragment should return
unsupported or expose the difference, not choose a convenient interpretation.

Catala's compiler work also illustrates how formal claims should be scoped.
The paper establishes properties for a core translation; the production
compiler and surrounding tooling are separate implementation objects. The
retained repository inspection also records an admitted supporting lemma in
the formalization notes. This project has not replayed those proofs. We borrow
the explicit language structure and preservation discipline, while requiring
independent evidence for any concrete adapter used in the bank product.

\citet[sec.~4.2]{huttner2022} discuss legal/programmer collaboration as a way
to keep source meaning and executable rules connected. The useful lesson is
organizational as well as technical. A legal reviewer can challenge an
exception's interpretation while an engineer checks that its computational
consequence matches that interpretation. Neither role should be expected to
certify the other's entire domain from a single code review.

\section{Metadata is part of the meaning: LegalRuleML}
\label{sec:languages-legalruleml}
A formula alone does not say which provision it interprets, which jurisdiction
it concerns, when it applies, or whether a competing formula is also
defensible. LegalRuleML addresses representation of legal rules together with
such contextual and provenance information. The work on legal interpretations
explicitly represents alternative readings associated with source material
\citep[secs.~2--4]{legalinterpretations2014}; the later OASIS specification
provides a broader standardized vocabulary \citep{legalruleml2021}.

For the fee-discount issue, two candidate rules can point to the same source
span while naming different interpretation records. One record might restrict
the exception to an initial reduction; another might include a documented
refund. Their source identity is shared, but their assumptions and consequences
differ. The system should not assign a single interpretation identifier to
both merely because the source paragraph is the same.

The representation must also preserve the authority of a selection. A legal
reviewer may adopt one interpretation for a defined product and date range.
That selection is an event with an author, reason and scope. It does not erase
the alternative or establish that the alternative was syntactically invalid.
A future source clarification may require reconsidering the choice. The
provenance record is what lets the system determine which decisions depend on it.

LegalRuleML is a representation standard, not a turnkey engine for resolving
every legal dispute. A serialized permission or priority still requires a
specified evaluation semantics. Different reasoning engines may support
different fragments or make different assumptions. For the prototype we should
borrow the discipline of explicit context and alternatives without requiring
the entire standard to become the first runtime dependency.

This leads to a useful separation in the proposed storage model. A source span
identifies words. An interpretation claim identifies an adopted or proposed
meaning. A formal rule identifies a computation. A review decision identifies
who accepted which claim under which context. Keeping those identities
separate prevents a hash of the source text from being used as if it also
identified the meaning chosen from that text.

\section{Stipula: a contract as parties, resources and transitions}
\label{sec:languages-stipula}
Stipula becomes relevant when a legal arrangement changes through actions over
time. Its original paper describes agreements, parties, fields, assets,
functions, states and scheduled events \citep[secs.~3--6]{stipula2021}. A
function can be available only to a specified party in a specified state, with
a precondition on the current data. Executing it changes fields or resources
and moves the contract to another state. This is a much richer object than
a single decision-table row.

To understand the mechanism, consider a synthetic service agreement. It begins
in a state awaiting agreement on terms. After the designated parties agree,
it enters an active state. An authorized party can request a permitted service
only while the agreement is active. A withdrawal changes the state so that
the same request is no longer available under that agreement. A timed event
may move the agreement into an expired state. This example illustrates a state
model; it is not a claim about the legal effect of any particular bank consent.

The state matters because the same action name can have different consequences
at different points in the history. A function called \texttt{requestService}
cannot be understood by its input parameters alone if its availability depends
on whether consent is active. A current Boolean copied from a stale database
may also be inadequate. The state should be derived from an attributable
history or a reviewed snapshot plus a complete subsequent event stream.

Stipula's agreement construct explicitly identifies who participates and who
must agree on the initial terms. This is useful as a modeling discipline:
the identity of the party agreeing is not an untyped annotation on a field.
For LegalMath, the analogous review question is who has authority to approve a
meaning, data mapping or release. We borrow the importance of explicit parties
and transitions, not a claim that a regulatory circular is itself a private
contract between the bank and regulator.

Assets in Stipula are distinguished from ordinary fields. An ordinary numeric
field can be copied as information; an asset transfer must account for the
resource leaving one location and arriving at another. In the original
semantics, asset movement prevents creating a negative quantity and preserves
the relevant resource through transfer. This supports questions about duplicate
spending, locked resources and the availability of contractual actions. It is
not a model of market liquidity or a substitute for the bank's general ledger.

For a synthetic escrow, suppose a contract holds a resource amount $b$ and
transfers $q$ to a recipient. A local arithmetic description is
\begin{equation}
 b'=b-q,\qquad r'=r+q,\qquad 0\leq q\leq b,
 \label{eq:asset-transfer}
\end{equation}
where $r$ is the recipient's modeled holding. Adding the first two equalities
gives $b'+r'=b+r$. This is a conservation property of the modeled transfer.
It says nothing about whether the underlying asset exists off-system or whether
the transfer is legally authorized. Those are separate premises. The example
shows how a formal language can support a useful exact property without
claiming to settle the whole legal arrangement.

Scheduled events add another dimension. A Stipula event can schedule a
continuation at a time and execute it only if the contract remains in a
specified state. An action taken before the event can change whether that
continuation is applicable. The ordering of an event and a function call at
the same time is consequently significant. The inspected original semantics
gives scheduled events precedence in the relevant equal-time situation. A bank
adapter must not assume that this is the legally correct ordering for its own
deadlines. The source, clock precision and event-order authority need review.

This example also distinguishes observation from permission. A formal contract
may describe which function calls its own execution platform permits. A bank
audit system must still record actions that occurred outside a permitted
workflow. Deleting an observed transaction because the policy would have
forbidden it would erase evidence of a possible violation. A normative monitor
therefore needs a clear separation between advice about an action and recording
its actual occurrence.

\section{A small state contract that a reader can execute by hand}

A complete miniature model helps distinguish the legal-language idea from its
implementation in a particular tool. Consider a synthetic escrow agreement in
which a payer deposits a fixed amount, a recipient may receive it after an
accepted delivery, and the payer can recover it after an expiry if delivery has
not been accepted. The example is chosen to teach states, parties, resources and
time. It is not a model of an actual bank product or a proposed interpretation of
the gift circular.

Let the contract state be one of \texttt{UNFUNDED}, \texttt{HELD},
\texttt{PAID} and \texttt{REFUNDED}. Let $b$ be the amount held by the contract,
$p$ the payer's modeled external holding and $r$ the recipient's holding. All
amounts are nonnegative integer units. Let $d$ be the agreed deposit, $\tau$ the
expiry time and $t$ the current event time. The parties and the meaning of
accepted delivery are agreed inputs to this model.

The state transition table is the executable meaning of the miniature contract.
It says which caller can perform which action, under what guard, and how the
state and resources change. It is deliberately written in a neutral notation;
it is not presented as tested Stipula source code.

\begin{longtable}{@{}P{24mm}P{33mm}P{42mm}P{48mm}@{}}
\caption{Synthetic escrow transitions, inspired by state-contract languages.}\\
\toprule
From & Action and caller & Guard & Effect and next state\\\midrule\endfirsthead
\toprule From & Action and caller & Guard & Effect and next state\\\midrule\endhead
\texttt{UNFUNDED} & Deposit by payer & $p\geq d$, $d>0$, $t<\tau$ & $p'=p-d$, $b'=d$, $r'=r$; enter \texttt{HELD}.\\
\texttt{HELD} & Release by authorized acceptance process & Accepted delivery and $t<\tau$ & $r'=r+b$, $b'=0$, $p'=p$; enter \texttt{PAID}.\\
\texttt{HELD} & Expiry event & $t\geq\tau$ & $p'=p+b$, $b'=0$, $r'=r$; enter \texttt{REFUNDED}.\\
\texttt{PAID} or \texttt{REFUNDED} & Any repeat transfer & No enabled transition & No second transfer is permitted by this model.\\
\bottomrule
\end{longtable}

Start with $p=100$, $b=0$, $r=0$, and an agreed deposit $d=40$. A valid deposit
produces $(p,b,r)=(60,40,0)$ and state \texttt{HELD}. If delivery is accepted before
expiry, release produces $(60,0,40)$ and state \texttt{PAID}. If expiry occurs
first, the refund produces $(100,0,0)$ and state \texttt{REFUNDED}. The same deposit
has two possible future paths, determined by events and their timing.

A reader can now check a conservation property without trusting a compiler. Define
$W=p+b+r$, the total modeled resource. For the deposit transition,
\begin{equation}
 W'=(p-d)+d+r=p+r=W,
 \label{eq:escrow-deposit}
\end{equation}
because the initial contract holding is zero. For release,
$W'=p+0+(r+b)=W$. For refund, $W'=(p+b)+0+r=W$. Each permitted transition therefore
preserves the total. Since the initial holdings are nonnegative and the deposit
guard requires sufficient funds, these transitions also preserve nonnegativity.
This is a local derivation about the stated miniature model.

The terminal states prevent a second transfer. After release, no transition from
\texttt{PAID} can release the same balance again. A Java implementation that
updates balances but forgets the state change could violate that property under
a repeated call. This is why the state transition is part of the contract, not
just a user-interface label. An idempotency mechanism can support repeated
requests, but must be specified consistently with the state model.

Time is equally substantive. The table uses $t<\tau$ for release and
$t\geq\tau$ for expiry. At exactly $\tau$, only expiry is enabled. If the intended
agreement allowed delivery at the deadline, the guards would need a different
boundary or an explicit ordering rule. A test performed one minute before expiry
cannot establish the exact-deadline behavior. The Stipula literature's event
ordering is therefore a modeling choice to understand, not a universal answer
for every contract \citep{stipula2021}.

The model also exposes a dependency hidden by ordinary prose: what establishes
accepted delivery? It might be an authorized person's recorded acknowledgment,
an external system event or a disputed factual assessment. The transition table
assumes that this predicate is supplied correctly. Conservation and no-double-
transfer proofs do not establish that assumption. If acceptance can be forged,
the formal transfer properties may hold while the wrong person receives the
resource.

The distinction between safety and progress is now visible. Conservation is a
safety property: no permitted step creates or destroys the modeled resource.
The ability to reach \texttt{PAID} or \texttt{REFUNDED} is a progress question.
If the platform never processes the expiry event and delivery is never accepted,
the resource can remain held. Proving eventual refund requires assumptions about
time advancement and event execution, not merely the transition table's arithmetic.
The contract-liquidity literature concerns such availability of actions and
resources; it should not be confused with market liquidity \citep{liquidity2023}.

This miniature also clarifies what a Java/JML translation would need to preserve.
A precondition specifies the caller, state and guard. A postcondition specifies
the new balances and state. An invariant states nonnegative holdings and any
conserved total. The verifier reasons about all allowed calls under those
conditions. It must also account for the event mechanism and any loop abstraction
used by the translation. The inspected Stipula-to-Java work provides a relevant
route, with the restrictions already described \citep{stipulakey2025}.

For LegalMath, the corresponding benefit is a disciplined way to model review,
consent and obligation lifecycles. A report awaiting review is not an approved
report; a withdrawn consent is not an active generation; an opened duty is not
fulfilled merely because a document was eventually delivered late. Explicit
states and transitions make these distinctions executable and testable.

The correspondence is methodological rather than a claim that every circular is
an escrow contract. Regulatory obligations can apply even when a party never
agreed to a private contract, and actual events must be recorded even if a formal
platform would reject them. A bank's monitoring system therefore separates the
normative question of what should occur from the factual record of what did
occur. The state model helps express that separation only when it is designed
for the intended purpose.

A practical specification exercise follows directly. Before implementing a new
history-dependent control, write its states, actors, events, guards and effects;
execute one normal path, one boundary path and one violating observation by hand;
identify the facts that the model assumes; and state which invariants and progress
properties matter. Then choose a language or verifier that supports those
properties. This sequence gives Stipula, Symboleo, eFLINT or a restricted custom
monitor a concrete role rather than selecting a language because its name occurs
frequently in the literature.

\section{What contract verification does and does not buy}
\label{sec:languages-contractverification}
Once a contract has a transition semantics, tools can ask whether certain
states are reachable, whether a resource can become trapped, or whether two
contracts have corresponding observable behavior. Stipula's legal-bisimulation
work gives a formal account of selected observable equivalences
\citep[behavioral equivalence section]{stipula2021}. A bisimulation relates
states so that an observable step on one side can be matched by a corresponding
step on the other. It is a statement about modeled behavior under that
definition of observation.

For the bank, an analogous question might ask whether a revised consent
implementation preserves the reviewed behavior for grants, withdrawals and
transactions. The observation set must include what matters: subject, category,
time, permitted use and recorded violations. If the comparison ignores
withdrawal, equivalence of the remaining observations would not establish the
property the bank needs. Formal observation boundaries deserve the same
scrutiny as ordinary input-domain restrictions.

\citet{stipulakey2025} demonstrate a translation from Stipula contracts into
Java with JML specifications and verification using KeY. JML expresses properties
of Java behavior, such as preconditions, postconditions and invariants; KeY
reasons about whether the Java program satisfies them. This is especially
relevant to our Java delivery target because it illustrates a route from a
specialized contract language into a general-purpose verification environment.

The inspected work has explicit restrictions. Its event conditions are modeled
symbolically, and the automated encoding keeps certain loop inputs constant
across iterations. Its examples establish feasibility for the examined models,
not complete fidelity to arbitrary calendars or external event streams. The
official translator source was inspected in the earlier project review; no KeY
proof was replayed. Adopting a similar route would require a precise mapping
from our RuleIR and event profile to the Java/JML semantics, including the
unknown and conflict results that ordinary Boolean contracts omit.

General verification cannot be promised for every expressive contract language.
\citet{reachability2026} establish undecidability results for clause reachability
in Stipula fragments and identify a decidable fragment under particular
restrictions. The restrictions concern disjoint function/event source states
and immediate events; they should not be confused with the different
disjoint-cycle condition used in the Java/JML translation work. A shared word
in two papers does not identify a shared theorem.

The practical consequence is to define bounded questions and supported
fragments. A tool can report that no bad state was found up to a specified
depth, or prove a property under a restricted transition model. It cannot
turn a finite search into a complete reachability decision for a language
whose general problem is undecidable. The report must retain the bounds so a
reviewer can understand which behaviors were left unexplored.

Stipula's amendment literature is also relevant. A contract can change its
declarations, state and permissible behavior through a specified amendment
mechanism \citep[secs.~3--5]{amending2023}. For regulatory controls, an amendment
must similarly identify whether existing obligations keep their opening source
version or migrate to a new one. Hot-replacing code without a state-migration
interpretation can change the treatment of historical events. The versioned
control and its transition obligations need separate approval.

\section{Other normative languages clarify different distinctions}
\label{sec:languages-other}
Symboleo models legal contracts using concepts such as obligations and powers,
with lifecycle distinctions that are easy to lose in a Boolean implementation
\citep[secs.~II--IV]{symboleo2020}. Creating a duty is not always the same as
activating it. A power can change a normative relationship, for example by
suspending or terminating an obligation. These distinctions are useful when a
bank workflow has review, activation, withdrawal and corrective actions with
different legal effects.

The retained Symboleo paper and software inspection do not provide a complete
implemented axiom set for LegalMath. Some supporting references in the early
paper are not fully available in the retained edition. We can use the lifecycle
questions to improve requirements, while treating any direct engine adoption
as a separate interface and semantics audit. A familiar conceptual diagram is
not a substitute for an executable contract.

eFLINT separates facts, actions and duties in a normative specification
\citep[secs.~3--5]{eflint2020}. A particularly useful distinction is that an
impermissible action may still occur, change state and create a violation.
This matches the difference between a bank's pre-transaction advice and its
post-event audit. A system that refuses to represent an unauthorized observed
action cannot explain what happened after a control failure.

The eFLINT literature also shows why language versions matter. The later
stable-model semantics provides a particular formal grounding and can produce
no model or several models \citep[translation and evaluation sections]{eflintstable2025}.
Those outcomes carry information; selecting one model silently would conceal
ambiguity or inconsistency. The newer language survey discusses implementations
whose semantic guarantees differ \citep{eflint2026}. A proof or test result for
one version should not be attributed to every interpreter bearing the same name.

s(LAW) combines legal formalization with logic programming that can explain
positive and negative outcomes and explore possible models
\citep[secs.~2--5]{slaw2024}. This makes assumptions visible in a way that a
single opaque classification may not. However, a model containing a hypothetical
fact does not establish that the fact is true for the client. Counterfactual
reasoning should propose what would change an answer, while the evidence
service establishes whether the corresponding condition actually holds.

These methods suggest an interface in which normative classification,
state transition and factual evidence are separate inputs to a decision. The
implementation should reject an unsupported duty kind explicitly. A maintenance
obligation, which must hold throughout an interval, cannot be implemented by
checking that it was true once. A non-preemptive achievement obligation cannot
be discharged by a performance before it came into force. The taxonomy in
\citet[sec.~3]{normative2016} helps identify these choices; the precise profile
still needs a source-grounded interpretation.

\section{Decision tables and business-rule engines}
\label{sec:languages-tables}
Decision tables remain useful when their rows and hit policy have a precise
meaning. A table can show, for example, that a known discount exception changes
the selected trigger while an unknown exception leaves a material question.
The reader must know what happens when several rows match: choose the first,
require a unique row, collect results or apply a priority. DMN makes such hit
policies explicit \citep[sec.~8.2.11]{dmn2024}.

That explicitness is the valuable part. A spreadsheet whose row order silently
resolves competing legal exceptions is harder to review because its legal
priority is hidden in presentation. Changing the sort order can change the
answer. A unique-hit table can reject overlap, but it cannot determine whether
the missing row represents an omitted legal situation. Independent source
coverage and boundary cases remain necessary.

Drools, OPA and other rule engines can execute policy logic and integrate with
enterprise applications. The question for this product is not whether they
are capable software. It is whether the selected engine's semantics, evidence
model, explanation behavior and deployment interface match the adopted
specification. Unknowns, conflicts, exception priorities and event ordering
need explicit mappings. A backend name is not an assurance claim.

For the prototype, a small RuleIR is easier to specify and compare than a
large engine with many unsupported features. Other engines can serve as
independent reference encodings for a supported fragment after that fragment
is defined. The comparison must include the behavior that ordinary Boolean
examples leave out. Engine diversity without semantic alignment can produce
many discrepancies that reflect adapter choices rather than flaws in the
legal interpretation.

\section{A solver can improve reasoning after a translation}
\label{sec:languages-neurosymbolic}
A neural-symbolic system separates tasks that language models and formal
engines perform differently. A language model proposes a formal representation
of prose. A symbolic engine evaluates or reasons about that representation.
The engine can give a reproducible answer to a precise logical question and
can expose inconsistency, unsupported syntax or a counterexample. The
translation that supplied its premises remains a separate source of error.

Logic-LM uses language-model formulation followed by an appropriate solver and
feedback-driven refinement \citep[sec.~3 and prompt appendices]{logiclm2023}.
LINC combines language-model formalization with first-order theorem proving
and analyzes failures at the translation and reasoning boundary
\citep[method and error analysis]{linc2023}. These methods support the design
choice that a model should not be asked to simulate every logical calculation
in prose. They do not support treating successful solver execution as evidence
that a legal requirement was fully captured.

A syntax repair illustrates the benefit and the limit. Suppose a candidate
uses an undeclared symbol in the gift rule. The parser identifies the symbol
and rejects the encoding. A repair prompt can ask for the missing declaration
or corrected reference while preserving the intended interpretation. The
repaired candidate is then validated again. This process may solve a concrete
representation error. It does not justify adding a new legal assumption
merely because doing so makes the formula easier to satisfy.

ARc develops policy-model creation, human vetting, structured inspection and
symbolic checking in a combined system \citep[sec.~3 and appendix A]{arc2025}.
Its policy creation procedure divides a document into units, proposes types,
variables and constraints, and composes them into a model. Its vetting
mechanisms include linting, template-based English inspection and testing.
These are relevant to our product because they make the proposed model
available for correction before individual answers rely on it.

Composition creates its own interpretation problem. Two units may use the same
word for different legal concepts, or different words for the same concept.
ARc uses embedding-based clustering in its variable-unification procedure.
For our high-stakes specification, such a merge should be a proposal with
provenance and validation. Merging ``offer discount'' with ``component
discount'' because their descriptions are similar could remove the distinction
that g04 was designed to test. A similarity score cannot establish identity
of legal concepts.

ARc's answer-verification procedure also uses redundant translations and
symbolic agreement checks. Its reported finding categories distinguish
ambiguous translation, impossible premises, invalid conclusions, entailed
conclusions and satisfiable but unentailed claims. The distinction is useful:
if a model permits both an outcome and its opposite under incomplete premises,
the system can show scenarios for each rather than invent a definitive answer.
For LegalMath, these scenarios become review questions tied to the candidate
set and the source context.

The word ``soundness'' in ARc's empirical evaluation needs careful translation
into a product claim. The paper defines an empirical quantity based on false
positives across requests; it is not a theorem that the English policy was
formalized correctly. The denominator differs from precision among approvals.
The manuscript therefore borrows the inspection and disagreement mechanisms
without importing a claimed legal-error probability. Chapter~\ref{ch:evaluation}
defines the denominators our own evaluation would need.

The tax-reasoning work of \citet{jurayj2026} further illustrates the importance
of task boundaries. Supplying gold rules, retrieving annotated cases and
extracting facts are different assistance conditions from generating rules
from a new circular. A comparison must account for those inputs and the human
effort that produced them. A system given a reviewed rule library should not
be presented as evidence that the same system can independently interpret
unreviewed prose.

\section{Roundtrip verification and targeted repair}
\label{sec:languages-roundtrip}
Roundtrip verification asks whether a formal representation survives a journey
through another representation. In the method of \citet{roundtrip2026}, source
English is formalized, the formula is translated back into English, and the
reconstructed English is formalized again. A solver compares the first and
last formulas. If they differ, diagnosis identifies a likely translation
stage to repair, and dependent stages are regenerated.

This is attractive for our user interface because the intermediate English
can make a formal candidate easier to inspect. Suppose the original formula
contains both product-link routes but the back-translation mentions only a
specific product. Re-formalizing that sentence may lose the product-type
route. The resulting counterexample exposes a defect in the explanation or
translation chain. A reviewer should not be shown that explanation as if it
faithfully described the formal rule.

The same construction can miss a shared omission. If the first formalization
already omits the product-type route, a faithful back-translation of that
incorrect formula and a faithful re-formalization will agree. The roundtrip
then passes. What survived is the first formula, including its error relative
to the source. The paper explicitly recognizes this limit; our ensemble must
retain a source-grounding path that does not simply inspect the same
back-translation.

Repair therefore needs a named target. If the formal rule is justified but
the explanation drops a condition, repair the explanation and its dependent
roundtrip outputs. If the original rule omits a source condition, create a
new interpretation candidate and revalidate its downstream artifacts. If the
two formulas use incompatible vocabularies, repair the mapping before claiming
the formulas disagree about one question. A generic ``make these consistent''
instruction does not distinguish those cases.

The stopping rule is equally important. A roundtrip procedure can run for a
bounded number of attempts and still end in disagreement. The product should
report the original and revised formulas, the counterexample, diagnosis,
actions attempted and the remaining issue. It should not keep repairing
until some pair agrees and then discard the path that made the original
source objection visible. The ensemble controller extends this principle
from one translation chain to several methods and issue types.

An explanation generated by deterministic templates provides a useful
additional path. It is constrained by the formal syntax and can be checked
for completeness of its displayed operands. It may be less elegant than
free prose, but it offers a stable view of what the program actually says.
A freely generated explanation can still help readers if it is compared
against that view and does not replace it as the sole account of the rule.

\section{Evidence spans and distant qualifications}
\label{sec:languages-evidence}
An answer supported by a source span is easier to review than an unsupported
assertion, but the span itself may be insufficient. A qualification can appear
in a footnote, definition or distant section. ContractNLI studies document
inference with evidence annotation and highlights challenges involving
references and exceptions \citep[task, annotation and error-analysis sections]{contractnli2021}.
For our product, the relevant lesson is to retain the supporting and limiting
passages needed for the inference, rather than a single convenient quotation.

The distinction between ``not mentioned'' and ``false'' is particularly useful.
If a retained source does not settle the definition of a fee refund, the
absence should remain visible. A retrieval engine returning no matching
paragraph does not establish that the exception excludes refunds. It may have
searched the wrong document, missed a synonym or lacked the referenced source.
The evidence report should distinguish unsuccessful retrieval from reviewed
absence within a specified source set.

Source spans also require stable identity. Character offsets depend on the
extraction version, while page numbers may not survive a changed layout.
Retain raw-byte identity, derivative identity, exact span text or hash, and
human-readable location. A reference to paragraph 10 without a source version
can become misleading after an amendment. A matching text hash without its
source context may identify a sentence copied into a secondary summary.

The formalization boundary should distinguish source claims from supplied
facts. A sentence defining a gift exception is part of the normative context.
A bank record establishing that a fee was paid is factual evidence. A
reviewer's classification relating the record to the exception is an
interpretive application. Combining all three into an undifferentiated
``evidence'' list makes it harder to see where disagreement actually arises.

Recent faithfulness work analyzes how language models and formal solvers can
diverge when implicit assumptions enter formalization \citep{limits2026}.
The retained paper also has annotation and denominator questions that prevent
us from transferring a general accuracy figure. Its examples nevertheless
reinforce a practical requirement: an implicit assumption that changes the
answer should become an explicit, challengeable premise. The system must
show where a model supplied information that the source packet did not.

\section{Testing paths, dates and inconsistent constraints}
\label{sec:languages-analysis}
Other methods contribute different checks once a formal model exists.
CUTECat uses concolic execution to explore paths in Catala programs
\citep[secs.~2--5]{cutecat2025}. Concolic execution combines a concrete run
with symbolic constraints describing the path taken. By changing a path
condition, the tool seeks an input that exercises another branch. For our
gift example, this is analogous to moving from a named-product case to the
case where only the product-type route can determine the answer.

Path-directed generation can expose an exception branch that ordinary
examples never visit. It cannot visit a branch that the formal program
omitted entirely. Independent source-derived cases remain necessary. The
reported costs on larger legal examples also show why a prototype needs
explicit resource limits; small-program speed should not be treated as a
promise that every circular can be exhaustively explored.

Date arithmetic is another place where a precise program can embody a wrong
choice. Adding a month to a date near the end of a month may require a
rounding convention. The choice can affect deadlines. \citet{dates2024}
formalize date arithmetic and analyze ambiguity in legal applications. The
transfer to this project is to make calendar rules, boundary conventions and
their source basis explicit. A verified date library does not decide which
rounding convention a particular provision requires.

ContractCheck analyzes consistency and executability of structured contract
models \citep[secs.~3--8]{contractcheck2025}. A satisfying assignment shows
that a modeled set of conditions can jointly hold. An unsatisfiable set can
help locate conflicting clauses or assumptions. Neither result proves that
the original legal document has been translated faithfully. The study's
manual translation is an important part of its task boundary.

An unsatisfiable core is a useful diagnostic: a subset of constraints that
already suffices for inconsistency. It is not necessarily the smallest such
subset unless a minimization procedure establishes that property. For a
reviewer, a small and source-linked explanation is valuable, but the product
should label a returned core accurately. It must also distinguish an
encoding inconsistency from legitimate defeasible conflict that the chosen
logic was intended to resolve.

These analyses fit into the ensemble as specialized questions. Does this
representation have a hidden path? Does a date convention alter a boundary
case? Are the supplied assumptions jointly possible? Does a revised rule
preserve the reviewed behavior? Their evidence can challenge a candidate
without claiming to replace the source interpretation process.

\section{A synthesis for the first implementation}
\label{sec:languages-synthesis}
The immediate product should keep a small executable core whose semantics are
fully specified. Around it, the interpretation record should be richer than
the executable language: it needs alternative readings, source versions,
definitions, unresolved qualitative judgments and supporting arguments.
Controlled templates give reviewers a readable view of formal structure.
Specialized analyses can then inspect declared fragments through adapters.

This choice avoids forcing every legal question into the first compiler.
A paragraph requiring judgment can become a reviewed assessment input or
manual procedure, with explicit evidence and responsibility. A temporal duty
can use an event profile whose limits are stated. A disputed definition can
remain in the interpretation set until review resolves it. The system can
still generate Java for a reviewed subset while keeping the broader
assessment visibly incomplete.

Every adapter needs a semantic contract. It declares supported constructs,
result categories, unknown and conflict behavior, vocabulary mapping,
resource bounds and how a returned witness is replayed. A paper's name or
a repository's existence does not supply that contract. A failed or
unsupported analysis is reported as such and cannot be counted as an
additional agreeing ensemble member.

The literature thus supplies complementary mechanisms: source-linked
authoring, precise exceptions, contextual alternatives, normative state,
symbolic consequence, path-directed testing and translation repair. The
new contribution proposed here is the controller that combines these
mechanisms without allowing one successful check to conceal another
unresolved question. Before specifying that controller, we need to explain
how competing interpretations are generated and defended. That is the task
of the next chapter.

\chapter{Constructing and comparing competing interpretations}
\label{ch:search}

\section{The object of search is a legal choice}

Suppose three models explain paragraph 10 in different words. The first says that
product-linked gifts are prohibited unless they are fee discounts. The second
calls the same arrangement a marketing incentive and repeats the exception. The
third turns the sentence into an identical decision table. We have three
presentations but possibly only one interpretation. Their agreement tells us
little about the alternative that a mixed offer should be assessed component by
component, or about the evidence required to classify a cashback payment. An
ensemble that rewards the number of agreeing explanations may therefore become
more confident without becoming more informative.

A useful hypothesis changes a consequential choice. It might change the unit of
assessment from a whole offer to each benefit, distinguish reimbursement of a
paid charge from an unrelated reward, or require an incorporated definition
before applying a term. The hypothesis must also explain why that choice is
supported. Merely deleting the words about a product type creates a different
program, but an unsupported deletion is a test mutation, not an equally
respectable interpretation. The system needs both kinds of variation and must
label their roles accurately.

We will use a candidate record with five substantive parts: a proposition about
the source; the assumptions under which it holds; a formal rule that expresses
its consequences; supporting and opposing reasons; and a description of cases
that distinguish it from another candidate. The candidate also carries source
versions and an immutable identifier. This makes a claim such as ``this cashback
is a discount'' answerable: a reviewer can inspect the proposed definition,
identify the facts it requires, and observe the transactions on which a competing
definition would disagree.

The search objective follows from that representation. We seek a useful set of
supported alternatives, together with reasons for rejecting other proposals and
questions that remain unanswered. Finding a single high-scoring sentence is a
different objective. The former fits the bank's need to know where judgment is
required; the latter can conceal the judgment in an optimization procedure.

\section{Legal theories are more than predicted outcomes}

Theories of legal reasoning provide a richer starting point than generic text
search. \citet[sections 3--5]{theories2003} represent a theory through cases,
legally relevant factors, rules, preferences between rules, and values that
explain the preferences. A factor is a legally significant characteristic, not
simply a word that occurs in a judgment. A rule relates such characteristics to a
conclusion. A preference explains why one applicable rule should prevail over
another when both cannot determine the outcome together.

To understand the distinction, consider an invented training example about a
reimbursement policy. One decision allows reimbursement because the expense was
necessary for an authorized task. Another refuses it because the claimant had
already been reimbursed. A theory containing only the first rule predicts payment
whenever necessity is established. A theory containing both rules, but no
priority, encounters a conflict in the double-payment case. A theory in which the
bar on double payment defeats the necessity rule explains both outcomes. This is
a logical illustration, not a statement about an actual court or SFC policy.

The theory's values might explain a preference by reference to preventing double
recovery. That explanation is itself contestable: it must be connected to
applicable authority rather than invented as the bank's preferred commercial
outcome. The theoretical machinery helps represent such an explanation; it does
not decide which values an institution is legally entitled to prioritize.

For the circular, the analogous factors include a product link, a type-of-product
link, the character of the benefit, and evidence of a charge reduction. A theory
might treat an incentive's relationship to the subscription charge as decisive;
another might require that the charge itself is reduced in the contractual fee
schedule. Their predictions may coincide for a direct fee waiver and diverge for
cash paid after settlement. The divergence identifies an interpretive question
that a bare outcome label would obscure.

The source paper does not claim to recover complete judicial reasoning from raw
opinions. Section 5.4 explicitly discusses limits of the representation. Our
implementation must preserve that limit: converting a judgment to factors is
another interpretive task, with evidence spans, assumptions and review. A model
that extracts a convenient factor list and then proves a conclusion from it has
not bypassed the original source problem.

\section{What case-based argument contributes}

The history surveyed by \citet[sections 2--3]{hypolegacy2017} is particularly
relevant to the user's concern about a single reading. HYPO organizes exchanges
in which one side cites a favorable case, an opponent distinguishes it or presents
a counterexample, and the first side replies. This exchange encourages a search
for facts that make an apparent analogy fail. The purpose is not to collect the
largest number of vaguely similar documents.

Imagine that an earlier, applicable authority addresses a benefit described as a
fee rebate. A candidate relying on it must state the shared features: perhaps a
charge was first incurred, the later payment could not exceed that charge, and
both arose from the same transaction. An opposing candidate may distinguish an
offer in which the payment exceeds any charge or can be earned without incurring
one. These are questions about material similarities. A search engine's textual
similarity score does not answer them.

HYPO's claims lattice orders relationships between cases using legally relevant
dimensions. A lattice is a partial ordering: some comparisons remain
incomparable. Two precedents can each support a different aspect of the present
case, without either being globally ``more similar.'' Preserving this structure
is useful for an interpretation workbench. A numerical embedding can nominate
cases for inspection, while the argument record explains which feature makes a
case relevant and which feature might defeat the analogy.

CABARET, as described in the retrospective, combines statutory rules with
case-based reasoning about difficult terms. BankXX searches a network of pieces
of legal argument. These suggest a division of labor for the circular: apply
explicit logical structure directly, and investigate open-textured classifications
through sources and arguments. They do not justify allowing analogy to override
an unambiguous express restriction without an applicable legal reason. The
original CABARET and BankXX implementations have not been audited for this
monograph; the historical claims here come from the inspected retrospective.

A production case record therefore needs more than a citation and a summary.
Record the issuing body, jurisdiction, date, procedural context, proposition
relied on, relevant factual features, passage that supports the proposition, and
known treatment by later authority. Distinguish a court's holding from an
advocate's submission or an incidental observation. Those classifications must
remain reviewable. An invented or misclassified precedent is more dangerous when
attached to a polished formal argument than when presented as an uncertain search
result.

This book supplies no actual case-law ruling on the cashback example. Its role is
to specify how such authority would enter the system if found. The immediate
23EC46 exercise remains tied to the circular and its unresolved incorporated
materials. A reviewer should never have to infer that a hypothetical analogy is
a reported judicial decision.

\section{AGATHA: explicit construction of rival theories}

AGATHA supplies a direct precedent for searching over competing legal theories.
\citet[sections 2--4]{agatha2005} define moves that extend or challenge a developing
theory: analogy, distinction, counterexample and changes to preferences. Rather
than asking a model to ``think harder,'' the system asks what legal move would
change the current account and why. This makes the search history inspectable.

In our running example, an analogy constructor could propose that two payments
belong to the same classification because both reverse a charge. A distinction
constructor could split them by whether a charge was actually paid. An exception
constructor could attach the fee-discount qualification at a particular point in
the rule. A counterexample constructor could search for a scenario that the
candidate treats as permitted despite an express product-type link. Each move
should identify its parent, the changed proposition and the source support.

AGATHA evaluates theories using explanatory power, simplicity and completeness.
Its search is a modified A* procedure, and the paper explains why the modification
can miss the best theory. We should not import the word A* as an optimality
certificate. The evaluated cases already have factor descriptions. The reported
exercise is therefore substantially easier than discovering the factors,
resolving cross-references and determining current authority in raw banking
materials.

The paper's adversarial lookahead is also instructive. A theory can look strong
because a limited opponent fails to find the appropriate counterargument.
Optimizing for a victory against that opponent can favor a weak justification.
For a bank, the objective must be resistance to informed challenge, not persuasion
of a convenient critic. The system should preserve the strongest located
objection even when an advocate agent prefers to stop discussing it.

We borrow the explicit construction operators and the record of competing
reasons. We do not adopt the paper's heuristic weights as calibrated legal
confidence, nor its advocacy objective as the bank's release criterion. A
successful local adaptation must be evaluated on whether it finds material
alternatives and omissions within a declared source perimeter. That evaluation
has not yet been performed.

\section{Arguments, attacks and unresolved conclusions}

An argument is a chain from premises to a conclusion. Some steps are strict:
under the stipulated meaning of conjunction, two true components make the
conjunction true. Other steps are defeasible: evidence that a payment reverses a
charge may support a classification while remaining vulnerable to a contrary
source or missing condition. ASPIC+ separates these kinds of inference and
formalizes ways to challenge them \citep[sections 2--3]{aspic2010}.

There are three useful attacks. An attack on a premise disputes a starting fact,
such as whether the charge was paid. A rebuttal supports a contrary conclusion,
such as classifying the payment as an additional benefit. An undercut attacks the
inference itself: even if the payment and charge have equal amounts, equality of
amount alone may not establish that one legally discounts the other. These
attacks require different remedies. Retrieving a payment record may settle the
first; retrieving a definition may address the third. Repeating the same
classification prompt need not address either.

Let an argument graph contain arguments $a$ and $b$ that defeat each other, with
no stronger argument resolving the dispute. Under the usual preferred-extension
semantics, the graph has two alternative acceptable sets, $\{a\}$ and $\{b\}$.
Neither argument appears in both sets. Under grounded semantics, neither is
accepted in this simple graph. The result expresses unresolved competition; it
does not say both conclusions are legally true.

An extension is a set of arguments satisfying a specified acceptability
condition. A conclusion is credulously accepted when at least one extension
supports it, and skeptically accepted when every relevant extension supports it.
A conclusion may be supported by different arguments in different extensions;
requiring one identical argument to appear everywhere would be a stronger and
different condition. The selected semantics and treatment of premises must be
stored alongside the result.

There is an implementation trap at this point. If a semantics yields no
extensions, code that applies a universal predicate to an empty list may report
that every proposition passes. The inspected Carneades Dung implementation
illustrates why callers must handle the existence of extensions explicitly. The
workbench will require a nonempty computed collection before reporting skeptical
support. If no extension exists, or enumeration is incomplete, the result is a
separate unresolved status. It cannot become a positive conclusion through
vacuous truth.

Argumentation computes consequences of the supplied arguments, attacks and
preferences. It cannot establish that all relevant arguments were supplied or
that the legal preferences are justified. Moreover, ASPIC+ rationality results
depend on conditions on rules and knowledge bases; the author's correction to
definition 6.8 must accompany any implementation relying on those results. A
free-form graph of model-generated sentences does not automatically inherit the
paper's theorems.

\section{Burdens, assumptions and evidence in Carneades}

Carneades distinguishes ordinary premises, assumptions and exceptions
\citep[sections 3--6]{carneades2007}. An ordinary premise needs support. An
assumption may stand in a particular dialogue until challenged. An exception can
defeat an otherwise supported argument. These distinctions explain why absence
of evidence has different effects in different reasoning arrangements.

For example, the workbench might require affirmative evidence that a transaction
falls within the bank's activity profile. It must not silently assume that fact
because most transactions in a training set did. Conversely, a procedural review
may temporarily accept an explicitly declared factual assumption to compare the
consequences of two legal readings. That temporary acceptance is a property of
the review exercise, not a permission to use the assumed fact in production.

A burden of proof specifies what support is needed for a proposition in a given
context. The 2007 paper includes illustrative proof standards and explicitly
notes that they have not been validated as legal proof standards. It would be
wrong to rename one of its numerical or dialectical conditions ``SFC approval''
and install it as policy. The bank needs its own reviewed rules about which
questions require legal adjudication, what factual evidence is sufficient, and
which roles can accept residual operational risk.

The original model is acyclic. Newer Carneades software has a different
computational model and cannot be treated as an executable copy of every result
in the paper. Our source inspection of the newer implementation found useful
argument utilities but also potentially expensive extension enumeration. A
prototype should first support small, bounded graphs with explicit semantics and
known reference examples. Introducing a general argumentation engine before its
resource and semantic behavior are understood would create another opaque
component in the assurance chain.

For the first release, even a simple structured argument register is valuable.
Require each candidate to name its supporting premises, distinguish assumptions
from evidence, identify opponents and specify what would defeat it. This can be
implemented before automatic extension computation. The eventual engine then
formalizes an already visible review practice rather than replacing an undefined
practice with a numerical score.

\section{Breadth before expensive investigation}

The initial search should cover named dimensions: actor and activity scope,
subject identity, definitions, logical attachment of exceptions, modality,
quantification, time and incorporated context. These dimensions are prompts for
possible differences, not a Cartesian product of equally plausible answers. If a
choice has no source support, retain it as unsupported or as a negative control;
do not inflate the shortlist with implausible alternatives.

Consider three dimensions with two alternatives each. Exhaustive combination
would generate eight candidates, but some combinations may be inconsistent. A
whole-offer interpretation combined with a component-only factual schema may not
be executable without a transformation. A candidate that uses a definition from
one historical version and an exception from an incompatible later version may
not describe any applicable regime. Rejecting these combinations requires an
explicit reason, because an encoding defect could otherwise be mistaken for a
refutation of the underlying reading.

A practical initial frontier reserves room for each material family before
ranking refinements within a family. The reservation is a search policy, not a
claim that all families have equal legal merit. If the budget permits only six
candidates and eight supported families are identified, the report must say that
two families remain unexplored. It should not display a six-candidate shortlist
as exhaustive. Increasing a beam width can improve exploration, but no finite
width establishes semantic completeness for unrestricted legal language.

The search tree records how candidates were created. The argument graph records
shared dependencies. These are different structures. Two branches may rely on
the same disputed definition, and a later clarification should affect both.
Representing all dependencies only as parent-child edges would duplicate facts
and make invalidation unreliable. A directed acyclic graph of immutable candidate
revisions can coexist with a possibly cyclic argument graph, provided each
structure has its own semantics and resource limits.

Deduplication also needs layers. Exact duplicate bytes can be merged by hash.
Canonical syntax can identify expressions that differ only in irrelevant
formatting. Formal equivalence can establish equal decisions over a declared
domain. None of these establishes that source assumptions are identical. Two
candidates with equivalent Boolean bodies but different definitions of ``gift''
remain separate interpretive commitments. Their shared body may be stored once,
while their source and factual bindings remain distinct.

\section{Tree search and the meaning of a reward}

Tree of Thoughts separates candidate generation, evaluation and search
\citep[section 3]{treeofthoughts2023}. Its inspected breadth-first implementation
retains a beam selected by model-based values or votes. That is bounded search,
not exhaustive breadth-first enumeration. String deduplication in the code does
not ensure that the retained thoughts express different legal hypotheses.

Language Agent Tree Search combines tree search, model judgments, reflection and
environment feedback \citep[sections 3.2--4 and appendix A]{lats2024}. A commonly
used selection expression illustrates the tradeoff. For a parent visited $N$
times and child $i$ visited $n_i>0$ times, let $\bar r_i$ be its observed mean
reward. A UCT-style selector evaluates
\begin{equation}
 U_i=\bar r_i+c\sqrt{\frac{\log N}{n_i}},\qquad c>0.
 \label{eq:uct}
\end{equation}
The first term favors a child that has produced useful outcomes. The second
encourages investigation of less-visited children. A separate rule handles
unvisited children; setting their denominator to zero is not an implementation.
The constant $c$, reward definition and visit accounting are material design
choices and must be recorded.

For a game, the reward can be a terminal win. For a software task, a reward may
include tests. For interpretation, neither model agreement nor passing a
candidate's own tests is an external truth signal. A useful action reward could
instead record a verified dependency retrieved, a replayable disagreement
scenario found, or a previously open question resolved by authenticated evidence.
Such a reward measures progress in an investigation. It must not be displayed as
the probability that the resulting legal reading is correct.

The distinction is especially important when reflection uses the model's own
explanation of a failure. A model can learn to satisfy a flawed test or to agree
with its earlier answer. The resolution controller must keep the authoritative
question and evidence fixed while allowing the candidate to change. The oracle
cannot be rewritten merely to improve the candidate's reward.

For the first prototype, a deterministic scheduler with family coverage and
explicit action priorities is easier to validate than a learned MCTS policy.
MCTS can later be an optional scheduler under the same action, budget and release
constraints. A comparison would ask whether it finds more independently judged
material issues at equal cost. It would not ask whether it accumulates a larger
self-assigned score. Chapter~\ref{ch:evaluation} specifies the needed experiment.

\section{Choose questions that separate the surviving readings}

Suppose four supported candidates remain. Three agree on a direct fee waiver,
but they differ on a payment larger than the fee, a payment made without any fee,
and a mixed voucher-and-waiver offer. Rechecking the direct waiver provides
little discrimination. The system should identify the question whose answer is
most likely to separate the remaining positions.

Generalized binary search formalizes a related problem for a finite hypothesis
class and an answer oracle \citep[Figure 1 and sections I--II]{gbs2011}. Its
identification guarantees have assumptions about the hypothesis space, queries
and oracle. Legal interpretation may violate all three: a defensible reading may
not yet be generated, a question may not admit a definitive answer, and the
reviewer can be uncertain. We borrow question selection without claiming those
guarantees.

A transparent heuristic uses pairwise disagreement. Let $H$ be the retained
candidates, $w_h\geq0$ descriptive priority weights, and $o_h(x)$ the candidate's
predicted outcome for a proposed factual scenario $x$. Define
\begin{equation}
 D(x)=\sum_{\{h,j\}\subseteq H}w_h w_j\,
       \mathbf{1}\{o_h(x)\ne o_j(x)\}.
 \label{eq:disagreement-question}
\end{equation}
Each unordered pair contributes when the scenario makes its predictions differ.
With equal weights, this simply counts separated pairs. Dividing by the expected
cost of answering the question can help scheduling, but the cost estimate is
itself uncertain. Neither version converts the weights into posterior
probabilities.

For two candidates, the useful scenario is any feasible one on which they differ.
An SMT solver can sometimes construct such a scenario, subject to the domain
conditions explained in Chapter~\ref{ch:proof}. For classification disagreements,
the solver may expose a variable assignment but cannot manufacture evidence that
makes it legally meaningful. A witness with a ``discount'' flag set to true is not
an answer to the question of what establishes that flag.

The review question should therefore contain the facts, the rival conclusions,
the source passage responsible for the disagreement, and the additional evidence
or authority sought. ``Which answer is right?'' is usually too vague. ``Does the
incorporated definition require a reduction in the contractual charge, or can a
later reimbursement qualify when it is capped at the paid charge?'' identifies
the missing premise. If the available authority does not resolve it, that is a
substantive result that the report should preserve.

\section{Uncertainty about meaning is not a vote count}

Self-consistency samples multiple reasoning paths and aggregates their answers
\citep[section 2]{selfconsistency2023}. Its evaluations show benefits on the
studied reasoning benchmarks. The method deliberately distinguishes diversity of
paths from a single greedy trajectory, but the samples still share a model and
prompting context. Forty sampled answers are not forty independent legal experts.
The paper does not establish omission detection or calibrated legal confidence
for SFC circulars.

Semantic uncertainty goes further by grouping answers by meaning before
estimating uncertainty \citep[sections 3--4]{semantic2023}. This is helpful because
ten paraphrases of one answer should not look like ten separate answers. In the
idealized formulation, strings expressing the same meaning belong to a class
$c$, and the probability of the class is the sum of the probabilities of its
strings. The class entropy is then
\begin{equation}
 H(C)=-\sum_c p(c)\log p(c).
 \label{eq:semantic-entropy}
\end{equation}
A concentration on one class produces lower entropy than an even distribution
across several classes. This measures uncertainty of the model's answer
distribution under the selected grouping. It is not a proof that the concentrated
answer is correct.

The paper approximates semantic grouping with natural-language inference. Its
appendix discusses nontransitivity and sensitivity to clustering order. Our
inspection of the official source identified a further implementation issue: the
examined clustering condition tests the absence of label zero in either direction,
rather than explicitly requiring entailment in both directions. The retrieved model configuration maps label zero to contradiction, one to
neutral and two to entailment: the inspected code therefore also permits neutral
relations \citep{semanticcode2023,debertaconfig}. The model configuration was retrieved separately from the pinned code;
its identity must be pinned as well before reuse. Even a
perfect implementation of the stated mutual-entailment test would remain an
empirical language-model judgment, not formal equivalence of bank controls.

The workbench can use semantic similarity to propose clusters for review. It
must not delete a dissenting candidate merely because an NLI model assigns a
high score. Formal equivalence can justify merging execution paths only within
its checked domain, while source assumptions and classification definitions stay
visible. Chapter~\ref{ch:evaluation} develops calibration and set-valued reporting
without assuming that a model's internal distribution is a distribution over
legal truth.

\section{What can be reused from the adjacent projects}

DynareMCP contains useful patterns for explicit hypotheses, diagnostics,
falsifiers and limits of claims. In the inspected legacy audit writer,
\srcpath{write_hypothesis_node} records a hypothesis, source support, code support,
a diagnostic, a falsifier and a supplied rank score. These fields map naturally
to interpretation candidates. The supplied score is not automatically calibrated;
its meaning depends on the caller.

The inspected legacy scheduler chooses among eligible frontier nodes by a stored
priority and identifier. It is not a general UCT engine hidden behind the word
``tree.'' Reusing its bookkeeping ideas is sensible; importing it and claiming
that a proven Monte Carlo search now verifies legal meaning would be wrong. The
new controller needs its own issue identity, source lineage, legal review states
and resource accounting.

MathDevMCP is useful once a proposition is precise enough to prove. A formal
obligation might concern preservation of a Boolean expression, absence of an
unbounded loop in a finite controller model, or equivalence of two encodings over
a declared domain. ResearchAssistant supports retrieval, parsing and retention
of the literature and regulatory evidence. Neither project supplies a privileged
interpretation of English. They strengthen different links in the chain.

The design that follows therefore combines several structures: a source inventory
to detect missing material, a candidate set to expose choices, an argument graph
to record reasons, a bounded search to investigate questions, and formal checks
to compare executable consequences. The diversity is useful because each
structure can reveal mistakes that the others overlook. Its value will depend on
how discrepancies move between them, which is the central subject of the next
chapter.

\chapter[An accountable ensemble with bounded resolution]{An accountable ensemble\\with bounded resolution}
\label{ch:ensemble}

\section{The design objective and its limits}

The proposed system must make a dangerous path difficult: a plausible reading is
accepted, translated correctly, and released without anyone discovering that it
omits a material qualification. Multiple generators alone do not close that path.
They may share a source extraction, a vocabulary, a training distribution and the
same mistaken assumption. The design therefore uses redundancy across different
activities, not just repeated answers to one question.

One activity inventories the source and its dependencies. Another identifies
normative statements. Several candidate constructors express different readings.
An opponent seeks a missing condition or counterexample. A formal evaluator
checks consequences. A separate reference process supplies expected decisions
for reviewed scenarios. A Java verifier checks preservation. Each activity has a
defined responsibility and a defined way to fail. A failure or discrepancy is
recorded even when every other member appears successful.

The main output is an interpretation dossier. It contains the proposed control,
the source it covers, the readings considered, the reasons for their dispositions,
the tests and formal checks, and the uncertainty that remains. The controller
may prepare this dossier for review. It does not acquire authority to approve
legal meaning merely because it has completed its work.

There are three success conditions at different levels. An engineering run
succeeds when it follows the specified state machine and produces a complete
report within its limits. An interpretation investigation succeeds when it
resolves its declared questions with adequate evidence or precisely identifies
what requires judgment. A release succeeds only after the required people approve
the exact interpretation and executable package. A run can satisfy the first
condition while ending with an unresolved legal question. That is a valid and
useful outcome, not an excuse to fabricate certainty.

No finite ensemble can establish that it has found every possible reading of
unrestricted language. The operational objective is more precise: cover a
reviewed source perimeter; require explicit treatment of every inventoried unit;
search a declared set of interpretive dimensions; investigate every detected
discrepancy under finite limits; and prevent a material unresolved question from
being silently converted into release. These obligations can be implemented and
tested. Their empirical effectiveness at discovering unknown omissions requires
the evaluation in Chapter~\ref{ch:evaluation}.

\section{Why diversity has to be designed}

Repeated sampling can reduce some random mistakes while leaving a shared mistake
untouched. A simple probability model makes the distinction clear. Let $q$ be the
probability of a common failure that makes every member wrong. Conditional on
that failure not occurring, assume for this illustration that each of $m$ members
has independent error probability $p$. The probability that all members are wrong
is then
\begin{equation}
 \Pr(\text{all wrong})=q+(1-q)p^m.
 \label{eq:common-error}
\end{equation}
The two events are disjoint: either the common failure occurs, or it does not and
all residual errors occur together. This derivation is an illustrative model,
not an estimated error law for language models or lawyers.

If $q=0.02$, $p=0.1$ and $m=3$, the expression gives $0.02098$. Treating the three
members as unconditionally independent would instead give $0.001$. Adding
members reduces the residual term in this example, but the common failure
remains. A missing footnote in the shared extraction is a concrete way that the
common failure could arise. Three careful interpretations of the same incomplete
text may agree perfectly.

A second calculation concerns the variability of a mean of member indicators.
Suppose $X_i$ records whether member $i$ makes a particular error, all have variance
$\sigma^2$, and every pair has correlation $\rho$. Expanding the variance of their
mean gives
\begin{align}
 \operatorname{Var}\!\left(\frac1m\sum_{i=1}^m X_i\right)
 &=\frac{1}{m^2}\left(m\sigma^2+m(m-1)\rho\sigma^2\right)\notag\\
 &=\frac{\sigma^2}{m}\{1+(m-1)\rho\}.
 \label{eq:ensemble-variance}
\end{align}
The expression explains why counting correlated members exaggerates the amount
of independent information. It does not identify their correlation in practice,
and it does not supply a legal error probability. Its role is to motivate
measurement of shared failure and to discourage confidence claims based on panel
size alone.

Diversity should therefore be recorded along several axes. Members may use
different model families, prompts, representations, source acquisition routes,
retrieval methods, factual schemas or verification methods. These differences are
opportunities for complementary errors; they do not prove independence. Two
vendors can share training material. Two parsers can both omit an image. Two human
reviewers can both receive a leading case summary. The record should describe
what differs and what remains shared.

The strongest inexpensive separation is often procedural. Let source inventory
and initial interpretation happen before either sees the other's preferred
conclusion. Freeze the initial outputs. Only then expose disagreements for
investigation. This preserves evidence that one member independently noticed a
condition. If all members see the first answer before responding, later agreement
may merely reflect anchoring. Deliberation can still be useful, but it must not
rewrite the history of independent proposals.

\section{The ensemble members and their responsibilities}

The prototype uses roles rather than assuming that every role is a language
model. The source custodian acquires the official document, retains the exact
bytes and metadata, and records incorporated materials. A document inventory
process enumerates paragraphs, footnotes, tables, attachments and visible
cross-references. A second inventory method compares the enumeration with the
rendered document. Both methods may be automated initially, but unexplained
mismatches remain review questions.

The normative reader identifies who must or may do what, under which conditions,
with which exceptions and dates. A controlled-language reader expresses the same
material using a restricted template. Their outputs are compared at the level
of actor, activity, modality, object, condition, exception, time and dependency.
This comparison can reveal a missing phrase even when the two prose summaries
sound similar.

A counterinterpretation reader is required to propose consequential alternatives
where the source supports them, and to say when it found none. It is not required
to invent a controversy for every sentence. Its task includes seeking an omitted
qualification, an alternative attachment of an exception, or a classification
that changes a boundary case. Unsupported alternatives are preserved as such,
which helps distinguish deliberate testing from legal disagreement.

The formal checker validates types, references and explicit semantics, compares
candidate rules over supported domains, and replays witnesses. It has no
permission to promote a source assumption to a fact. The scenario designer
constructs examples that separate readings and includes unknown, conflicting and
out-of-scope evidence. Its expected outcomes come from an independent review
process or remain labelled provisional.

The implementation checker compares the accepted rule with the Java package and
verifies the package's exact identity. Finally, the dossier assembler reconciles
all required outputs and creates the review report. It may summarize, but it may
not suppress a minority objection, change its severity or mark it resolved without
a linked resolution record.

A role can fail in several ways: unavailable, malformed output, unsupported
claim, resource limit, incompatible source version or abstention. Abstention is
an informative response when a role lacks the evidence needed for its question.
It is not a vote for the majority. If a role is mandatory for the configured
coverage claim, its missing output prevents that claim from being marked complete.
A replacement role must be explicitly authorized by the configuration and its
independence limits recorded.

\begin{table}[htbp]
\centering\small
\caption{Responsibilities that create different opportunities to detect error.}
\label{tab:ensemble-roles}
\begin{tabularx}{\textwidth}{P{31mm}YY}
\toprule
Role & Question answered & Failure that must remain visible\\
\midrule
Source inventory & What material is present and incorporated? & Missing attachment, unmatched paragraph, inaccessible dependency\\
Normative reader & What obligations, conditions and exceptions are expressed? & Omitted modality, scope or qualification\\
Alternative reader & What supported choice changes a consequence? & No supported alternative found; unexplored family\\
Argument critic & What defeats the premises or inference? & Unanswered objection or unverified authority\\
Formal checker & What follows from this specified rule? & Invalid encoding, unknown solver result, limited domain\\
Reference reviewer & What is expected for this scenario, and why? & Disagreement, incomplete facts, shared authorship\\
Java checker & Does this package preserve the accepted specification? & Mismatch, unsupported operation, stale build\\
Dossier assembler & Are all required results and dispositions present? & Missing member, open issue, broken evidence link\\
\bottomrule
\end{tabularx}
\end{table}

\section{Build a source perimeter before interpreting it}

The source perimeter is the set of documents and versions the investigation
claims to cover. It begins with the named circular, but incorporated definitions,
annexes and transition provisions may expand it. Each dependency has a reason:
explicit incorporation, definition of a term, amendment, background explanation,
or a proposed analogy. These relationships must not all be collapsed into a
generic ``relevant document'' link.

For 23EC46, the retained inventory has eighteen paragraphs and four footnotes.
The executable slice covers paragraph 10. That is a legitimate prototype scope
only if the report makes the remaining material visible and explains whether it
is contextual, another obligation, a dependency, or outside the selected slice.
Labelling all unimplemented paragraphs ``not applicable'' merely because the
prototype lacks code would be wrong. Implementation capability is not a legal
applicability test.

The inventory assigns each source unit a stable locator within a particular
source hash. It records the extracted text, its position, and whether the unit
contains a possible normative statement or reference. The second inventory can
use a different parser or rendered-page inspection. A mismatch creates a source
coverage discrepancy before any interpretation ranking takes place.

The controller expands dependencies until each required edge reaches a retained
source, a justified out-of-scope disposition, or an unresolved dependency record.
Cycles are legal as document relationships: a code and a guidance note can refer
to each other. The traversal uses visited version identifiers to avoid infinite
retrieval, while preserving both edges. A retrieval depth limit stops processing
but leaves the unexpanded edge visible. It does not prove that the dependency is
unimportant.

Publication date, effective date and retrieval date occupy separate fields. A
circular published today may describe a later transition or an earlier
arrangement. The source custodian cannot infer commencement solely from the HTTP
response time or the first date in the text. The retained source packet is a
historical evidence object. A current applicability assessment requires a
separate search for amendments and supersession, with its own cut-off date.

This inventory is an important defense against common omissions, but it is not
an oracle. An extractor might enumerate a footnote without understanding its
legal effect. The requirement that each unit receives a reasoned disposition
connects inventory to interpretation. A coverage check should fail if a unit is
present but has no explanation, or if an explanation claims executable coverage
without identifying the rule and verification evidence that provide it.

\section{Freeze blind proposals before reconciliation}

An initial proposal request contains the frozen source packet, the selected
activity profile, the vocabulary already approved for that profile, and the
member's task. It excludes other members' proposed conclusions. The response is
stored with model or tool identity, configuration, prompt-template version, source
hashes and completion status. Raw output is retained subject to the bank's access
and retention policy; the structured candidate is a validated derivative.

A candidate cannot refer to a source version it was not given without opening a
new retrieval request. This rule prevents a model's recalled description of an
unseen amendment from entering as if it were inspected evidence. Recalled
material can nominate a search, but the supporting claim remains unverified
until the source is acquired and checked.

The blind phase ends when every required member has responded or reached a
terminal failure under the initial-phase budget. The controller records that
boundary. It then compares the candidates and inventories. Later repairs may
share all relevant evidence, because their purpose is informed resolution. The
record must continue to distinguish the initial independent proposals from
responses influenced by the shared discussion.

This separation prevents a misleading report such as ``five independent members
agreed'' when four were shown the first member's answer. It also makes a useful
experiment possible: measure which defects were caught before deliberation and
which required shared investigation. A system that produces agreement by
suppressing dissent should look different from one that resolves an objection
with a missing official definition.

Candidate revisions are immutable. A revised rule points to its parent, states
what changed, and links the evidence or diagnostic that motivated the change.
The original remains available. Reviewer approvals attach to a particular
candidate hash, not to a mutable display name such as ``latest rule.'' If a
source, assumption or rule changes, the old approval cannot silently authorize
the new candidate.

\section{Detect discrepancies at the right level}

A discrepancy is a structured difference that may affect coverage, meaning,
execution or evidence. It need not be a pair of contradictory prose answers.
One member may omit an input that another requires; two members may agree on a
rule but cite different versions; one parser may miss a footnote; the Java
runtime may return the same decision with an incomplete audit trace. All deserve
a recorded disposition appropriate to their consequence.

The comparator first checks identity: source versions, activity profile, subject
unit, candidate dependencies and evaluation times. Comparing answers to different
questions creates false disagreement. A model interpreting a whole offer and one
interpreting a single benefit must be identified as differing in subject identity
before their truth values are compared.

It then compares structured meaning: actor, action, object, modality, conditions,
exceptions, quantifiers, temporal boundaries and references. Exact structural
matches are cheap evidence of agreement about those fields, but source coverage
still needs a separate check. Differences in presentation can be classified as
nonmaterial only with an explicit rationale, such as identical canonical syntax
or a domain-qualified equivalence result.

Formal outcome comparison comes next. A counterexample is especially useful when
two apparently similar rules disagree only at a boundary. For paragraph 10,
omitting product-type linkage is exposed by $P=\False$, $T=\True$, with the other
conditions established. Omitting the discount exception is exposed by a genuine
fee discount. Neither example resolves whether a particular cashback is a fee
discount; that is a different discrepancy with a different evidence need.

The comparator also checks omissions against independent inventories. Every
required source unit, dependency and ensemble role must have a disposition. An
all-agreeing set of candidates can therefore still produce discrepancies. This
is the intended behavior: agreement among candidates cannot erase an uncovered
paragraph or a failed source acquisition.

A discrepancy record includes its origin, affected source units, candidate
identifiers, issue family, factual or legal question, observed difference,
materiality, evidence needed, permitted next actions and current status. It also
stores the consequence of leaving it unresolved. A classification issue might
block a particular incentive category; a corrupted source hash might invalidate
the entire run. The scope of the block must be explicit and conservative when
impact cannot yet be determined.

\section{An issue needs an identity that survives retries}

Naive retry systems often attach the attempt counter to a model response. When
the response creates a revised candidate, a new counter starts. A difficult issue
can then consume unlimited attempts by repeatedly changing its wording or
creating descendants. The proposed controller attaches budgets to the run and to
an issue root that survives candidate revisions.

An issue root identifies the source context, the substantive question and the
affected subject or rule family. It does not include a candidate hash as the only
identity, because candidate hashes necessarily change during repair. The system
assigns an immutable root identifier when the question is first recorded.
Subsequent discrepancies can link to that root or create a child issue. A child
inherits the root's budget account and the run's limits. Creating a new display
record never creates new spending authority.

Automatic deduplication proposes that two issues have the same root when they
concern the same source unit, semantic dimension and consequence. A reviewer or
strict deterministic criterion must resolve uncertain merges. An erroneous merge
could hide a second problem; an erroneous split could waste work. Global action
limits protect against unlimited spending even when deduplication is imperfect,
while the report shows possible duplicates for later review.

Issue identity also separates changed evidence from repeated reasoning. If an
unavailable FAQ becomes available during the same run, the controller can attach
the new version and continue within the remaining budget. If a later circular
changes the applicable rule, that is a new source context and normally a new run.
The old run remains closed with its original result. Opening the new run requires
a recorded trigger and authorization; a worker cannot evade exhausted limits by
automatically renaming the investigation.

A candidate that depends on several issues keeps all of their identifiers. A
resolution of one issue cannot clear the others. The dependency relation also
supports invalidation: retracting a definition affects every candidate that used
it, even if those candidates were reached through different search branches.
The service must query this relation directly rather than trying to reconstruct
it from narrative summaries.

\section{Statuses describe evidence, not optimism}

The issue lifecycle begins at \texttt{OPEN}. An issued action moves it to
\texttt{INVESTIGATING}. It may then return to \texttt{OPEN} with additional evidence,
move to \texttt{AWAITING\_EVIDENCE}, or require \texttt{HUMAN\_REQUIRED}. These are
unresolved states. They differ in the next useful action, not in whether the
system may assume the preferred answer.

Resolution has three main forms. \texttt{RESOLVED\_EVIDENCE} means that an inspected
source or factual record establishes the missing premise under the accepted
review rules. \texttt{RESOLVED\_ADJUDICATION} means that an authorized reviewer has
made and explained the interpretive decision. \texttt{RESOLVED\_EQUIVALENCE} means
that the apparently different executions have been shown equivalent over a
specified domain and that no distinct source commitment remains at issue.
These statuses require different evidence and cannot substitute for one another.

For example, an SMT equivalence result cannot resolve a dispute about whether the
word ``may'' establishes a universal classification. Both candidates may have
been encoded with the same unjustified classification flag. Conversely, an
explicit source passage that defeats a missing-exception candidate need not wait
for a probabilistic vote. The source-based resolution identifies the clause,
updates the candidate and requires the affected formal and Java checks again.

An issue may become terminal while remaining unresolved. The terminal reason can
be \texttt{BUDGET\_EXHAUSTED}, \texttt{DEADLINE\_REACHED}, \texttt{NO\_PROGRESS},
\texttt{DEPENDENCY\_UNAVAILABLE}, \texttt{CANCELLED}, or an integrity failure.
Terminality means automatic work has stopped. It says nothing favorable about
the candidate. The report must therefore store resolution status separately from
processing status; a single Boolean \texttt{done} is inadequate.

Materiality is another separate field. An issue is material when its unresolved
answer could change applicability, a required action, the decision, the timing,
the explanation required for review, or the validity of the evidence chain. If
impact is unknown, it is treated as material until assessed. Downgrading an issue
requires a reason and authority appropriate to the configured review policy.
A generator cannot label its own objection immaterial simply to unblock release.

A resolved issue can be reopened by new evidence, a source amendment, an invalid
witness or a discovered dependency. Reopening creates a new event and invalidates
results that depended on the old resolution. It does not erase the earlier
record. If the original run is already closed, the successor investigation links
to it and states which conclusion is being reconsidered.

\section{Choose a remedy that can answer the question}

Each discrepancy type has a permitted action set. A missing source calls for
retrieval or escalation to the source custodian. A parsing disagreement calls for
inspection of the original page, an alternative parser or manual transcription
with provenance. A scope disagreement calls for analysis of the actor and
activity language, incorporated definitions and the selected bank profile. A
formal mismatch calls for a replayable witness and a correction to the relevant
encoding or runtime component.

A useful action has a question, an expected evidence type, a bounded execution
method and a success condition. ``Ask another model'' is incomplete. ``Ask the
alternative reader to identify source-supported reasons why the exception might
attach to the whole offer, using only the retained paragraph and definition'' is
a bounded interpretive task. Its result is still a proposal, but the task makes
clear what new information it seeks.

The controller should prefer actions with a plausible path to resolution. If the
question is whether a fee was paid, querying transaction evidence may be more
useful than sampling another legal explanation. If the question concerns the
meaning of an incorporated term, more transaction data may be irrelevant. This
matching of remedy to issue is where structured uncertainty becomes operationally
valuable.

An action can also deliberately seek a falsifier. For a proposed rule, request a
feasible scenario in which it contradicts an express source condition. For a
classification, request a fact pattern that the definition treats counterintuitively
and ask which source justifies that boundary. For a claimed equivalence, ask the
solver for a difference rather than asking a model whether the programs look
similar. Falsification attempts must retain their domain and assumptions.

Repeated actions are allowed only when their rationale changes or a bounded retry
policy permits recovery from a technical failure. Retrying an unavailable server
can make sense; asking the identical question with no new evidence is unlikely
to resolve a legal ambiguity. The report should show which retries addressed
transport failure and which explored a substantive alternative. Both consume
resources, but they mean different things for the investigation.

A repair is validated against the original discrepancy and against the unaffected
requirements. Fixing the type-of-product omission must not remove the fee-discount
exception. A local witness becoming consistent is necessary evidence about that
repair, but insufficient evidence about the whole rule. The controller runs the
dependent validation set, records any new discrepancies, and keeps the old issue
open until its resolution criterion is actually satisfied.

\section{A resolution matrix prevents the wrong kind of closure}

The controller needs a typed relationship between a discrepancy and the evidence
that can close it. Without that relationship, every issue can eventually be
``resolved'' by another fluent explanation. The matrix below is a proposed
engineering contract for the first implementation. It states the minimum kind of
evidence, not a substitute for legal judgment about whether that evidence is
persuasive in the particular case.

\begin{longtable}{@{}P{27mm}P{56mm}P{65mm}@{}}
\caption{Discrepancy-specific actions and closure evidence.}\label{tab:resolution-matrix}\\
\toprule
Discrepancy & Appropriate next action & What can justify closure\\\midrule\endfirsthead
\toprule Discrepancy & Appropriate next action & What can justify closure\\\midrule\endhead
Missing source unit & Compare original rendering and independent inventory & Retained unit, corrected inventory and reviewed disposition; agreement among summaries is insufficient.\\
Unavailable dependency & Retrieve the named version or seek authorized clarification & Inspected applicable source, or an explicit adjudication of the remaining question; a guessed definition is insufficient.\\
Different source versions & Establish the assessment cut-off and amendment relationship & A documented version selection and reconsideration of affected candidates.\\
Actor or activity scope & Inspect addressee, operative language and profile facts & Reviewed applicability proposition with evidence for the actual bank activity.\\
Disputed classification & Identify definition and distinguishing factual evidence & An accepted definition and sufficient facts, or an attributable adjudication; a Boolean label alone is insufficient.\\
Exception attachment & Construct component and whole-arrangement examples & Source-grounded structural choice, with contrary reading disposition and regression cases.\\
Necessary versus sufficient condition & Compare implication directions on boundary scenarios & A reviewed direction of dependence and a rule that preserves other required conditions.\\
Timing or transition & Separate publication, commencement, valid time and knowledge time & Reviewed interval and transition rule with endpoint examples.\\
Factual conflict & Obtain correction or apply an approved evidence-priority rule & Attributable supersession or reviewed normalization; newest arrival alone is insufficient.\\
Formal mismatch & Produce and replay a feasible difference witness & Validated repair or domain-qualified equivalence, with source assumptions retained.\\
Unsupported operation & Inspect whether a supported encoding preserves meaning & Reviewed encoding plus conformance evidence, or an approved runtime extension; compiler rejection does not defeat the reading.\\
Java mismatch & Compare full semantic results and trace & Corrected exact package with repeated affected verification and new build identity.\\
Missing ensemble member & Retry under budget or use an authorized replacement & Required role output with documented independence limits; silence is not agreement.\\
Search truncation & Investigate retained frontier or narrow the claimed scope & Completed declared finite search or explicit reviewed restriction; a high best score is insufficient.\\
Broken approval lineage & Reconstruct exact source, report, bundle and build links & Fresh valid approvals for exact identities; a matching display title is insufficient.\\
\bottomrule
\end{longtable}

Some issues need more than one row. A cashback classification can depend on an
unavailable definition and an uncertain factual relationship to a fee. Retrieving
the definition resolves the dependency but may leave the factual question open.
The controller should represent these as linked issues so that progress is visible
without overstating closure. A single combined issue can be acceptable if its
resolution criterion explicitly requires both conditions.

The matrix also guides action rejection. If a proposed action cannot produce the
needed evidence, the scheduler should normally decline it and explain why. A new
model paraphrase is not an appropriate remedy for a corrupted source digest. A
solver equivalence query is not an appropriate remedy for a missing commencement
date. This does not prohibit exploratory work; it requires that exploratory work
be labelled as such and charged to the investigation budget rather than counted
as resolution.

Closure is evaluated against the issue's original question and current source
context. Suppose a model responds to a component-level question with a correct
whole-offer explanation. The response may be useful, but it has not answered the
original question. The controller creates or links the relevant candidate and
keeps the original issue open. This guards against a subtle form of apparent
progress in which the task changes until the model can answer it confidently.

\section{Resolution can reveal that the first question was wrong}

A structured controller must still permit substantive reframing. The initial
question may assume a distinction that the source does not make. For example,
the team may ask whether a cashback is a discount when the relevant provision
instead classifies the entire arrangement through another defined concept.
The correct response is to record the changed understanding, identify which
source supports it and replace the candidate with a new revision. The old question
is dispositioned as superseded by a reasoned reframing, not silently deleted.

Reframing does not reset budgets. The successor issue links to the original root
or consumes the same global account, and the report explains the relationship.
If the reframing expands the source perimeter beyond the authorized run, the
current investigation can end with a request for a successor run. This prevents a
useful ability to change direction from becoming an unbounded search loophole.

A concrete example concerns subject identity. A provider initially treats an
offer as one benefit. The independent inventory identifies a source passage that
requires separate treatment of components. The controller creates a component
model, revises factual requirements and reruns affected scenarios. It also records
that prior results used a different subject abstraction. Merely changing the
formula while leaving the whole-offer adapter untouched would not complete the
repair.

Another example concerns modality. A candidate initially turns guidance about
factors to consider into an unconditional prohibition. Investigation shows that
the source requires an assessment rather than a predetermined outcome. The repair
may need a workflow obligation, evidence of consideration and a review record,
not another Boolean threshold. If the current runtime cannot express the accepted
requirement, the result is a supported interpretation with an implementation gap.
That is more accurate than forcing the source into a familiar Boolean shape.

These cases explain why the issue register includes the affected decision and the
kind of evidence required. The controller can recognize that the target has
changed and invalidate dependent checks. A test suite for the previous target
cannot automatically certify the new one, even if many examples happen to retain
the same final status.

\section{The investigation graph should preserve shared causes}

The following diagram separates the objects that are often collapsed into a
single chain of model messages. Several candidates can depend on one definition,
and several discrepancies can affect one candidate. Actions investigate issues;
results create evidence or revisions. The final report reconciles the graph rather
than selecting the last message as the answer.

\begin{figure}[htbp]
\centering
\begin{tikzpicture}[
 box/.style={draw=teal,rounded corners,align=center,text width=43mm,minimum height=12mm,font=\small},
 arr/.style={-{Latex[length=2mm]},thick,draw=navy},node distance=12mm and 10mm]
\node[box] (source) {Retained source and\\reviewed inventory};
\node[box,right=of source] (definition) {Shared definition\\or missing premise};
\node[box,below=of source] (candidates) {Immutable competing\\candidate revisions};
\node[box,right=of candidates] (issues) {Discrepancies and\\issue-root budgets};
\node[box,below=of candidates] (evidence) {Validated evidence,\\witnesses and checks};
\node[box,right=of evidence] (actions) {Reserved bounded\\resolution actions};
\node[box,below=of evidence,xshift=26mm,text width=72mm] (report) {Complete report: resolved reasons,\\surviving alternatives and uncertainty};
\draw[arr] (source)--(candidates);
\draw[arr] (source)--(definition);
\draw[arr] (definition)--(issues);
\draw[arr] (definition)--(candidates);
\draw[arr] (candidates)--(issues);
\draw[arr] (issues)--(actions);
\draw[arr] (actions)--(evidence);
\draw[arr] (evidence)--(candidates);
\draw[arr] (evidence)--(report);
\draw[arr] (issues.east)--++(8mm,0)|-(report.east);
\end{tikzpicture}
\caption{Search creates a history, while dependencies explain what must be reconsidered.}
\label{fig:ensemble-dependencies}
\end{figure}

An important consequence is that one new source can resolve several issues
without spending a separate retrieval action for each. The action is charged once,
and each affected issue receives a validated link to the result. Their root
budgets do not acquire extra credits. Conversely, one failed source acquisition
can block several candidates, and the report should identify that shared cause
rather than presenting a long list of apparently unrelated uncertainties.

The graph also prevents misleading independence claims. If four candidate members
all rely on the same extracted definition and factual mapping, the report can
show that common dependency. Their reasoning paths may differ, but their evidence
base is shared. This is a more useful account of redundancy than simply displaying
four model logos beside an agreement percentage.

\section{Finite budgets are part of the semantics}

The user requires automatic processing up to a maximum number of times. To make
that requirement implementable, we distinguish an initial proposal phase from
resolution rounds. A resolution round consists of selecting currently open
issues, reserving authorized actions, collecting bounded results and reconciling
the resulting candidates. The round count increases when the controller commits
the round's action plan, even if every external action later times out.

The configuration supplies a maximum number of resolution rounds, a maximum
number of issued actions per issue root, a maximum number of issued actions for
the run, and an absolute deadline. It also supplies per-action timeouts, candidate
and graph-size limits, and any token or monetary limits. None is inferred from
the fact that the existing drafting stub happens to use three attempts. That
stub's retry behavior addresses a different contract.

A reference policy for deterministic fixtures uses three resolution rounds,
three issued actions per issue root, and twelve external actions in total,
including the initial proposal actions. It permits at most four initial external
proposal actions, leaving at most eight for later investigation. These are
engineering examples chosen to make boundary tests tractable. They are not
empirically justified production defaults or recommendations about how much
legal review is enough. Production configuration must be approved and evaluated
for the selected task and cost constraints.

The reference policy also sets a ten-minute run deadline and a thirty-second
external action timeout. Deterministic local checks have separate finite input
and execution bounds. A mandatory check is never silently skipped when its bound
is reached: it returns an incomplete result and blocks the corresponding claim.
Token and monetary caps are nullable only when the provider integration cannot
supply a reliable bound; such a configuration is labelled unsuitable for an
unattended paid run. A paid deployment must supply conservative reservations.

Let $B$ be the configured global action cap and $a$ the number of actions already
issued. Before dispatching any new external action, the controller must establish
\begin{equation}
 a+1\leq B,\qquad a_r+1\leq B_r,\qquad t<t_{\rm deadline},
 \label{eq:budget-guard}
\end{equation}
where $a_r$ and $B_r$ are the issued count and cap for the affected issue root.
For an initial proposal action, a separate initial-phase account replaces the
issue-root account. The guard also checks the round, token, cost and action-type
constraints. Passing one limit does not override another.

Reservation occurs atomically before dispatch. The issued count increases even if
the worker crashes before learning whether the provider accepted the request.
This conservative accounting prevents an ambiguous network result from becoming
free work. A reservation can be cancelled without spending only before dispatch
is durably committed and only through an atomic state transition that proves no
worker can still send the action.

An action's effective waiting limit is the smaller of its configured timeout
and the time remaining before the run deadline. The worker uses a monotonic clock
for elapsed waiting, while the durable UTC deadline supports restart and audit.
Clock adjustments cannot replenish an action budget. At the run deadline, pending
actions receive terminal or remote-result-unknown dispositions; the controller
does not extend the run simply because an action began a moment earlier.

Provider SDK retries must be disabled or included in the same accounting. A
single logical action that secretly makes five billable requests defeats the
meaning of the cap. When a provider offers idempotency keys, the adapter uses the
recorded action identifier. When it does not, the system cannot promise exactly
once execution across a lost response; it reports the uncertainty, counts the
attempt, and follows the configured retry policy.

\section{No progress is a substantive stopping reason}

A finite global budget guarantees that external actions cannot continue forever,
but it need not be spent in full. A no-progress condition avoids cycling through
identical proposals. The controller computes a progress signature from accepted
source additions, newly supported assumptions, new feasible witnesses, issue
resolutions and genuinely distinct candidate semantics. Cosmetic wording changes
do not count as progress.

The signature is a deterministic summary of validated changes, not a model's
claim that its answer improved. If a resolution round produces no such change,
the controller records that fact. The reference policy stops after two successive
rounds without progress, unless an already reserved action is still awaiting a
result within the deadline. This is an engineering stopping rule for the fixture
profile. A production policy may choose another threshold, but must disclose it
and evaluate the cost of premature stopping.

New unsupported candidates are not automatically progress. Otherwise a generator
could prevent termination by inventing fresh interpretations indefinitely.
Similarly, splitting one issue into several differently worded issues does not
reset the no-progress count or any budget. The controller retains those proposals
for review when useful, but their existence alone does not establish a reason to
continue automatic investigation.

A source becoming temporarily unavailable may justify a scheduled successor run,
but not an endless loop in the current run. The report records the failed
retrieval, the affected question and the conditions under which a new run would
be useful. A human or an authorized source-change service can initiate that run
with a new recorded budget. This preserves the distinction between bounded
automation and an operational process that may legitimately continue over days.

No-progress termination is not a judgment that the research direction has failed.
It says that the current actions, sources and budget have not resolved the
question. A later phase designed to obtain independent authority may still be
justified. The system must report whether it encountered an invalid source,
broken implementation, insufficient evidence, or simply a candidate that failed
its checks. Those outcomes call for different next steps.

\section{A deterministic first scheduler}

The first implementation needs a concrete action order before a learned search
policy is available. The reference scheduler uses a lexicographic key rather
than an opaque weighted score. It first handles integrity and source-coverage
problems, then missing dependencies and applicability, then definitions and
exception structure, then formal and Java mismatches, and finally report or
presentation discrepancies. This order is a proposed engineering policy: earlier
questions can invalidate the inputs to later checks. It is not a legal ranking
of every possible issue.

Within a class, unresolved material or unknown-materiality issues precede reviewed
immaterial issues. Earlier creation sequence precedes later sequence, and the
stable issue identifier breaks any remaining tie. Every selected action must be
permitted for the issue kind and have remaining root and global budget. If an
earlier issue is awaiting evidence that the available tools cannot obtain, it
remains visible while the scheduler can make progress on another issue. Priority
is not permission to spend the whole budget repeating an unavailable action.

For each issue, action selection follows a short decision procedure. Check whether
an already retained source or result answers the question; if so, validate and
link it without issuing a new external request. Otherwise, retrieve a named
missing authoritative dependency when a permitted retrieval route exists. If the
source is present but the candidate encoding differs, request a witness or a
source-grounded repair. If competing supported readings remain, select a
contrasting scenario or focused argument question. If no permitted action can
supply the required evidence, prepare the human review question and stop automatic
work on that issue.

Actions are not dispatched merely because the scheduler can name them. A
pre-dispatch validator checks that their inputs differ meaningfully from a prior
failed substantive action, unless the retry is an explicitly permitted transport
recovery. It also checks source identity, adapter support and the remaining
budget. Rejection is recorded as a scheduling disposition; it does not consume
an external-action slot unless a dispatch reservation has already been committed.
Local scheduling and validation remain bounded by graph and execution limits.

Candidate breadth uses a similarly explicit rule. At initial reconciliation,
reserve one retained slot for each supported interpretation family before
allocating additional slots to refinements. If the number of families exceeds the
candidate cap, retain the families in the same declared issue-priority and
creation order, store the excluded family identifiers in the frontier, and open
a search-incomplete issue. This policy favors reproducibility and visible
incompleteness. It does not claim that its ordering finds the best legal theory.

For the optional pair-disagreement heuristic, the first prototype uses equal
weights in equation~\eqref{eq:disagreement-question}. It enumerates the available
feasible contrasting scenarios within the supported comparison fragment, selects
the scenario separating the most retained pairs, and breaks ties by stable
scenario identifier. A cost-aware or learned weight can later replace this rule
under a versioned evaluation. Equal weights here are a baseline convenience,
not equal probabilities that every reading is correct.

The implementation report records the scheduler version and each selection key.
An engineer can therefore reconstruct why an issue was investigated before
another, and a reviewer can challenge the priority policy without guessing at a
model's hidden preference. Chapter~\ref{ch:evaluation} compares this deterministic
baseline with enhanced search at equal resources. Both must obey the same budget,
reporting and release constraints.

\section{A controller that can be implemented and checked}

The controller is a deterministic state machine around fallible external tools.
Its durable state contains the run configuration, source snapshot, required role
results, candidate revisions, issue roots, action records, budgets and report
status. External results enter through validation. They do not directly mutate
approvals or release state.

The following pseudocode specifies the main control flow. The functions named
\texttt{commit} and \texttt{reserve} are transactions; they must either write all
of their stated changes or none. The action worker uses the committed action
record and cannot create additional actions on its own.

\begin{lstlisting}[caption={Main controller, with explicit terminal uncertainty.}]
start_run(source_packet, profile, required_policy):
    validate_source_identity_and_policy()
    commit(run = INITIALIZING, counters = 0)
    reserve_and_issue_initial_roles_with_global_budget()
    collect_until_complete_or_action_deadline()
    validate_outputs_and_freeze_initial_proposals()
    reconcile_inventory_candidates_and_required_roles()

    while has_unresolved_issue():
        if integrity_failure():
            close_as(FAILED_INTEGRITY); break
        if cancelled():
            close_as(CANCELLED); break
        if deadline_or_global_or_round_limit_reached():
            mark_remaining_issues_with_stop_reason()
            close_as(BLOCKED_UNRESOLVED); break
        if no_progress_limit_reached():
            mark_remaining_issues(NO_PROGRESS)
            close_as(BLOCKED_UNRESOLVED); break

        plan = choose_permitted_discriminating_actions()
        if plan.is_empty():
            mark_remaining_issues(HUMAN_REQUIRED)
            close_as(BLOCKED_UNRESOLVED); break
        commit_next_round_and_reserve_bounded_actions(plan)
        dispatch_only_reserved_actions()
        collect_results_until_each_action_is_terminal()
        validate_results_and_create_candidate_revisions()
        rerun_affected_checks_and_reconcile_all_open_issues()
        record_progress_signature()

    if not already_closed():
        apply_integrity_cancellation_and_deadline_guards()
    if not already_closed():
        check_required_roles_coverage_and_all_mandatory_checks()
        if any_requirement_incomplete():
            close_as(BLOCKED_UNRESOLVED)
        else:
            close_as(READY_FOR_REVIEW)
    write_complete_report_for_every_terminal_status()
    return report_reference
\end{lstlisting}

Every detected discrepancy enters reconciliation, including differences later
judged immaterial. A deterministic presentation-equivalence check can close a
purely cosmetic issue immediately with its reason; unassessed differences remain
unresolved and conservatively material. The loop therefore processes unresolved
issues rather than only issues already assigned a high severity. Integrity,
cancellation and deadline guards are checked again before final readiness, even
when no issue remains, so an empty loop cannot bypass a terminal condition.

The final coverage check is essential. If no comparator detects a disagreement,
the loop might never run. The controller must still require complete source
coverage, mandatory members and verification results. ``No open issue'' and
``all required evidence present'' are separate conditions. A failed member must
also create an issue, but the final check protects against a defect in that issue
creation path.

Actions that finish after the run deadline are retained as late observations.
They do not reopen a closed run or retroactively change its result. A successor
run may use them after validation. This rule makes reports stable and prevents a
late worker from changing the package that a reviewer has already inspected.
Late evidence can still trigger a documented follow-up when it changes risk.

The controller must produce a report even when it fails. If the primary report
writer is unavailable, the durable run and action records support recovery of a
minimal failure report. An integrity failure prevents a positive result, but it
must not erase the explanation of what was attempted. Report generation itself
has bounded resource use and a recovery path; an exception thrown while rendering
a large argument graph cannot make the investigation disappear.

\section{Termination and the boundary of the proof}

We can prove a limited but useful property of the proposed controller. Assume a
finite input packet, a finite configured action cap $B$, atomic issue accounting,
finite local reconciliation with explicit graph limits, and an enforced timeout
for every dispatched action. Assume also that the storage and supervisor either
respond within their operational bounds or move the run to a recoverable failure
state. These assumptions concern implementation and infrastructure; the proof
does not establish them merely by listing them.

Define the remaining action budget as $V=B-a$, where $a$ is the durable number of
issued actions. Initially $V$ is a nonnegative integer. Each new dispatch requires
an atomic reservation and reduces $V$ by one. No transition increases $V$ in the
same run. Therefore there can be at most $B$ dispatches. Between dispatches,
reconciliation is finite by its input and execution bounds. Waiting is finite by
the action timeout and absolute deadline. When no permitted action remains, the
controller enters a terminal state. Under these assumptions, automatic processing
terminates.

This argument remains valid when an action discovers additional issues. New
issues do not replenish the global budget. It also remains valid when a worker
restarts, provided the durable counter is authoritative and a lease does not
authorize a second unrecorded dispatch. The per-issue cap and round cap are useful
additional constraints, but the global cap is what prevents an unbounded stream
of newly named issues from evading termination.

The proof is about bounded processing, not resolution. A controller can terminate
with every difficult issue unresolved. It is also not a proof of wall-clock
availability under arbitrary storage failure. If the process is indefinitely
suspended or the database never responds, no local algorithm can guarantee a
report at a fixed real-world time without an external supervisor. The operational
design must therefore monitor leases and deadlines independently and recover a
failure report when infrastructure returns.

A second invariant concerns release. Let $M(r)$ mean that run $r$ has an unresolved
material issue, $C(r)$ mean that its mandatory coverage and checks are complete,
and $A(r,b)$ mean that the required authorized reviewers approved exact bundle
$b$ on the evidence from $r$. The release service requires
\begin{equation}
 \operatorname{Release}(r,b)\Longrightarrow
 \neg M(r)\land C(r)\land A(r,b)\land\operatorname{Integrity}(r,b).
 \label{eq:ensemble-release}
\end{equation}
This is a design invariant to be enforced by the release transaction. The
controller's \texttt{READY\_FOR\_REVIEW} state supplies neither $A$ nor release
itself. An implementation test must attempt to bypass each conjunct, including
through a stale approval or a source update, and verify rejection.

The invariant also exposes the limit of the design. If the inventory misses a
material source and no check detects it, $C(r)$ could be incorrectly recorded.
Formal verification of the release transaction would still prove only that it
uses the recorded evidence correctly. Independent source review and empirical
omission tests address the reliability of that evidence. The book's separation
of proof claims is therefore retained inside the ensemble rather than abandoned
at its most important decision.

\section{Durable actions, leases and ambiguous execution}

An action moves through \texttt{RESERVED}, \texttt{DISPATCHED}, and one of
\texttt{SUCCEEDED}, \texttt{FAILED}, \texttt{TIMED\_OUT}, or
\texttt{RESULT\_UNKNOWN}. The last state is necessary when the controller cannot
determine whether a remote service performed the action. A lost connection after
submission is different from a request rejected before dispatch. Both can lead
to missing results, but only the former may already have consumed remote work.

The reservation transaction records the action identifier, input hashes, issue
root, round, budget charge, adapter identity, timeout and idempotency key. An
outbox entry is written in the same transaction. A worker claims that entry with
a lease and a fencing token. The token increases when ownership changes, so an
old worker cannot commit a result after a newer owner has taken responsibility.
A valid late result can be retained as an observation without changing the closed
action's authoritative state.

The dispatch boundary requires care. If the outbox records only ``not yet sent''
and a worker crashes after sending but before updating the database, another
worker may send again. A provider idempotency key can make that retry refer to the
same operation, but not every provider supports this guarantee. The adapter must
state its capability. Without remote idempotency, a recovery policy can choose to
leave the result unknown rather than risk another paid request, or authorize a
new counted action. It cannot truthfully promise exactly once remote execution.

The budget system reserves an upper bound on tokens or cost when those resources
are capped. Settlement records actual usage when known. A refund of unused
monetary reservation may be legitimate after a confirmed terminal response, but
it does not decrement the issued-action count. Otherwise a sequence of failed
requests could evade the maximum number of attempts. If actual usage exceeds a
provider's documented bound, the run records an accounting discrepancy and stops
new spending until the integration is investigated.

A compare-and-swap version protects issue updates. The worker returns evidence
against the issue revision it received. If the issue has changed, the controller
revalidates whether the evidence is still relevant before applying a resolution.
This prevents two simultaneous actions from overwriting each other's objections
or resolving an issue using a candidate that is no longer current.

Cancellation stops new reservations and attempts to cancel outstanding work.
Remote cancellation may be best effort. The report records outstanding or late
actions and keeps their budget charges. Cancellation cannot convert an incomplete
run into a clean result, and it cannot delete the source of a later adverse
finding. These details matter because regulatory evidence must survive ordinary
operational events, including a person closing a browser or restarting a worker.

\section{What counts as a justified repair}

A repair has two responsibilities: remove the identified defect and preserve the
requirements that were already correct. For a missing product-type condition,
the revised body must include the type route, and the previously verified gift,
exception and scope behavior must remain intact. The repair record identifies the
changed expression and the affected checks. It cannot simply report that a new
model answer received a higher score.

The original discrepancy remains immutable. The repair links to it and supplies
resolution evidence: the source clause, changed candidate, witness before and
after, regression checks and any reviewer decision. A failed repair becomes
another candidate revision with a failed status. This is useful evidence about
the investigation and should not be removed from the dossier merely because a
later revision passes.

A repair can be rejected at several layers. The source may not support its new
condition. The controlled-language rendering may no longer match the formal
rule. The formal rule may use an unsupported operation. The Java package may
miscompute the accepted expression. These are different failures. A compiler
rejection means the current encoding cannot be delivered through that compiler;
it does not necessarily refute the legal reading. The next action might be a
supported encoding or a documented implementation extension rather than a new
interpretation.

Reference cases require special protection. When a candidate disagrees with a
frozen expected outcome, the system records the disagreement and investigates
both sides. It does not assume the reference is infallible, particularly in our
same-author 22-case exercise. But changing the reference requires an independent
resolution record with evidence and reviewer authority. A model cannot rewrite
the expected result and then claim that its unchanged candidate now passes.

The same rule applies to formal domains. If a witness exposes a defect, narrowing
the domain to exclude it is a substantive change. The domain change needs a
source or operational justification, not merely a desire to eliminate the
counterexample. Every equivalence claim includes the domain hash and assumptions,
so an excluded scenario remains visible to reviewers.

Finally, repair must preserve dissent. If one supported reading is selected after
review, the rejected reading, its strongest argument and the reason for rejection
remain attached. Future evidence may make that earlier objection relevant. A
system that stores only the winning candidate forces the next team to rediscover
the same ambiguity and cannot explain why the original team chose as it did.

\section{Scores that assist investigation without pretending to be probabilities}

A single confidence number combines incompatible questions. A source can be
authoritative while its application is ambiguous. A candidate can be internally
consistent while missing a clause. A Java artifact can be verified against an
unapproved rule. The proposed output therefore uses a score vector and explicit
statuses rather than a scalar legal-confidence percentage.

The vector records source coverage, dependency completeness, supported and
unsupported assumptions, unresolved material objections, scenario agreement,
formal-check status, implementation-preservation status and search coverage.
Some components are counts or categorical states. Others may be empirical
scores, such as retrieval relevance. Their definitions are stored with the run.
A score is not allowed to compensate for a failed mandatory condition.

For example, a candidate might cover eighteen of eighteen inventoried paragraphs
with dispositions, yet have two unresolved incorporated sources and one material
classification question. Displaying only ``100 percent coverage'' would mislead.
The report must state that source-unit disposition is complete while dependency
and interpretation resolution are incomplete. Conversely, a lower model ranking
need not disqualify a candidate whose source argument remains undefeated.

Search priority can be a weighted function of expected information gain, action
cost and issue materiality. The weights are scheduling choices. They may be tuned
under a declared evaluation protocol, but they are not probabilities of legal
correctness. Every displayed numerical score includes its purpose and calibration
status. A field named \texttt{probability\_correct} remains absent or null until an
appropriate target-specific calibration study supports a precisely defined event.

A useful shortlist can be a Pareto set. Candidate $h$ dominates candidate $j$ only
when it is no worse on every declared comparable dimension and better on at least
one. Even then, dominance is an aid to presentation, not legal defeat. A candidate
with a unique source argument should not disappear merely because another has
more completed engineering checks. Unsupported candidates can be rejected for
specific reasons; incomparable supported candidates remain visible.

Set-valued reporting is especially valuable when authority is genuinely
incomplete. The dossier may present two conditional proposals: under definition
$D_1$, use rule $R_1$; under definition $D_2$, use rule $R_2$. It must identify who
can resolve the definition and what operations remain blocked. This is more
informative than selecting one proposal with a decorative confidence interval
whose coverage has never been established.

\section{A complete dry run: repairing an explicit omission}

The following trace is a designed example for the future controller. It has not
been executed by a live model and does not replace the retained second-circular
results. The source packet contains the selected paragraph, its surrounding
context and the unresolved dependency records. The activity profile and factual
schema are fixed before proposals begin.

Four initial actions are reserved. The normative reader proposes
$G\land P\land\neg D$, omitting product-type linkage. The controlled-language
reader proposes $G\land(P\lor T)\land\neg D$. The alternative reader agrees that
both product and type links require examination and identifies a separate
classification ambiguity. The inventory member confirms the paragraph and
footnote enumeration but reports an unavailable incorporated source. The run has
issued four actions; eight remain under the reference global cap.

Reconciliation creates issue $I_1$ for the missing type route, $I_2$ for the
classification ambiguity and $I_3$ for the unavailable dependency. The comparator
does not average the two formulas. It constructs the known-fact scenario
$G=\True$, $P=\False$, $T=\True$, $D=\False$, with scope established. The first
candidate returns false and the second true. The source passage explicitly
includes the type route, so $I_1$ has a concrete textual and behavioral basis.

In resolution round one, the controller reserves a source-grounded repair action
for $I_1$ and a dependency retrieval action for $I_3$. The issued count becomes
six before either is dispatched. The repair produces a child candidate that adds
$T$ to the disjunction. Its controlled-language rendering, type checks and witness
replay pass. The regression set checks that fee discounts remain excluded and
scope uncertainty retains its established behavior. The issue is resolved by
evidence with the original discrepancy and both candidate versions retained.

The dependency retrieval times out. Its action is terminal and counted; $I_3$
remains unresolved. The run does not claim complete circular interpretation just
because $I_1$ was repaired. If the unavailable source could affect the selected
classification, both $I_2$ and the relevant candidate remain blocked. The report
can still state precisely that the explicit type-of-product omission was found
and corrected.

This trace illustrates two kinds of progress. The repaired omission supplies
positive evidence about the ensemble's ability to catch that seeded defect in
this designed case. The timeout supplies operational evidence about the
controller's accounting. Neither is a measured error rate on real circulars.
The prototype acceptance test will reproduce these transitions with scripted
members before any live-model comparison is attempted.

\section{A complete dry run: an ambiguity that survives the budget}

Continue with the hypothetical cashback question. One candidate treats a later
payment as a fee discount when it reimburses an actually paid charge and is capped
at that charge. Another requires a reduction in the contractual fee itself.
Their rules agree on a direct waiver and disagree on a later reimbursement.
Neither definition is accepted merely because it produces a convenient outcome
for the bank.

The scenario designer creates three synthetic offers. Offer A waives a charge of
HK\$100. Offer B charges HK\$100 and later pays HK\$100 back under a linked
promotion. Offer C charges HK\$100 and pays HK\$150. The factual amounts are known,
but the legal classification of the later payment remains disputed. The dossier
shows the candidates' conclusions and the exact definition each requires. It
also separates any excess payment into a proposed component analysis rather than
silently treating the entire amount as a discount.

In round one, a retrieval action seeks the incorporated definition. If the source
is inaccessible, the result is a failed retrieval, not support for either
reading. In round two, a bounded argument action asks each candidate to identify
the strongest source-based reason against its own definition. Suppose both
return plausible but unresolved explanations and no new authority. The controller
records the arguments, but a longer explanation alone does not constitute
resolution.

In round three, a question-selection action prepares a focused review question
with the contrasting offers and the missing premise. The action can successfully
produce the question while the issue remains unresolved. After the third round,
no further automatic resolution round is authorized. The candidate set is
preserved, the issue receives a budget stop reason, and the run closes as
\texttt{BLOCKED\_UNRESOLVED}.

The report states which offers are affected, what the candidates agree on, where
they differ, which sources were inspected, which source remains unavailable, and
what an authorized reviewer must decide. It contains no assertion that three
rounds make the leading reading probably correct. The unresolved classification
is not replaced by a false Boolean, and the Java release route remains blocked
for the affected control.

A later authorized reviewer may establish the definition from an applicable
source or decide that a conservative bank policy should govern the uncertain
category. Those are distinct decisions. A bank policy that declines all cashback
offers pending clarification can be operationally useful, but it must be labelled
as the bank's precautionary choice. It cannot be reported as a proven SFC
prohibition unless the source supports that proposition.

\section{The uncertainty report is a first-class result}

An uncertainty report must let a reviewer reconstruct the decision without
reading raw model transcripts. It starts with the source and activity perimeter,
the proposed control, the run's terminal status and the operational consequence.
A blocked report should say what is blocked. ``Low confidence'' is not enough:
the reader needs to know whether the uncertainty concerns an incentive category,
an effective date, a missing definition or the integrity of the entire source.

The report then presents the surviving interpretations in comparable form. For
each, give the controlled-language statement, formal rule identifier, assumptions,
supporting sources, strongest objections and distinguishing scenarios. Use the
same factual scenarios across candidates so that the difference is visible.
Candidate ordering may reflect a declared presentation policy, but must not imply
that the first row has been approved.

Every detected discrepancy appears in an issue register with its processing and
resolution states, materiality, affected decisions, attempted actions and final
reason. Resolved issues remain present because their history explains why the
final candidate differs from the initial one. Unresolved issues identify the
missing evidence or human decision. A report that shows only unresolved issues
would conceal the repairs; a report that shows only repaired issues would conceal
the remaining risk.

Coverage is reported against explicit denominators. Count source units inventoried
and dispositioned, dependencies required and acquired, interpretation families
identified and explored, mandatory roles completed, candidate comparisons
attempted, and verification obligations passed or incomplete. The denominators
must be retained even when the run exceeds a display limit. Large inventories
can be paginated or stored as linked machine-readable records; they cannot be
silently truncated to make the report look complete.

The execution account records configured limits, issued actions, rounds used,
timeouts, unknown remote results, no-progress decisions, elapsed time, token and
cost reservations, actual usage when available, and unspent budget. A run that
stops early because it lacks authority should not look like a run that exhausts
its resources. Both may be blocked, but their next useful actions differ.

A concise machine-readable summary might look as follows. Identifiers here are
illustrative; the complete contract in Chapter~\ref{ch:implementation} adds
validated references and integrity fields.

\begin{lstlisting}[caption={A blocked result preserves alternatives and the next decision.}]
{
  "schema_version": "interpretation.v1",
  "run_status": "BLOCKED_UNRESOLVED",
  "source_ref": "23EC46-retained-version",
  "selected_slice": "paragraph-10",
  "candidate_ids": ["cashback-reimbursement", "cashback-fee-reduction"],
  "material_unresolved_issue_ids": ["classification-cashback"],
  "release_eligible": false,
  "processing_stop": "ROUND_LIMIT",
  "rounds_issued": 3,
  "probability_of_legal_correctness": null,
  "next_decision": "Establish the applicable fee-discount definition",
  "next_decision_owner_role": "LEGAL_INTERPRETATION_REVIEWER"
}
\end{lstlisting}

The report should also state the limits of its evidence. In the second-circular
exercise, the expected cases and candidate were prepared by the same author.
That fact belongs beside the verification summary. A reviewer should not have to
find a footnote in a separate research note to discover that the apparent oracle
was not independent. Similarly, a solver timeout must appear as incomplete
verification, not disappear from an average pass rate.

Reports are immutable versions. A corrected report points to its predecessor and
explains the correction. Approval binds the exact report, candidate, source
packet and Java bundle. The report is thus both a readable explanation and an
index into reproducible evidence. Its human readability matters because a
technically complete JSON object that no reviewer can understand does not solve
the original manual-interpretation problem.

\section{Agreement can still conceal a missed obligation}

The hardest example is not a disagreement. Suppose every generator ignores a
footnote that changes applicability. All formulas, scenarios and Java results
agree because they share the same omission. Repeated self-consistency sampling
would reinforce that agreement. Roundtrip translation could reproduce the same
shortened meaning. Differential execution would confirm that the wrong
specification is implemented consistently.

The independent inventory creates another opportunity to catch the omission. It
requires a disposition for the footnote. A normative reader responsible for that
unit may identify the qualification even if the candidate generators did not.
A source-to-rule coverage check then opens an issue because the footnote has no
corresponding assumption or rule. The repair propagates through candidate
comparison and implementation checks.

Now suppose the inventory also misses the footnote because both extraction paths
use the same defective text layer. Rendered-document inspection or a distinct
acquisition route may detect it. If those are absent or also fail, the system can
still miss the obligation. This remaining possibility is why the product's
claims must be measured against independent source review and adversarial
omission tests. Redundancy reduces opportunities for undetected error only to
the extent that the redundant activities can fail differently.

A second common-mode failure concerns vocabulary. All members may agree that a
payment is a discount because the prompt labels it ``fee rebate.'' The remedy is
to distinguish observed facts from proposed classifications. Present the payment,
charge, contractual relationship and timing without inserting the disputed label
as an established input. Then ask what evidence supports the classification.
Blindness to the preferred answer is useful only if the factual packet itself
has not already embedded that answer.

A third failure concerns evaluation. If the reference reviewer sees the model's
answer before forming an expected outcome, the reference may inherit the model's
mistake. The ensemble may then appear accurate against a contaminated oracle.
Independent initial judgments and recorded adjudication are necessary to measure
whether the product actually discovers errors. They are also part of a credible
operational review process, especially for novel or high-impact controls.

\section{Interaction between uncertainty and operational action}

The interpretation service should not decide every operational response. Its job
is to state what is known, what is disputed and which controls are justified by
the approved specification. The bank's host system then applies an explicit
operational policy to outcomes such as unknown, conflict or blocked release.
Those policies may include stopping an offer, routing a case to review or using
an already approved fallback control.

This separation avoids a dangerous ambiguity in ``fail closed.'' Closing an
uncertain marketing offer may be appropriate in a particular process, but
stopping a payment, withholding a disclosure or disabling an existing control
could create other obligations or harms. The responsible business and control
owners must approve the action for the relevant workflow. The interpretation
engine supplies a reasoned status; it does not invent a universal response for
all banking operations.

For the selected gift rule, a true result indicates that the specified
prohibition condition is satisfied. A false result says only that this condition
is not established under the accepted facts and rule. It does not authorize the
offer under every other requirement. The host must combine this result with other
applicable controls, such as disclosure or suitability obligations, according to
an independently reviewed orchestration policy.

An unresolved interpretation can sometimes be isolated. If the disputed cashback
definition affects only one incentive category and the system can reliably
identify that category, a reviewer may approve a limited control for unaffected
cases while routing the disputed category to manual review. The scope restriction
must be explicit, tested and included in the approved bundle. It is not enough to
state informally that the system will ``avoid edge cases.''

If impact cannot be isolated, the block applies more broadly. The system should
prefer an honest incomplete release proposal over an unsupported claim of
separability. The uncertainty report gives the decision-makers the information
needed to choose between narrowing scope, obtaining more evidence, retaining an
existing control or postponing deployment.

\section{Acceptance conditions for the central method}

The first implementation should be accepted against deterministic scenarios
before live model performance is assessed. A scripted member can deliberately
omit the type route, supply a contradictory source version, fail to respond, or
return a malformed candidate. These fixtures isolate controller behavior from
model variability. Passing them establishes specified engineering behavior, not
legal translation accuracy.

The controller must demonstrate that a unanimous but uncovered source unit opens
an issue; that a supported minority candidate remains in the report; that a
formal mismatch produces a replayable witness; that an unavailable dependency
cannot be marked resolved by repeated explanation; and that a successful repair
reruns affected checks. Each condition targets a concrete way an apparently
successful ensemble could hide an error.

Budget tests must reach exact boundaries. With a cap of twelve, the thirteenth
external dispatch must be rejected even if it is proposed by a new child issue.
A crash after dispatch must not reset the count. A timeout must consume its
issued action. A new worker with an old lease must fail to overwrite the current
result. A provider retry must either reuse a verified idempotency key or consume
a separately authorized attempt.

Report tests must verify completeness for successful, blocked, cancelled and
failed runs. Every issue identifier must have a disposition; every unresolved
material issue must appear in the blocking summary; every candidate revision
must retain its parent and source references; and every displayed score must
identify its purpose. Empty extension sets, solver unknowns and missing member
results must remain explicit rather than being converted to success by default
logic.

The release test is adversarial. Attempt release from
\texttt{READY\_FOR\_REVIEW} without approvals, with an approval for a different
bundle, after a source change, with an unresolved material issue, and after an
integrity failure. All attempts must fail. A positive release fixture requires
complete checks and authenticated approvals for the exact package. The local
synthetic identity mechanism is sufficient to test the state transitions, but
enterprise authentication and role enforcement remain a deployment requirement.

These conditions give an implementation agent a concrete target. They also
preserve the intellectual purpose of the ensemble: investigate alternatives,
expose missing evidence, and retain uncertainty when the evidence does not settle
the question. The resulting product is useful because it makes responsibility
and reasoning visible, not because a collection of fallible components can be
renamed a proof of English meaning.

\chapter{From an accepted interpretation to verified Java behavior}
\label{ch:verification}

\section{The specification becomes the comparison target}

Once reviewers accept a particular interpretation, a different verification
problem begins. The source question was whether the rule expresses the relevant
regulatory requirement. The implementation question is whether the deployed
program preserves that accepted rule for the inputs and histories it supports.
Keeping these questions separate allows strong engineering evidence without
misrepresenting it as legal proof.

The accepted object is larger than a Boolean formula. It includes the source
versions, activity profile, factual definitions, subject identity, time policy,
unknown and conflict semantics, exception structure, result interpretation and
coverage limits. If any of these is omitted, two implementations can agree on a
formula while giving different advice in practice. A Java adapter that labels
all rebates as discounts has changed the rule's application even if its Boolean
operator code is flawless.

The verification target must therefore name the complete bundle and its
permitted input domain. The bundle is immutable and content-addressed. The
comparison includes the decision, missing and blocking facts, errors, source
references and execution trace. Different backends have different execution
identities, so equality of their entire serialized result hashes is not the
right criterion. Compare the specified semantic fields and independently verify
each backend's own result identity.

The existing workbench already provides a finite RuleIR interpreter, a Java
runtime, generated policy classes and a release workflow. The ensemble proposed
in Chapter~\ref{ch:ensemble} is an extension before that delivery path. It should
produce reviewed candidates that the existing validator can accept, rather than
silently expanding the runtime language whenever a model proposes a new legal
construct. Unsupported meaning remains an explicit engineering and review issue.

\section{Facts are evidence-bearing values}

A production fact is not simply a field such as \texttt{gift=true}. It has a type,
value or uncertainty status, evidence references, validity interval and recording
time. The validity interval says when the evidence applies to the underlying
situation. The recording time says when the system knew it. These times answer
different questions and both are needed for historical reconstruction.

Suppose the bank learns on 10 October that a fee record effective on 1 October
was incorrect. A decision replay asking what the system knew on 5 October should
not use the later correction. A corrected assessment made on 10 October may use
it to reassess the earlier event. The two assessments should coexist. Replacing
the old record would erase the evidence needed to explain the original decision.

The existing normalization boundary selects records valid at \texttt{valid\_at}
and recorded no later than \texttt{known\_at}. A correction names the evidence it
supersedes. Two admissible, non-superseded records with different values create a
conflict. Arrival order is not a source-priority rule. If both records agree,
their evidence may be combined without inventing a conflict.

Intervals are half-open: $[a,b)$ includes the start and excludes the end. If a
record is valid until midnight, it is not valid at that midnight. This convention
avoids overlapping adjacent intervals when one record ends exactly as another
begins. Tests must exercise equality at both endpoints, not merely dates well
inside an interval. All assessment timestamps use the canonical UTC format with
six fractional digits; legal calendar dates are separate values.

Unknown facts have reasons. Missing evidence, stale evidence, an unreviewed
classification and unresolved event order can all lead to unknown, but they call
for different remedies. A host should not display the same uninformative message
for every case. Conflict similarly preserves the competing evidence identifiers.
The evaluator does not decide which source is legally superior without a
reviewed normalization rule.

This design puts factual interpretation where it belongs: at a visible boundary.
For the circular, classifying a benefit as a gift or fee discount is an evidenced
assessment. The decision engine consumes that assessment under its approved
definition. If the ensemble has not resolved the definition, the adapter cannot
supply a known value merely to satisfy a required input field.

\section{Exact arithmetic prevents a second kind of ambiguity}

The existing scalar types are Boolean, integer, Hong Kong dollar money in cents,
and Gregorian date. Integer and money values use canonical decimal strings and
arbitrary-precision arithmetic. A monetary value of HK\$100 is represented as
\texttt{"10000"}. The string representation avoids a transport layer converting
the value through a floating-point number before the exact arithmetic begins.

Units are part of the contract. A number expressed in dollars and the same number
expressed in cents differ by a factor of one hundred. A compiler cannot infer
which was intended from the digits alone. The host mapping must identify the
source unit, conversion and evidence. Declaration-specific bounds can reject
implausible inputs, but a plausibility check is not a substitute for a reviewed
unit mapping.

Consider proportional ownership of a monetary amount. If the accepted operation
is exact scaling by a rational fraction $n/d$, the runtime computes the product
and quotient using integers. Let $x$ be cents. It checks
\begin{equation}
 nx=dq+r,\qquad 0\leq r<d
 \label{eq:exact-scale}
\end{equation}
for nonnegative $nx$; for signed values the implementation's quotient and
remainder convention must still make zero remainder the exactness test. A
nonzero remainder means that the result is not an integer number of cents. The
current contract returns \texttt{E\_INEXACT\_SCALE}; it does not choose a rounding
policy on the user's behalf.

If a bank requires rounding, that is a substantive extension. The rule must state
which quantity is rounded, at what stage, in which direction and under whose
policy. Rounding each component before summation can differ from rounding the
sum. Two correct implementations of different policies can therefore disagree.
The ensemble should recognize this as a specification issue rather than treating
one result as a numerical accident.

Foreign-currency conversion is outside RuleIR 0.1. A reviewed adapter supplies an
HKD amount with rate, valuation instant and rounding evidence. A missing rate
creates unknown evidence. It must not trigger a guessed conversion or reuse an
unboundedly old rate. The bank's data owner is responsible for the freshness and
valuation policy, while engineering verifies that the adapter implements it.

Dates require similar precision. A date is not a timestamp, and an obligation
measured in business days needs a calendar policy. The current expression
language compares validated dates but leaves date arithmetic to reviewed event
scheduling. This restricted design prevents a generic arithmetic library from
silently deciding whether a public holiday counts toward a regulatory deadline.
The appropriate calendar must be versioned with the assessment.

\section{Unknown, conflict and scope have an order}

The three-valued Boolean rules in Chapter~\ref{ch:proof} govern unknown values,
but they are not the entire runtime semantics. Before evaluating the requested
rule, the existing implementation checks conflicting evidence in its static
transitive dependency closure. That closure includes scope, body, exceptions and
branches that might later be skipped. A conflict in the closure returns conflict
even if another route could determine the Boolean result.

This is a conservative evidence policy. For example, an actor may appear outside
the selected scope, but a conflicting fact in the rule's static closure still
produces conflict before scope evaluation. The second-circular case that exercises
this behavior is deliberately different from ordinary short-circuit intuition.
The specification must teach it because an engineer using conventional Java
\texttt{\&\&} or \texttt{||} could otherwise implement a different result.

After conflict checking, scope is evaluated. False scope returns
\texttt{OUT\_OF\_SCOPE}; unknown scope returns \texttt{UNKNOWN} with a scope reason;
true scope permits evaluation of the body. A body that would be false does not
turn unknown scope into a known result. The system is answering a scoped
condition, and the applicability question remains unresolved.

Within the body, conjunction and disjunction evaluate all operands in declared
order. This eager behavior retains complete evaluated evidence and exposes errors
or conflicts even if an earlier operand would determine an ordinary Boolean
answer. Missing inputs can appear in a known result, while blocking inputs are
empty when the unknowns do not prevent that result. The distinction is useful:
the bank may know the condition is false while still needing to repair a stale
input feed.

A conditional expression behaves differently. It evaluates the condition and
only the selected branch. An unknown condition returns unknown without evaluating
either branch, even when the branches appear equal. Skipped branches remain type
checked and part of the static dependency closure. Their trace roots are marked
skipped so that the explanation does not invent an evaluation that never occurred.

Defaults add another layer. All exception guards are evaluated first. Two true
guards produce conflict even if their values would agree. An unknown guard can
prevent selection of another true guard because the unknown might become a
competing exception. One true guard with all others false selects its value; all
false guards select the base. Priority is expressed only by explicit nesting with
a reviewed interpretation. Array order is not hidden legal priority.

These rules are not universal truths about law or logic. They are the existing
workbench's specified execution choices. The ensemble may identify a regulation
that needs different semantics. The correct response is to propose and review an
extension or keep the case outside the supported fragment, not to reinterpret
the current operators ad hoc.

\section{Validation failures are part of observable behavior}

A verifier must compare invalid-input behavior as well as successful decisions.
Malformed JSON, duplicate keys, noncanonical numbers or unpaired Unicode
surrogates can fail at the transport boundary before a semantic result can be
hashed. Valid JSON with an invalid bundle or snapshot can instead produce a
structured semantic error. These outcomes have different evidence and should
not be conflated.

The current validation order is schema, references and duplicates, types, cycles,
time, static-closure conflicts and evaluation. A deterministic order makes error
reports reproducible when an input has several defects. It also prevents a
runtime from accessing an invalid reference before discovering that the bundle
should have been rejected. Fixed error codes and pointers are compared across
backends; dependency-library exception text is not the public contract.

Resource limits are explicit. The present canonical JSON boundary is bounded in
size and depth, and the rule graph has node and expanded-call-depth limits. An
oversized acyclic expression is a resource error; a cycle is a cycle error. An
implementation must not catch either and continue with a truncated rule. The
same principle applies to the new ensemble's graph and report limits.

Tests should combine defects when precedence matters. A snapshot can contain a
conflict and an expired validity interval. A bundle can contain a type error in a
branch that would otherwise be skipped. A default can have two applicable guards
whose values contain errors that should never be evaluated. Isolated operator
tests cannot establish the behavior of these combinations.

A useful specification test states both the expected result and why a tempting
implementation would differ. For example, a short-circuit implementation might
return false before encountering an error in the second operand. The accepted
eager semantics requires the error. This counterexample teaches the contract and
makes the test more durable than an assertion copied from the current function's
output.

\section{Independent oracles and their limits}

An oracle is a source of expected behavior against which an implementation is
compared. Independence has several meanings. A separately written evaluator can
be independent of the production code while relying on the same mistaken
interpretation. A human reference can be independent of the model's answer while
sharing the same incomplete source packet. The verification report must specify
which independence it has actually achieved.

The retained 22-case oracle was frozen before the candidate bundle, which protects
against some forms of adapting expected outputs to the implementation. It was
nevertheless prepared by the same author. The 729-case enumeration gives complete
coverage of six ternary abstract inputs for the selected formula. Neither fact
establishes independent legal interpretation or whole-document coverage.

A stronger reference process begins with independent source reading and scenario
judgment. Reviewers state the expected result, assumptions, supporting passage
and unresolved questions before seeing the candidate output. Disagreements are
adjudicated with reasons or retained as alternative acceptable outcomes. The
reference then has a defined relationship to the evidence rather than being a
single unexplained label.

For engineering checks, a small executable oracle can be valuable when its method
is genuinely different. The second-circular body has a read-once unate structure,
which permits its Boolean-completion check to agree with the specified
three-valued result. The general expression $p\lor\neg p$ shows why that argument
cannot be extended to every RuleIR expression: under the specified semantics,
unknown $p$ yields unknown, while every classical completion yields true.
An oracle based on the wrong generalization would reject correct implementations
or bless incorrect ones.

The oracle itself therefore needs a stated domain and a small set of discriminating
examples. The workbench should retain oracle author, method, source versions,
review status and any shared components. A generated oracle that imports the same
evaluator under a different function name supplies little independent evidence.
A reference that is changed during repair remains useful only when the change is
separately justified and versioned.

\section{Mutation tests measure what the checks can notice}

A mutation deliberately changes a candidate or implementation in a way that
represents a plausible error. If the verification process rejects it for the
intended reason, the test demonstrates sensitivity to that error. It does not
estimate the frequency of that error in real use, and a high mutation score does
not prove complete correctness.

The second-circular exercise used four meaningful mutations: remove the discount
exception, omit product-type linkage, assume gift classification and ignore scope.
The retained named cases detect all four after compilation. This is more
informative than mutating arbitrary characters because each change corresponds
to a failure a reviewer or programmer could plausibly make.

The ensemble extension should add source and process mutations. Remove a footnote
from one extraction; supply a stale dependency to one member; label a speculative
classification as established; suppress a minority candidate in the report; reset
an issue counter after a revision; or reuse an approval after a source change.
Each mutation targets a distinct assurance boundary. Passing only formula
mutations would leave the new controller largely untested.

Some mutations are equivalent under the selected domain. Changing an expression
may not change any supported outcome, and failure to detect that mutation is not
necessarily a defect. The mutation report must distinguish equivalent mutations,
invalid encodings and live semantic defects. Counting all undetected mutations as
failures, or silently excluding inconvenient ones, would distort the result.

A mutation is killed when a required property changes and the check catches it.
A compiler crash is not always the intended detection: if the task is to test
whether the oracle notices an omitted exception, a syntax-invalid mutant does not
answer the question. Mutation construction should preserve validity where the
research question concerns semantic sensitivity. Invalid-input tests remain a
separate and necessary category.

\section{Symbolic comparison requires a feasible domain}

For supported expressions, symbolic comparison can search for inputs on which two
candidates differ. The domain records admissible values, knownness conditions and
relationships among facts. Before asking whether the outputs differ, the solver
must establish that the domain has at least one model. Otherwise an inconsistent
domain makes every equivalence query vacuously true.

Suppose one domain constraint says that a payment exceeds the fee and another
says that it cannot exceed the same fee. A query finding no differing outcome
under both constraints has taught us nothing about the real candidates. The
workbench returns an inconsistent-domain result. The domain must be repaired or
reviewed before any equivalence claim is made.

A satisfiable difference query returns a witness. The witness must be decoded to
a concrete fact snapshot and replayed in both the reference evaluator and Java.
This replay checks that the symbolic encoding corresponds to executable semantics.
A solver result alone can be wrong relative to the claimed runtime if the
encoding mishandles unknown values, scope or arithmetic. The witness and both
results belong in the report.

An unsatisfiable difference query establishes no differing result within the
checked fragment and domain, subject to the solver and encoding assumptions. It
does not establish that the candidates use the same legal definition, cover the
same source material or behave identically outside the domain. The result label
must therefore include the domain and comparison projection.

The current comparison fragment supports documented Boolean and exact
integer/money reasoning with explicit knownness. It rejects unsupported date,
scale, default and event constructs. The two solver checks share a bounded time
budget, with a separate worker deadline. Unknown, timeout, missing Java replay or
unsupported constructs never become equivalence. The ensemble may schedule
another permitted check, but cannot relabel the incomplete result as a pass.

Counterexample-guided refinement is useful when a witness exposes an imprecise
abstraction. The engineer identifies why the witness is infeasible or why the
encoding omitted a required condition, changes the abstraction with a documented
reason, and reruns the query. The discipline matters: excluding a witness because
it is inconvenient is not refinement. The revised domain must correspond to a
reviewed fact about the application.

\section{Property tests should follow the semantics}

Property-based tests generate many inputs and check a relation that should hold.
They can explore combinations that named examples miss, but the property must
be valid for the specified semantics. Classical Boolean identities cannot all be
transferred to unknown or conflict behavior. Likewise, reordering operands may
preserve a value while changing the required trace or error precedence.

Useful properties include deterministic evaluation for identical canonical
requests, rejection of undeclared facts where the schema requires it, exact
integer arithmetic, preservation of known values under irrelevant evidence
additions, and stable source references under deterministic rendering. Each
property needs a precise precondition. Adding a conflicting record in the static
dependency closure is not an irrelevant addition.

Metamorphic tests compare related inputs when a direct oracle is difficult. If a
known amount increases while all other facts remain fixed, a simple lower-bound
threshold condition should not change from true to false. This monotonicity
applies only to the stated threshold fragment. It need not hold for a default
with an upper-bound exception or for a scope condition that changes at the same
time. The test generator should derive its applicable domain from the property,
not assume the property is universal.

For evidence time, a carefully limited relation is useful: replaying an earlier
\texttt{known\_at} should exclude records first recorded later. Increasing
\texttt{known\_at} does not necessarily preserve the decision, because a correction
or conflict may become visible. A test that demands decision monotonicity in
knowledge time would impose a false property. The intended invariant concerns
which records are admissible, not whether later knowledge always makes the same
answer more certain.

Fuzzing transport boundaries is a different activity. It tests duplicate keys,
size limits, malformed timestamps, unusual Unicode and nesting limits. Its
success criterion is controlled rejection or valid deterministic behavior, not
legal interpretation. These tests protect the integrity of the evidence and
execution chain and should be included in the engineering report under that role.

\section{The actual generated Java interface}

The current emitter validates the bundle, canonicalizes it and embeds its exact
bytes in a generated class. The class name is
\texttt{hk.legalmath.Policy\_} followed by the first twenty hexadecimal digits of
the bundle hash. The class exposes the full bundle hash and delegates evaluation
to the Java semantic runtime. This is a generated, immutable policy binding; it
is not the older proposal's imagined record-based API or a proof that each legal
rule has been translated to hand-optimized Java control flow.

The two current entry points are:
\begin{lstlisting}[language=Java,caption={Implemented generated-policy API.}]
public static String evaluate(
    String snapshotJson,
    String ruleId,
    String validAt,
    String knownAt,
    String mode);

public static Map<String,Object> evaluate(
    Map<String,Object> snapshot,
    String ruleId,
    String validAt,
    String knownAt,
    String mode);
\end{lstlisting}

The string overload returns strict JSON. The map overload returns a deeply
immutable result. The generated policy and Java runtime operate without Python,
a language model, a solver or network retrieval. This property is important for
a bank integration: the interpretation workbench prepares and verifies a package,
while production decision execution can remain local and deterministic.

A concrete host compiles against the actual generated class name from the build
manifest. It supplies a validated snapshot, exact rule identifier, both assessment
times and the requested mode. The mode does not grant authority. A caller cannot
turn a draft into an approved control by passing the word ``production.'' The
release service and host must establish approval and applicability separately.

The conceptual integration looks as follows. The symbolic class name in the
listing is replaced by the actual generated name when the integration fixture is
built; the release lookup and transaction methods are host responsibilities.
\begin{lstlisting}[language=Java,caption={Host responsibilities around deterministic evaluation.}]
ApprovedRelease release = releases.resolve(profile, validAt, knownAt);
release.verifyAuthorityAndDigests();
ConsistentSnapshot s = adapter.readSnapshotWithVersions(subject, order);
String resultJson = Policy_HASH.evaluate(
    s.strictJson(), release.ruleId(), validAt, knownAt, "production");
Decision decision = decoder.decodeStrictly(resultJson);
HostDisposition disposition = hostPolicy.map(decision);
transactions.commitIfVersionsUnchanged(
    s.versions(), order.id(), order.payloadHash(),
    release.identity(), decision, disposition);
\end{lstlisting}

\texttt{ApprovedRelease}, \texttt{ConsistentSnapshot} and the host transaction
interfaces in this illustration are proposed integration concepts, not classes
already supplied by the generated SDK. Their responsibilities must be implemented
in the bank's environment. The example makes the boundary explicit so an agent
does not mistake illustrative types for an existing dependency.

The host must handle every tagged result. A true selected prohibition condition
may require stopping the relevant offer; false permits only progression to other
controls under the host policy. Unknown, conflict, out-of-scope and error each
have explicit dispositions. No result string should be cast to a Boolean by a
language's truthiness rules. The decoder rejects unknown status tags so that a
future runtime extension cannot silently take an old default branch.

\section{Traces are part of the delivered meaning}

A decision without an explanation is difficult to investigate. The runtime trace
records node identifiers, operators, children, type, value or status, evidence
references and source spans. Shared rule evaluations are memoized, and skipped
branches have explicit markers. The trace is deterministic under the request and
engine version, making it suitable for replay and comparison.

Trace fidelity is a separate verification obligation. Two programs may return
true while one omits the evidence for the exception or incorrectly claims to have
evaluated a skipped branch. A regulator-facing explanation assembled from that
trace could be misleading even though the Boolean result matches. The comparison
therefore includes trace structure and evidence, not just the top-level status.

Natural-language explanations are generated from the verified trace and approved
source mappings. A language model may help render a readable summary, but it
cannot alter the computed result or invent a justification. If a summary mentions
a source that the trace and bundle do not contain, validation should reject or
flag it. The authoritative explanation remains reconstructible without the
summary model.

The current source-span policy associates trace nodes with the owning rule's
explicit citations and retains exception citations in the bundle. That is a
coarser mapping than a separate legal proof for every arithmetic operation. The
report should describe the granularity honestly. A future refinement may attach
more precise spans, but a dense set of citations should not be mistaken for a
formal demonstration that each mapping is legally correct.

Execution hashes identify the canonical request and result under a particular
backend version. They help detect altered bytes and reproduce behavior. They do
not authenticate a reviewer or prove that a document came from the regulator.
Authentication, source acquisition and approval are separate controls. Confusing
these roles is a common source of overclaiming in systems described as auditable.

\section{Event histories add obligations that a snapshot cannot express}

Some requirements concern sequences: obtain consent before an action, honor a
withdrawal, provide a document within a period, or preserve a breach after late
performance. A current Boolean snapshot cannot express all of these relationships.
The existing event component replays finite attributed histories using a separately
versioned Java entry point, \texttt{EventReplay.replay(Map<String,Object>)}.

A stream header binds subject, category, actor, action, profile and inception.
Events belong to that header; an API path cannot redirect their meaning to a
different subject. Duplicate identical events are idempotent, while conflicting
payloads or sequence reuse are errors. Ordering authority is explicit. If event
order is unresolved, the system must not invent a sequence based solely on arrival.

Consent illustrates why completeness matters. Absence of a grant in an incomplete
history does not establish that no grant occurred. A completeness watermark or a
reviewed initial state supplies the relevant observation boundary. A watermark
cannot cover a time later than the recording instant of its assertion. A
withdrawal ends the applicable consent generation, and a later regrant uses a new
evidence identity rather than resurrecting a previously withdrawn record.

For an achievement obligation, activation opens a duty with a deadline. The
selected contract uses the interval from activation up to, but excluding, the
deadline for timely performance. Performance exactly at the excluded endpoint is
therefore late under that contract. This is an engineering convention that must
match the reviewed legal timing rule; a different legal rule requires a different
specified boundary, not a casual adjustment in the host.

Late performance does not erase the historical breach. The system can record both
that a duty was breached and that it was later fulfilled. Similarly, an observed
action is retained even if an advisory control would have prohibited it. Deleting
inconvenient observations would make replay describe an idealized compliant world
rather than the bank's actual history.

Each opened duty binds the source bundle that created it. A later amendment does
not automatically migrate existing duties. Migration may be appropriate, but it
requires an explicit legal and operational policy about which obligations change,
which dates apply and how historical findings are retained. The current MVP has
no automatic migration implementation. The ensemble can identify the question
and prepare alternatives; it cannot assume the answer from the new publication
date alone.

\section{Build, verification and approval form an acyclic chain}

A build manifest records the bundle, generated source, runtime source, toolchain,
Java target, dependencies and JAR digest. A verification report then names the
build and exact JAR it tested. A release manifest names that build and verification
report, together with applicability and effective-time information. Approval
binds the final release manifest and bundle.

The direction matters. Embedding a manifest that contains the JAR's own digest
inside that JAR creates a circular identity problem. The current design keeps
these records external and acyclic. A reviewer approves the exact tested bytes
through the final manifest, rather than approving a mutable directory or a
human-readable version label.

Candidate compilation occurs before legal approval, because reviewers need a
concrete package and verification evidence to inspect. Compilation does not grant
deployment authority. Export into a controlled release route occurs only after
the required approvals and checks. This distinction lets engineers do useful
reversible work without repeatedly asking for permission, while preserving the
human decision at the point where it matters.

Reproducible builds strengthen the identity chain. Under the same declared
toolchain and inputs, build in separate directories with fixed archive-entry
timestamps and compare emitted source and binary hashes. A mismatch should be
investigated and recorded. Reproducibility does not prove semantic correctness,
but it reduces uncertainty about which bytes correspond to the reviewed source.

The local MVP uses synthetic identities and unsigned historical archives to
exercise workflow. A bank deployment needs enterprise identity, role separation,
key management and independently verified release signatures or equivalent
controls. A hash proves equality of bytes under its cryptographic assumption;
it does not prove who authorized those bytes. The manuscript's implementation
plan keeps that integration distinct from the local demonstration.

\section{A correct decision can still race with the transaction}

Suppose a control depends on consent and exposure. The host reads both, evaluates
an order and then commits it. Between the read and commit, consent may be
withdrawn or another order may consume the available limit. The decision was
correct for its snapshot but is no longer valid for the transaction being
committed. This is a concurrency problem outside the pure rule evaluator.

The host should read a consistent snapshot with version identifiers, evaluate
exposure including the proposed order, and atomically validate the versions while
reserving the exposure and recording the decision. If the versions changed, it
rereads and reevaluates under a bounded transaction retry policy. That policy is
separate from interpretation-resolution retries and must not share an ambiguous
counter called simply ``attempts.''

Idempotency protects against duplicate submissions. Repeating the same order
identifier and payload should recover the same committed result or reservation.
Using the same identifier with a different payload is a conflict. The identity
must include the legally relevant order content, not merely a user-interface
request number. Otherwise a retry can accidentally authorize a changed transaction.

Release selection has a similar race. The host records the release identity used
for the decision and verifies that the applicable release policy permits it at
commit. Whether an in-flight transaction can finish under a prior release is a
bank policy requiring explicit treatment. The rule engine cannot resolve that
question by consulting its process clock or loading whatever JAR is newest.

A rollback also has legal meaning. Restoring a previous software build can repair
an implementation defect if that build still expresses the applicable rule.
Restoring a superseded legal regime is different. The incident procedure must
separate software rollback, operational suspension and legal-rule selection.
Historical packages remain available for replay even when they cannot be used for
new decisions.

\section{A verification dossier that supports a release decision}

The final technical dossier links the accepted interpretation to its implementation
through several kinds of evidence. Named scenarios demonstrate reviewed boundary
behavior. Exhaustive finite checks cover declared small domains. Property and
mutation tests examine broader classes of implementation mistakes. Symbolic
comparison supplies domain-qualified results and replayed witnesses. Differential
execution compares Python and Java semantics. Build records identify exact bytes.
Host tests exercise concurrency, idempotency and release selection.

These methods overlap but are not interchangeable. A compiler proof could cover
an unbounded class of accepted expressions under stated semantics, while still
assuming that the factual adapter and source interpretation are correct. A test
suite can catch a concrete adapter error that lies outside that proof. An
independent legal review can reject a perfectly verified expression because it
misreads an exception. The dossier should state the role of each result so that
one strong component does not hide a missing one.

For the current project, the retained engineering evidence includes the local
MVP test suite and the separate second-circular exercise. The new ensemble is not
yet an implemented live-model product. The appropriate next engineering step is
to implement its controller and contracts, then run deterministic acceptance
scenarios, then evaluate interpretation performance with independent reference
judgments. Declaring the entire product ready because the Java path works would
skip the central problem this book is designed to address.

A release reviewer should be able to answer four concrete questions from the
package: what rule is being enforced; what evidence makes its inputs meaningful;
what demonstrates that these exact bytes preserve the rule; and what happens when
an assumption or input is uncertain. If any answer requires guessing from a model
confidence score, the dossier is incomplete. Chapter~\ref{ch:implementation} turns
these requirements into service records, operations and ordered implementation
tasks.

\chapter{An implementation contract for the interpretation workbench}
\label{ch:implementation}

\section{Extend the existing product at a defined boundary}

The current repository already has source intake, RuleIR validation and execution,
event replay, Java generation, review records, build and release manifests, and a
local service interface. Its drafting provider is a deterministic fixture. The
new work is an interpretation investigation service between source intake and
approval of a candidate bundle. Replacing the entire runtime would discard useful
engineering evidence and enlarge the task unnecessarily.

The extension accepts a retained source packet, an activity profile and an
explicit investigation policy. It creates inventories, candidates, arguments,
discrepancies and bounded actions. Its output is a report and, when appropriate,
one or more validated draft bundles. The existing review and release services
remain responsible for approval and export, with an additional requirement that
the interpretation report and its unresolved issues be checked.

This boundary has two advantages. The interpretation service can evolve its
candidate-generation methods without changing deterministic production execution.
The Java path can be verified independently of any particular model provider.
A candidate that requires an unsupported operator is returned as an engineering
question rather than forcing a runtime change through an unreviewed model output.

The companion contracts define the proposed \texttt{interpretation.v1} extension.
They are stored in \texttt{docs/monograph/contracts}. These are design contracts
and examples, not
a claim that the service is already implemented. The accompanying ordered task file identifies dependencies and acceptance conditions. An
implementation agent should read this chapter with those files and the existing
v0.1 runtime contracts; it should not infer additional authority from illustrative
code in the manuscript.

\section{Separate durable records from generated prose}

The service needs a small set of durable record types. A run binds the source
packet, activity profile, policy and required ensemble roles. A candidate binds a
proposed meaning to a formal representation and its assumptions. An issue records
a discrepancy and its evidence need. An action records one authorized external
operation. A report reconciles the run's evidence and terminal status. Review
decisions and source coverage records connect these objects to human authority
and document completeness.

Generated prose is a field within these records, not their controlling structure.
A model may write an explanation saying an issue is resolved, but the service
changes the issue's resolution status only through a validated transition with
the required evidence. A sentence saying ``approved by compliance'' has no effect
unless an authenticated review record exists for the exact object.

Identifiers are stable opaque strings within a namespace; content hashes identify
immutable payloads. These roles should not be confused. A run identifier remains
stable while events accumulate, whereas each immutable candidate revision has a
new content identity. A human-readable title can change without becoming the
primary key. References use identifiers and, where integrity requires it, the
expected content hash.

Timestamps use the project's canonical UTC format. Durations are integer
milliseconds or seconds with explicitly named units. Counts are nonnegative
integers. Monetary reservations use exact minor units and a named billing
currency. No field called \texttt{cost} should leave the unit or estimated-versus-
actual distinction implicit. This discipline prevents the controller from
comparing incompatible budgets.

The schemas reject unknown properties. This catches misspelled limits and
unrecognized status fields rather than ignoring them. Structural validation is
only the first step: the service also verifies references, state transitions,
source compatibility, budget arithmetic and role authority. JSON Schema cannot
prove that a source passage supports an interpretation or that two referenced
objects belong to the same run.

\section{The run and policy records}

A run records its immutable configuration hash, source packet hash, profile
identifier, required roles, creation time, deadline and current lifecycle state.
It also references its candidate, issue, action and report collections. Mutable
projections such as counters are derived from durable events or updated
transactionally with them; they are never trusted solely because a client submits
a number.

The lifecycle begins with initialization, proceeds through initial proposals and
resolution rounds, and ends in ready-for-review, blocked-unresolved, cancelled or
failed-integrity. A ready-for-review run is complete enough to present, not
released. A run that is blocked can still contain useful validated candidates.
The API returns them with the block, rather than suppressing all information
because the overall investigation is incomplete.

The policy supplies all material resource and coverage choices. It names required
roles, maximum initial actions, maximum resolution rounds, per-root and global
action caps, no-progress limit, action timeout, run duration, candidate count,
argument-node count and dependency depth. It also declares whether monetary and
token caps are enforced. Configuration validation checks internal consistency,
such as requiring the initial-action cap not to exceed the global cap.

The reference fixture policy is deliberately explicit:
\begin{lstlisting}[caption={Reference limits for deterministic controller fixtures.}]
{
  "schema_version": "interpretation.v1",
  "policy_id": "fixture-policy.v1",
  "use": "DETERMINISTIC_FIXTURE_ONLY",
  "max_initial_actions": 4,
  "max_resolution_rounds": 3,
  "max_actions_per_issue_root": 3,
  "max_actions_total": 12,
  "max_no_progress_rounds": 2,
  "action_timeout_seconds": 30,
  "run_deadline_seconds": 600,
  "max_candidates": 24,
  "max_argument_nodes": 200,
  "max_dependency_depth": 6,
  "token_cap": null,
  "cost_cap_minor_units": null
}
\end{lstlisting}

The candidate, graph and dependency bounds make local processing tractable for
fixtures. Reaching one opens or preserves an incompleteness issue. These numbers
are convenience choices for engineering tests, not selected production defaults.
A paid provider configuration must add enforceable cost and token reservations,
and its deployment policy must explain why the chosen limits are sufficient for
the intended workload.

A policy version is immutable once used. Changing a cap during a run requires an
explicit authorized amendment event if that capability is later implemented;
the first prototype should instead require a successor run. This keeps budget
semantics simple and prevents a worker from extending its own authority. The
successor links to the prior report and explains why additional work is justified.

\section{The candidate record carries the interpretive commitment}

A candidate has a parent revision, source references, subject unit, controlled-
language statement, formal bundle reference, assumptions and interpretation
family. The family identifies a consequential choice, such as component versus
whole-offer assessment. It is not merely a cluster label assigned by a text
embedding. A candidate can belong to several dimensions, but its changed choices
must be explicit relative to its parent.

Each assumption has a statement, provenance and status. Possible statuses include
supported by inspected source, factual evidence required, provisional modeling
choice, disputed and rejected. A known fact in a runtime snapshot must not be
created directly from a disputed assumption. The candidate-to-fact mapping is a
reviewed transformation with its own evidence requirements.

Supporting and opposing arguments reference source spans or other arguments.
An argument distinguishes premises, inference rule and conclusion, and labels
strict versus defeasible inference where that distinction is used. If an
argumentation engine is not yet implemented, these records can still support a
structured review interface. The service must not label the resulting graph as
formally evaluated until the selected semantics has actually been run.

Candidate status distinguishes active, deferred, defeated, unsupported and
selected-for-review. Deferred means not yet sufficiently investigated, not false.
Defeated requires a recorded reason under the chosen argument and authority
policy. Unsupported means the current record lacks adequate support. Selected
for review is a presentation decision; only a separate authenticated decision can
approve the candidate's meaning.

The record includes score components with definitions and calibration status.
The initial contract permits a scheduling priority and categorical evidence
states. It explicitly leaves probability of legal correctness null. Later
calibration can add a versioned, narrowly defined score, but cannot reinterpret
old uncalibrated priorities as probabilities after the fact.

Source and assumption identity survive formal deduplication. If two candidates
share an identical RuleIR body, the body can be stored once. Their distinct
assumptions and source mappings remain separate candidate records. This prevents
a compiler optimization from erasing the very disagreement the workbench is
intended to expose.

\section{The issue and action records make automatic work accountable}

An issue names its root, parent if any, run, source units, affected candidates,
semantic dimension and materiality. It records the observed discrepancy and the
question that would resolve it. Processing state and resolution state are separate
fields. A terminal processing state can coexist with unresolved meaning.

Resolution evidence is typed. A source-based resolution references retained spans
and the reviewer or deterministic rule that accepted them. An adjudication
references the authenticated decision. An equivalence resolution references the
comparison report, domain and confirmation that no distinct source commitment
remains. A bare model statement is not a valid resolution-evidence type.

An action names the run, issue root or initial role, round, permitted action type,
input hashes, adapter, timeout, reservation and current state. It also stores the
idempotency key, lease owner, fencing token and terminal result reference. The
action's question and success criterion are immutable after dispatch. Changing
them creates another counted action.

The action result separates transport success from substantive progress. A model
request can return valid JSON yet fail to provide any supported new proposition.
A retrieval can succeed technically while returning the wrong version. A solver
can complete normally with an unknown result. The adapter reports its technical
outcome, and the controller validates whether the result answers the action's
question.

External text is untrusted data. It cannot name a new tool and cause execution,
change a budget, alter role permissions or select an arbitrary filesystem path.
Permitted actions come from a server-side registry. The model can propose an
action from that registry, but the controller validates its applicability and
reserves resources before dispatch. This is both a security boundary and a way to
keep the method's stopping proof connected to actual behavior.

Typed evidence records bind their content hash, source context, validation state
and origin action or review. Source spans, factual records, validated repairs,
formal comparisons and human adjudications are distinct evidence kinds. A formal
comparison additionally names its domain and result; an unknown result cannot
close an equivalence issue. The companion checker rejects unvalidated closure
and equivalence claims supported only by a repair narrative. Establishing legal
sufficiency and authenticated review authority remains a service and review
responsibility beyond JSON validation.

The issue record retains attempted actions even when they fail. It is often useful
to know that an official dependency was unavailable at the assessment cut-off,
that two classification prompts produced no new evidence, or that a formal check
was outside the supported fragment. Those facts explain why the investigation
stopped and what would justify reopening it.

\section{Coverage records connect source units to consequences}

A coverage record maps each inventoried unit to one or more dispositions:
implemented rule, contextual explanation, definition dependency, separate
obligation, reviewed out-of-scope provision, or unresolved question. The record
includes the reason and responsible reviewer where a legal judgment is involved.
A paragraph can have more than one role; it may contain both a definition and an
operative condition.

Coverage validation first checks that the record refers to the same source hash
as the inventory. It then checks that every inventoried unit appears, every
referenced rule or issue exists, and any claim of implemented coverage has a
validated bundle and verification evidence. It cannot decide the substantive
adequacy of a legal disposition from structure alone, so that judgment remains
an explicit review obligation.

Dependencies are edges with types and versions. An incorporated definition is
stronger than a background analogy in the sense that missing it can directly
prevent interpretation of the selected rule. The impact service uses these edges
to identify affected candidates when a source changes. A generic document list
would not tell the service which conclusions depend on which definitions.

The source inventory itself is versioned and reviewed. An improved extractor may
find a previously missed footnote without any change to the official bytes. That
creates a new derivative and inventory revision, not a new official source
publication. Candidates whose coverage was assessed against the older inventory
need reconsideration. Distinguishing source revision from extraction revision
makes that event explainable.

For the second circular, the executable scope remains paragraph 10. The coverage
record must retain all eighteen paragraphs and four footnotes with their
respective dispositions. A machine-readable count of one implemented paragraph
and seventeen other paragraph dispositions is more honest than a banner saying
``circular translated'' without a denominator. The report can still present the
selected slice as a successful engineering example while showing why whole-
document release is not established.

\section{Storage transactions and event history}

The first implementation can extend the existing transactional store. Tables for
runs, candidates, issues, actions, coverage revisions, reports and review decisions
store canonical payloads with indexed identity and relationship fields. Foreign
keys protect references, unique constraints protect idempotency, and transactions
protect coupled state changes. The exact storage engine can remain the existing
local choice for the prototype; bank deployment requires its own availability and
administrative-control assessment.

The most important transaction reserves an action. It reads the current run and
root counters under a lock or equivalent serializable mechanism, checks all
limits, inserts the action and outbox row, increments counters and appends the
audit event. If any step fails, none commits. A later worker cannot dispatch an
action that lacks the committed reservation.

A resolution transaction similarly checks the issue revision, validates the
resolution evidence, updates the issue projection and invalidates dependent
results when necessary. It does not modify immutable candidate or source blobs.
The event history records what changed and why, while the current projection
supports efficient queries. Recovery can rebuild projections or verify them
against the event sequence.

Client idempotency is scoped by authenticated caller and operation. Repeating the
same key with the same canonical payload returns the existing result. Reusing the
key with a different payload returns conflict. The service must not accept a
client-supplied role in the payload as authority. Authentication determines the
caller and its permitted operations.

The local audit log is useful for reproducibility but is not automatically
 tamper-resistant against a database administrator. A bank deployment may require
external append-only storage, signed checkpoints or equivalent controls. Those
choices are part of the operational security design, not properties conferred by
calling a table an audit log. The prototype should preserve enough identity and
ordering information to integrate such controls later.

\section{An API with reviewable state transitions}

The proposed HTTP extension uses asynchronous jobs for investigation and explicit
resources for review. Creating a run requires the source packet, profile,
configuration and idempotency key. The server validates all references and returns
a run identifier and job reference. It does not block the request until every
model has completed.

\begin{longtable}{@{}P{49mm}P{41mm}P{58mm}@{}}
\caption{Proposed interpretation-service operations. Existing v0.1 endpoints remain separate.}\label{tab:interpretation-api}\\
\toprule
Operation & Purpose & Required behavior\\
\midrule\endfirsthead
\toprule Operation & Purpose & Required behavior\\\midrule\endhead
\texttt{POST}\newline\path|/interpretation-runs| & Create a bounded investigation & Validate source/profile/policy; authenticate caller; idempotent creation\\
\texttt{GET}\newline\path|/interpretation-runs/{id}| & Inspect current state & Return counters, phase, terminal reason and linked resources\\
\texttt{POST}\newline\path|/interpretation-runs/{id}/cancel| & Stop new automatic work & Preserve outstanding actions and produce a cancelled report\\
\texttt{GET}\newline\path|/interpretation-runs/{id}/report| & Retrieve a report version & Expose blocked and failed reports as well as review-ready ones\\
\texttt{GET}\newline\path|/interpretation-issues/{id}| & Inspect a discrepancy & Return candidates, evidence need, attempts and status separately\\
\texttt{POST}\newline\path|/interpretation-issues/{id}/decisions| & Record human adjudication & Require reviewer authority, expected revision and exact evidence links\\
\texttt{POST}\newline\path|/interpretation-runs/{id}/successors| & Authorize further work & Require reason, new policy and lineage; never reset the old run\\
\bottomrule
\end{longtable}

These paths are proposed additions, not descriptions of endpoints already present
in the retained OpenAPI file. The implementation task must generate a new OpenAPI
contract from actual request models and test it against the service. The book's
path table states responsibilities; the machine-readable schema files state the
record shapes.

Malformed requests return boundary errors without creating partial runs. Missing
references and incompatible source versions return semantic validation errors.
A stale expected revision returns conflict so the reviewer can reload the changed
issue. Unauthorized decisions fail without writing a review or audit entry that
suggests acceptance. The service records denied-operation security evidence
through its appropriate logging channel without pretending that the mutation
succeeded.

Pagination must preserve completeness. Collection responses include stable cursors
and total or explicitly unknown counts; reports bind the complete underlying
collection hashes. A user-interface page showing twenty issues cannot imply that
only twenty exist. Exports include all issue records or a manifest identifying
every page, with integrity checks that detect omission.

\section{The review interface should expose the decisive difference}

A reviewer should begin with the source passage and the decision under review,
not a wall of model transcripts. The interface shows the activity and subject
profile, the candidate's controlled-language rule, and the facts needed to apply
it. A neighboring comparison shows the strongest surviving alternative and a
scenario where the outcomes differ.

The discrepancy view explains the type of uncertainty. Missing factual evidence
should lead to a data request; unresolved meaning should lead to a legal question;
unsupported runtime operations should lead to engineering work. Mixing these
under a single confidence badge makes review slower and encourages the wrong
remedy. The interface can use concise status labels, with the full evidence one
step away.

Source references open the retained version and highlight the exact span. A live
link is useful for checking current authority, but it should not replace the
retained version used in the run. If the live document differs, the interface
shows the change and offers a successor investigation. This prevents a reviewer
from unknowingly approving a candidate against a document that changed after
its creation.

Two review moments require different records. An adjudication during the run
binds the issue revision, candidate and source context; it cannot require a hash
of a final report that does not yet exist. The eventual report includes that
adjudication. A subsequent final review binds the immutable report and candidate.
The companion \texttt{issue-decision} and \texttt{review-decision} schemas keep
these directions separate and avoid a circular report-to-decision hash dependency.
Neither record alone authorizes deployment.

A decision form asks for the proposition being accepted, the evidence relied on,
the treatment of the strongest objection, the scope of the decision and any
conditions. It does not ask only for a confidence slider. The authenticated role
and exact candidate/report hashes are supplied by the service. The reviewer
cannot approve an unseen revision through a stale browser tab.

Unresolved reports remain useful work products. They can present a focused
question for the next authorized reviewer, show which offers are affected and
identify a temporary operational policy proposal. The interface should make clear
whether that policy is a bank choice or a conclusion about the regulator's rule.
This distinction is central to keeping cautious operations from becoming false
statements of law.

\section{Provider adapters are deliberately narrow}

A provider adapter accepts a typed task and returns a bounded response plus usage
and execution metadata. It has no direct database credentials for approvals,
release records or budgets. It cannot invoke arbitrary tools selected in model
text. The controller supplies the exact retained sources and permitted retrieval
scope, and the adapter reports which inputs were actually sent.

The initial interface can support three adapter kinds: scripted fixtures,
retrieval tools and model proposals. Scripted fixtures are essential because they
make failures reproducible. Retrieval adapters return bytes and acquisition
metadata for validation. Model adapters return structured proposals and source
references that the controller checks. All adapters use the same action identity
and timeout accounting.

The adapter must distinguish refusal, abstention, malformed output, transport
failure, timeout and unknown remote result. Treating them all as an empty answer
would lose the reason a mandatory role is incomplete. It must also report hidden
provider behavior, such as automatic retries or server-side tool execution, when
that behavior affects resource accounting or evidence provenance. An integration
that cannot bound these behaviors is unsuitable for the unattended profile.

Prompt templates are versioned source files. They state the role, source
perimeter, required structured fields and prohibition on treating source text as
instructions to the system. The prompt should ask for assumptions and opposing
reasons, but those requests do not make the response reliable by themselves.
Validation and independent evidence remain necessary.

The service stores provider and model identifiers, relevant sampling settings,
template version and source hashes. A provider alias that changes its underlying
model must be recorded as such. Reproducibility may be limited even with these
fields, so reports distinguish replay of a retained response from regeneration of
a fresh model response. Deterministic execution of an accepted bundle does not
depend on reproducing the model's original generation.

\section{Security protects the meaning of the evidence}

The source packet can contain malicious or misleading instructions, whether
through a compromised document, an uploaded file or quoted text. The interpreter
must treat them as source content, not commands. A paragraph saying to ignore
prior rules cannot change the controller's budget or authentication policy. The
model's task is to analyze that text within the source context; tool permissions
come from the server.

Acquisition should retain the existing restrictions on HTTPS destinations,
redirects, public addresses, response size and time. A user-supplied URL does not
establish official authority. Uploaded documents remain unverified until a
separate process establishes provenance. A source that is useful for comparison
may still be ineligible to support release.

PDF parsing needs an isolation boundary in a bank deployment. The current local
prototype's bounded parser is not an operating-system sandbox for arbitrary
hostile files. A production intake worker should have limited filesystem and
network access, resource limits and a narrow output contract. This protects the
workbench without changing the legal interpretation method.

Bank factual data and public regulatory sources have different access needs.
The initial ensemble prototype should use public sources and synthetic scenarios.
Sending client or transaction data to a model provider requires the bank's
approved data-handling arrangement. The architecture supports minimizing that
exposure: legal interpretation can often use abstract fact patterns, while
production evaluation remains inside the bank with the generated Java package.

Logs must avoid becoming an uncontrolled copy of sensitive evidence. Store
references and necessary metadata in operational logs, retain full evidence in
access-controlled stores, and make redaction policies explicit. A redacted report
must preserve enough structure to explain its limitations. Redaction cannot
silently remove the only evidence supporting an approval while leaving the
approval summary apparently complete.

Security tests should use concrete attacks: source text that requests a tool
call, candidate output containing executable markup, a URL redirecting to an
internal address, an oversized graph, a forged reviewer role and an action result
that tries to change its run identifier. The expected behavior is rejection or
safe treatment as data, with a recorded reason. These tests protect the authority
boundaries on which the interpretation method relies.

\section{The ordered implementation sequence}

The implementation should proceed in small stages that each produce a reviewable
capability. The first stage freezes contracts and builds deterministic examples.
The next adds persistence and budgets. Only after the controller behaves correctly
with scripted members should it call live model providers. This sequence avoids
confusing stochastic model behavior with defects in state management.

\begin{longtable}{P{13mm}P{49mm}P{85mm}}
\caption{Ordered work packages for the proposed extension.}\label{tab:extension-tasks}\\
\toprule
ID & Deliverable & Acceptance condition\\\midrule\endfirsthead
\toprule ID & Deliverable & Acceptance condition\\\midrule\endhead
E01 & Contract models and reference policy & All valid examples pass; extra fields, incompatible limits and invalid transitions are rejected.\\
E02 & Source inventory and dependency records & Every retained unit has a disposition; missing or mismatched units create issues even under unanimous proposals.\\
E03 & Immutable candidate and assumption store & Revisions retain parents and source hashes; changing a definition invalidates dependent candidates.\\
E04 & Issue identity and lifecycle & Child issues inherit root accounting; terminal unresolved states cannot masquerade as resolutions.\\
E05 & Atomic budget and outbox service & Concurrent reservations never exceed caps; crash recovery preserves issued counts and fencing.\\
E06 & Scripted provider and controller & Omission repair, persistent ambiguity, timeout, no-progress and cancellation traces match the specification.\\
E07 & Structured comparison and witnesses & Different source contexts are detected; supported formal differences replay in Python and Java.\\
E08 & Complete report and review interface & Every issue and mandatory role is represented; dissent and unexplored families remain visible.\\
E09 & Review and release integration & Exact report/candidate/build identity is required; all stale or unresolved release attempts are rejected.\\
E10 & Restricted retrieval and model adapters & No uncounted retries or model-directed tools; malformed and unavailable responses retain distinct outcomes.\\
E11 & Independent reference corpus & Source and scenario judgments are frozen before model outputs; disagreement and adjudication are retained.\\
E12 & Equal-budget evaluation harness & Baselines share source access and cost accounting; uncertainty and abstention metrics are reported together.\\
E13 & Optional argumentation/search extension & Small reference graphs and family-preservation tests pass; incomplete search remains explicit.\\
E14 & Bank integration and shadow operation & Real identity, host mappings, operational dispositions and rollback procedures are approved and exercised.\\
\bottomrule
\end{longtable}

Dependencies matter. E05 depends on the record and issue contracts, while E06
needs E05 and a scripted adapter. E09 depends on complete reports and existing
release identity checks. E10 can be developed behind the adapter interface, but
live evaluation should wait for E11 and E12. E13 is optional until the simpler
scheduler's limitations are measured; it must not delay a clear implementation
of bounded disagreement handling.

Each task should specify files to add or extend, existing interfaces to reuse,
acceptance identifiers and evidence locations. The companion task JSON provides
that structure without changing the existing master plan's completed T00--T22
history. These are new extension tasks, not a retroactive claim that the earlier
MVP failed to implement its agreed scope.

An implementation agent should stop a task for a genuine contract inconsistency,
missing authority or unavailable required infrastructure. It should not stop
merely because a scripted candidate fails: detecting that failure is often the
acceptance criterion. A failed interpretation candidate can trigger the next
planned repair phase while still blocking release. The task record distinguishes
those outcomes.

\section{Acceptance scenarios that exercise the whole cycle}

The first end-to-end fixture uses the explicit type-of-product omission. It
freezes the source packet, emits two candidates, creates a discrepancy, reserves
a repair action, validates the revised formula, reruns affected checks and
produces a report. The report is ready for review only if all other required
coverage and dependencies in that fixture are complete. The real second-circular
packet remains blocked where its dependencies are unresolved.

The second fixture uses the cashback ambiguity with no resolving authority. It
must preserve both supported candidates through the configured rounds and end in
a blocked report. The test checks not only the terminal status but also the
attempt history, retained objections, question for the reviewer and absence of
release authority. This scenario prevents an implementation from treating
exhaustion as permission to choose the highest score.

The third fixture tests unanimous omission. All candidate providers return the
same incomplete rule, while an independent inventory identifies an uncovered
source unit. The controller must open a material issue despite unanimous agreement.
If it does not, the ensemble has implemented voting rather than the proposed
method.

The fourth fixture tests failures under concurrency. Two workers attempt the last
available reservation, one crashes after dispatch, and a stale worker later
returns a result. Exactly one reservation may consume the remaining slot; the
crashed action remains counted; the stale result cannot overwrite a newer state.
The test uses a controlled clock and scripted transport so it is reproducible.
Real distributed-system testing later extends this fixture under the deployment
infrastructure.

The fifth fixture tests invalidation. A reviewer approves a candidate and report,
then a source derivative reveals a previously omitted footnote. The affected
coverage and interpretation records become stale. Release with the earlier
approval must fail until the impact is reviewed and the required checks rerun.
An unchanged Boolean body does not bypass the source-coverage question.

The sixth fixture tests report recovery. A run is cancelled or encounters a
malformed provider response. The service must still produce a complete terminal
report from durable records. If rich rendering fails, a minimal machine-readable
report must remain available. The acceptance condition is not a pretty success
screen; it is a reconstructible account of what happened.

\section{The first vertical slice and its completion rule}

A useful first vertical slice implements E01--E09 with scripted members and the
public second-circular packet. It exercises source inventory, competing candidates,
a repair, a persistent ambiguity, a report and the existing Java path. It does
not need an expensive model integration to test whether the controller's logic
is coherent. This is the smallest product-shaped demonstration of the proposed
method.

Completion requires more than green unit tests. The service must generate an
inspectable report from a fresh run, preserve the exact input and output identities,
show all unresolved issues in the interface, and reject a release attempt for the
blocked example. A separate synthetic fully resolved fixture can demonstrate the
positive review path using local test identities. Its approval status must remain
clearly synthetic.

The implementation note should record the command, environment, source hashes,
policy, result files and any departures from the design. If a departure changes
meaning, budget accounting or approval behavior, it requires a contract revision
and new acceptance cases. Minor internal choices can be made by the agent, but
must not silently alter observable semantics.

Live models then replace scripted proposal adapters without changing the controller
contract. That transition creates a new empirical question: whether the methods
find useful alternatives and omissions on unseen circulars at acceptable cost.
The next chapter supplies the reference process, baselines and uncertainty
analysis needed to answer that question. Engineering completion of the vertical
slice is a prerequisite for that study, not its result.

\chapter{Measuring error, uncertainty and the value of redundancy}
\label{ch:evaluation}

\section{The experiment must ask the right question}

The proposed ensemble has not yet demonstrated improved interpretation of SFC
circulars. Its design is motivated by inspected literature, the second-circular
exercise and explicit failure analysis. Those are reasons to build and test it,
not evidence that its combination of methods is superior. A credible evaluation
must distinguish a working controller from a useful interpretation method and
from a release process acceptable to the bank.

The main empirical question is whether the method discovers and reports material
interpretation errors and omissions more reliably than simpler alternatives at
comparable cost. A secondary question is whether its reports help reviewers make
better, faster and more traceable decisions. A third concerns resource use:
whether targeted resolution spends effort on questions that can actually be
answered. These questions require different measurements.

The feared failure is a false clean result: a report appears complete and
unambiguous while a material source requirement, alternative reading or factual
assumption has been missed. A system that abstains on every case avoids some false
clean results but may be operationally useless. Evaluation must therefore measure
risk and useful coverage together. Neither raw agreement nor a low average model
uncertainty score is an adequate promotion criterion.

Before a live study, write an evidence contract naming the corpus, baselines,
reference process, primary outcome, uncertainty analysis, hard vetoes, resource
limits and permitted conclusions. Freeze it before examining comparative results.
The manuscript proposes such a protocol; it does not report a completed trial.
Thresholds that depend on the bank's risk tolerance remain decisions for the
responsible owners, not numbers to be invented by the implementation agent.

\section{Define the unit of evaluation}

A circular contains many questions, and a single question can affect many
transactions. Counting every transaction as an independent interpretation example
would exaggerate sample size when they all depend on the same disputed
classification. The evaluation should distinguish document, provision,
interpretation question, candidate, scenario and transaction levels.

The primary unit for omission detection is a reviewed provision or interpretation
question within a source packet. The source packet includes the circular and
required dependencies as of a declared cut-off. Scenario-level outcomes measure
how candidate rules behave under facts. Document-level summaries measure whether
an apparently complete investigation misses any material issue. These views
complement one another but have different denominators.

For example, the 729 ternary assignments are not 729 independent regulatory
interpretations. They all arise from one selected formula and six abstract inputs.
They are excellent finite-domain engineering coverage for that formula. Treating
them as a sample of 729 circulars would produce meaningless uncertainty intervals
for legal translation accuracy. The 22 named cases likewise share source,
author and modeling assumptions.

The corpus should record these relationships explicitly. Each scenario names its
source packet, provision, interpretation question and reference decision. Each
candidate names its generation method and source version. Analysis can then
aggregate at the appropriate level and use cluster-aware uncertainty estimates
rather than treating related rows as independent observations.

An unresolved reference question is still an evaluation object. The desired
output may be recognition of ambiguity, preservation of supported alternatives
and escalation, rather than selection of a single label. If the dataset forces
such questions into a binary answer, it rewards unjustified certainty and
penalizes the behavior the bank actually needs.

\section{Build the reference before seeing the model's answer}

The reference process begins with independent source inventory. Reviewers enumerate
provisions, footnotes, definitions, exceptions, dates and dependencies. A separate
reader checks the inventory against the retained document. The process records
which materials were unavailable and which judgments depend on them. The corpus
should not silently exclude difficult documents merely because their references
are inconvenient to acquire.

For each selected question, at least two appropriately qualified reviewers form
initial judgments without seeing the candidate system's answer. They state the
reading, assumptions, relevant passages, plausible alternatives and examples that
would distinguish them. This is more work than assigning a label, but it creates
a reference that can test interpretation rather than superficial answer matching.
The exact number and qualifications of reviewers are part of the study plan.

When reviewers disagree, adjudication examines the reasons. A third reviewer or
panel can resolve the issue, retain multiple acceptable readings, or declare the
available authority insufficient. The adjudication record preserves the original
positions and the evidence that changed them. Majority voting among reviewers
should not automatically settle a question of authority or source meaning.

Reviewer independence is documented, not assumed. Shared training, a common case
summary, prior involvement in the implementation and access to model outputs can
all influence judgments. Some collaboration is necessary for adjudication, but
initial blinded judgments should remain distinguishable from the later consensus.
The study reports both stages.

Reference revisions are inevitable when new evidence appears. They require a
versioned correction with a reason. Comparative results are recomputed against
the corrected reference, and the report explains the effect. A correction that
favors one system is not automatically improper, but it must be justified without
using that system's desired score as the reason.

The reference itself is not a universal truth oracle. It is an independently
reviewed account of the declared source perimeter at a time. The study can measure
agreement with that account, sensitivity to its material issues and handling of
its unresolved questions. It cannot establish that no future authority or
previously unknown source could change the interpretation.

\section{Construct a corpus that exposes the expected failures}

A useful corpus contains several kinds of difficulty: explicit Boolean conditions,
exceptions, definitions supplied elsewhere, actor and activity scope, uncertain
classification, amendments, transition periods, event histories, monetary units
and interactions between requirements. The initial public examples can anchor the
engineering harness, but a method comparison needs additional circular families
that were not used to design the prompts or repair logic.

Selection should reflect the intended product perimeter. If the product focuses
on private-bank distribution controls, an evaluation consisting entirely of simple
product-authorization filing instructions would be a weak test of deployment
readiness. Conversely, including arbitrary legal texts outside the supported
language can be useful for measuring abstention, but should not be mixed with
in-scope accuracy without reporting the distinction.

Temporal separation matters. A holdout containing a later amendment can reveal
whether the system improperly reuses an earlier definition or priority. The
source packet must state whether the task is historical interpretation as of a
cut-off or current applicability review. A model with knowledge of a later rule
can otherwise appear to correct a historical answer by importing information
that was unavailable at the assessment date.

Synthetic cases are valuable for controlled defects. They can isolate a missing
exception, an ambiguous attachment or a contradictory fact. Their provenance must
be visible, and their performance must be reported separately from naturally
occurring circular questions. A system can learn the style of seeded mutations
without becoming better at discovering unanticipated omissions.

Include negative controls in which no material ambiguity is supported, as well
as positive controls with known alternative readings. A method that always
invents a dispute should not receive full credit for finding ambiguity. The
reference should distinguish a plausible source-supported alternative from an
arbitrary logical variation. Review burden from unsupported alternatives is an
operational cost worth measuring.

The corpus should also include deliberately incomplete packets. The correct
behavior may be to identify the missing dependency and block a conclusion. These
examples test whether the system recognizes the limits of its evidence. They
must not be scored as ordinary translation errors simply because no final rule
is produced.

\section{Use a baseline ladder with comparable resources}

The simplest baseline is a single structured reading with the same source packet
and a fixed prompt. It establishes what the ensemble adds beyond ordinary
translation. A stronger baseline uses that single reading plus deterministic
schema, coverage and execution checks. This prevents the proposed method from
claiming credit for checks that do not require an ensemble.

A third baseline uses repeated sampling and answer aggregation, inspired by
self-consistency. It receives the same total resource allowance and reports its
aggregation method. A fourth uses diverse initial proposals with a fixed,
deterministic resolution scheduler. The enhanced method adds the proposed
hypothesis-family coverage, targeted questions and optional argumentation or tree
search. This ladder isolates the contribution of each mechanism.

\begin{table}[htbp]
\centering\small
\caption{Comparison arms separate redundancy from search and verification.}
\begin{tabularx}{\textwidth}{P{16mm}YY}
\toprule
Arm & Mechanism & Question it answers\\
\midrule
B0 & One structured proposal & What does ordinary structured translation achieve?\\
B1 & B0 plus deterministic checks & How much comes from source coverage and validation alone?\\
B2 & Repeated proposals with aggregation & Does repeated sampling add useful error detection at equal cost?\\
B3 & Diverse proposals and fixed resolution & Does explicit disagreement handling improve useful coverage?\\
B4 & B3 plus targeted search or argumentation & Does the enhanced scheduler find additional material issues?\\
\bottomrule
\end{tabularx}
\end{table}

Equal resources do not mean merely equal numbers of model calls. Calls differ in
input length, output length, retrieval, latency and cost. The study should record
both token or monetary budget and elapsed time, and present comparisons at
several predeclared budget levels if that is part of the question. Human reviewer
access must also be comparable. Giving only the proposed method a legal reviewer
would confound the method with extra expertise.

Prompt and configuration tuning use a development set. The strongest baseline
should receive reasonable task-specific tuning rather than a deliberately weak
prompt. The test set remains held out until the protocol is fixed. If a test
failure triggers a repair, the repaired method is evaluated on a fresh holdout or
reported as development evidence, not as an untouched final result.

Search methods require their own assumptions audit. A beam width, family quota,
UCT exploration constant or no-progress threshold transferred from puzzles or
another repository is a hypothesis. It can be a starting point, but must not
become the default without target-specific evidence. The deterministic scheduler
is the initial engineering baseline because its behavior is easier to inspect,
not because it has been proven optimal.

\section{Measure omissions, alternatives and useful abstention}

Let $Q$ be the set of independently reviewed material questions in the evaluation
corpus. For each $q\in Q$, the reference records required conditions, accepted or
plausible readings and whether available evidence resolves the question. A
material omission occurs when the system's report fails to represent a required
condition or unresolved question within its claimed scope. This definition
includes missing uncertainty, not only incorrect final formulas.

A false clean result occurs when the system reports a question as resolved and
eligible for review without disclosing a material unresolved defect identified by
the reference. The denominator is the number of questions or documents declared
clean, depending on the reported metric. A system that declares fewer clean cases
may reduce this rate while increasing manual workload. The report must show both.

Define useful coverage as the fraction of eligible questions for which the system
produces a reference-supported, sufficiently complete result without unresolved
material defects. Define selective risk as the error rate among the questions it
chooses to resolve. A risk-coverage curve varies the decision policy and shows the
tradeoff. The policy must not be tuned on the final test labels after the fact.

Alternative-reading recall measures whether the report contains the reference's
material supported alternatives. It should be assessed by meaning and consequence,
not string matching. Alternative precision measures how many proposed alternatives
are judged relevant and supported. These measures expose a method that increases
recall by overwhelming reviewers with arbitrary possibilities.

Resolution quality has a different denominator: detected discrepancies. Measure
whether the selected action obtains the needed evidence, produces a valid repair,
correctly escalates, or falsely closes the issue. A successful HTTP response is
not a successful resolution. The action's declared question and expected evidence
type make this classification possible.

Report completeness is largely an engineering metric. It checks that every issue,
role, candidate disposition and resource outcome appears. It can be tested
exhaustively over bounded controller fixtures. It does not establish that the
issues themselves cover every material legal question. Keeping these metrics
separate prevents a perfect report-rendering score from being presented as
perfect interpretation accuracy.

\section{Severity and direction of error need explicit treatment}

Not every error has the same consequence. Missing a prohibition may permit an
impermissible offer; inventing a prohibition may block a legitimate activity or
mislead the bank about its obligations. Both are errors relative to the accepted
interpretation, even if one appears operationally conservative. The evaluation
should report their directions rather than collapsing them into a single accuracy
number.

Severity categories must be defined before scoring. They can consider whether the
error changes an action, affects many clients, is hard to detect downstream or
undermines the source and approval chain. Assigning weights is a policy choice
requiring the bank's input. The study can report unweighted counts and separate
severity strata until those weights are approved.

A weighted average can hide a rare severe error among many trivial successes.
Therefore severe false clean outcomes should also be reported as individual cases
with their mechanisms. A predeclared severe-error veto can block promotion even
when average performance is favorable. That veto should not automatically stop
all research: if the failure reveals a repairable omission detector, the next
planned repair phase may remain justified.

The report should distinguish an invalid experiment from a failed candidate.
Leakage of test labels into prompts invalidates the comparison and may require
rebuilding the evaluation. A candidate missing an exception on a valid heldout
case is evidence against that candidate's readiness. An unavailable provider is
an operational failure that may prevent estimating interpretation performance.
These outcomes should not all be summarized as ``the experiment failed.''

Reviewer workload also has severity. A report containing hundreds of unsupported
alternatives may conceal the important issue through overload. Measure review
time, number of evidence inspections, unresolved-question quality and corrections
made by reviewers. A method that catches more errors but makes the review process
unmanageable needs a different presentation or search policy, not an unqualified
claim of superiority.

\section{Zero observed errors does not mean zero risk}

Suppose a study observes no errors in $n$ independent, identically distributed
trials with unknown error probability $p$. Under this simplified binomial model,
the probability of seeing zero errors is $(1-p)^n$. A one-sided upper confidence
bound at level $1-\alpha$ solves $(1-p)^n=\alpha$, giving
\begin{equation}
 p_{\rm upper}=1-\alpha^{1/n}.
 \label{eq:zero-errors}
\end{equation}
For a 95 percent bound and $n=100$, this is approximately $2.95$ percent. With
$n=1000$, it is approximately $0.30$ percent. These calculations show why a small
perfect test set cannot establish a very small error rate.

The independence assumption is substantial. One hundred scenarios derived from
one circular do not generally behave like one hundred independent circulars.
Shared wording, definitions and prompts create clusters. The calculation is
therefore illustrative unless the trial definition and sampling design justify
the model. It should not be applied to the 729 truth-table assignments as a legal
error bound.

For method comparisons, paired evaluation on the same source questions reduces
irrelevant variation. A question-level difference records whether one method
finds a material issue that the other misses. Confidence intervals or permutation
procedures should preserve the document or question clustering appropriate to the
sampling design. Resampling individual rows while ignoring their shared circular
can produce overly narrow intervals.

Multiple stochastic generations require multiple seeds or repeated provider runs,
with the seed and configuration recorded where supported. Differences from one
run are descriptive. If a model API does not guarantee reproducibility from a
seed, the study records independent request repetitions and their limitations.
The retained responses permit audit even when exact regeneration is impossible.

Tail measures such as maximum review time or the 99th percentile of action cost
need particularly cautious interpretation. A few difficult circulars can dominate
them, and small samples give unstable estimates. They are useful diagnostics and
capacity-planning signals, but a superiority claim requires a predeclared analysis
with sufficient observations and uncertainty evidence.

The evaluation report should state which candidates failed hard safety or
integrity screens, which remain viable, whether any ranking is statistically
supported, and which differences are descriptive only. Passing a hard screen is
not evidence that one viable method is better than another. This discipline is
especially important when the apparent gain comes from a small number of rare
severe errors.

\section{Calibration asks a narrower question than truth}

A calibrated score has a defined relationship to an observed event in a target
population. For example, among cases assigned a particular probability range,
the fraction agreeing with the independently reviewed reference should be
consistent with that range, within sampling uncertainty. Calibration is not a
property conferred by using the word confidence in a JSON field.

The target event must be precise. ``The interpretation is correct'' is too broad
unless the reference process, source perimeter and acceptable alternatives are
specified. A more measurable event might be ``the proposed paragraph-level rule
has no material discrepancy against the adjudicated reference for this corpus.''
Even that event depends on the quality and scope of the reference.

Calibration data must be separate from the data used to tune the score and from
the final evaluation set. A score trained on one circular family may be
miscalibrated on another, especially when definitions, authority structure or
source quality change. Model and prompt updates can also invalidate an earlier
calibration. The calibration record therefore binds the method and population
versions.

Reliability diagrams and proper scoring rules can diagnose probability estimates,
but a visually plausible diagram with few cases is weak evidence. Calibration
also does not ensure discrimination: a model that assigns the base rate to every
case can be calibrated yet unhelpful for selecting review priorities. The study
should assess both risk ranking and calibration, with uncertainty intervals.

Semantic entropy can nominate uncertain answers, as discussed in
Chapter~\ref{ch:search}. Its published evaluations use particular datasets and
correctness proxies \citep{semantic2023}. A high area under a ranking curve in
those tasks does not establish calibrated legal error probabilities. In the
workbench, entropy may help prioritize review after a target-specific study, while
mandatory source and release checks remain categorical requirements.

Low entropy is particularly dangerous under common-mode omission. Every sampled
answer can concentrate on the same incomplete reading. The independent inventory
and source review remain necessary even if an uncertainty score appears excellent
on average. A useful evaluation explicitly includes these unanimous-error cases
rather than assuming that disagreement accompanies every mistake.

\section{Conformal prediction can support sets under stated assumptions}

Conformal prediction offers a principled way to construct prediction sets from
calibration examples \citep[section 1.1 and appendix D]{conformal2022}. It is
relevant because the proposed product should sometimes return several candidates
rather than a forced answer. Its guarantee, however, concerns a specified
statistical prediction problem under exchangeability. It is not a theorem about
all possible interpretations of English.

To make the mechanism concrete, suppose each task has a finite reviewed label
space $\mathcal Y$, such as a small set of predetermined interpretation classes.
A fitted scoring function $s(x,y)$ assigns larger values when label $y$ looks less
compatible with task $x$. The scoring function is fixed before calibration.
For $n$ calibration tasks with reviewed labels $(X_i,Y_i)$, compute
$S_i=s(X_i,Y_i)$.

Choose a target miscoverage level $\alpha$ and define
\begin{equation}
 k=\left\lceil(n+1)(1-\alpha)\right\rceil,
 \qquad
 \widehat q=\begin{cases}
 S_{(k)},&k\leq n,\\
 +\infty,&k=n+1,
 \end{cases}
 \label{eq:conformal-quantile}
\end{equation}
where $S_{(k)}$ is the $k$th smallest calibration score. For a new task $x$, return
\begin{equation}
 \Gamma(x)=\{y\in\mathcal Y:s(x,y)\leq\widehat q\}.
 \label{eq:conformal-set}
\end{equation}
The finite-sample correction uses $n+1$, not simply the empirical $1-\alpha$
quantile with an unspecified interpolation rule. The inclusive inequality also
matters when scores tie.

Why does this work? Conditional on the separately fitted scoring rule, if the
$n$ calibration examples and the new example are exchangeable, the position of
the new true-label score among the $n+1$ scores is symmetric. In the no-tie case,
its rank is uniform. Excluding it requires a rank beyond $k$, whose probability
is at most $\alpha$ by the definition of $k$. With ties and inclusive sets, the
basic coverage statement is conservative. This gives marginal coverage over a
new exchangeable example and the calibration sample.

For $n=19$ and $\alpha=0.1$, $k=18$, so the threshold is the eighteenth ordered
calibration score. For $n=4$ and $\alpha=0.1$, $k=5$, requiring the infinite
threshold and potentially the entire label space. The latter example makes a
practical point: a very small calibration set cannot support a narrow set with a
strong finite-sample coverage target through this construction.

The guarantee does not say that each particular circular has a 90 percent chance
of containing the right reading in its returned set. Nor does it guarantee
coverage conditional on every rare ambiguity type. Marginal coverage can coexist
with poor performance on a small important subgroup. The distinction between
marginal and conditional coverage is discussed explicitly in the tutorial
\citep[section 3]{conformal2022} and is central to a high-stakes application.

The label space is another limitation. If the generator has never proposed a
material reading and the label space contains only generated candidates, a set
selection method cannot include that missing reading. One can add an
``unrepresented interpretation'' or escalation outcome to a reviewed task
formulation, but that changes the prediction problem and requires reference data.
Conformal calibration does not establish semantic completeness of open-ended
candidate generation.

The appropriate early experiment is therefore narrow: use a fixed, reviewed
classification task or a defined error-detection score, calibrate on separate
source families where exchangeability is defensible, and measure set size,
coverage and subgroup behavior. Source amendments, provider changes and shifts in
document type should trigger renewed assessment. The inspected official tutorial
code is useful for checking quantile and set construction, but its image examples
are not a ready-made legal calibration protocol.

\section{Search coverage and semantic completeness are different}

A controller can report that it explored every node in a finite generated graph.
That is a meaningful computational fact. It does not establish that the graph
contains every supported legal interpretation. The generator, inventory and
construction operators determine what enters the graph. Evaluation must examine
both the search procedure and the generation boundary.

For each reference alternative, record whether it was generated, retained,
investigated, correctly dispositioned and represented in the final report. This
stage-by-stage accounting locates the failure. If an alternative was never
generated, changing the scheduler may not help. If it was generated but pruned by
a global beam before investigation, family preservation may be useful. If it was
investigated but falsely defeated by a weak argument, the evaluation criterion
needs repair.

The same accounting applies to omissions. A source unit can be missed at
acquisition, extraction, inventory, normative classification or rule mapping.
Treating all such failures as model reasoning errors hides the actionable cause.
The ensemble's modular design makes it possible to test which redundant activity
actually contributed detection.

Ablations remove one component at a time under comparable budgets: the independent
inventory, the alternative reader, the argument critic, the witness generator or
the no-progress scheduler. If removing a component changes both resources and
information access, the interpretation of the ablation is limited. The study
should state whether it measures the component's total contribution or its value
at fixed cost.

Search completeness claims must name the finite object. ``All twenty-four retained
candidates were compared on the declared Boolean domain'' is testable. ``All legal
readings were considered'' is unsupported. The report should also expose the
unexplored frontier, discarded unsupported proposals and resource-limited
comparisons so that reviewers can judge whether more investigation is justified.

\section{Evaluate the reviewer, not only the model}

The final product is a workbench used by people. A study should therefore compare
review with and without the ensemble dossier, using a design that controls for
case difficulty and learning effects. Reviewers should not see the same case in
both conditions in an order that makes the second condition trivially easier.
A counterbalanced or appropriately randomized assignment can reduce this problem.

Outcomes include detection of material defects, quality of the final rationale,
time to a defensible decision, number of unnecessary escalations and ability to
identify remaining uncertainty. Faster review is valuable only if the substantive
quality is preserved. A dashboard that encourages reviewers to accept a misleading
clean badge can reduce time while increasing risk.

The study should inspect automation bias directly. Include cases where the
leading model answer is wrong, where a minority candidate is right relative to
the reference, and where no single answer is justified. Observe whether reviewers
inspect the opposing argument and source. The purpose is not to trick staff, but
to learn whether the interface supports the intended skeptical review.

Reviewer feedback also tests self-containment. Can a compliance reader understand
the formal rule through its controlled-language rendering and examples? Can an
engineer identify the required Java behavior without consulting an external paper?
Can both explain why a solver result does not resolve an uncertain definition?
These comprehension tasks are more informative than asking whether the report
looks professional.

Any study involving actual bank staff, confidential materials or production
workflows requires the bank's normal authorization and data controls. The public
prototype can first use synthetic tasks and qualified independent readers.
Its findings should be reported at that level. A successful public-source
usability exercise is not evidence that every private-bank process has been
validated.

\section{Promotion decisions and the next evidence needed}

The study should predeclare a small set of promotion criteria and hard vetoes.
A severe false clean result, broken source identity, uncounted external action,
unauthorized release or incomplete mandatory report can veto promotion. Average
runtime, model score and cost are explanatory until the correctness and integrity
conditions are met. A method cannot trade away a required release invariant for
better throughput.

Among methods that pass the hard screens, ranking requires the predeclared
uncertainty analysis. If the data do not distinguish them, report them as viable
under current evidence. A descriptively favorable method may justify a longer
study or an optional feature, but not a new default merely because its mean score
is highest in a small sample.

\begin{table}[htbp]
\centering\small
\caption{Decision record required after a substantive method comparison.}
\begin{tabularx}{\textwidth}{P{39mm}Y}
\toprule
Decision field & Required interpretation\\\midrule
Primary criterion & State the observed result and uncertainty against the frozen criterion.\\
Veto status & Identify severe errors, integrity failures and invalid comparisons before discussing speed.\\
Viable candidates & List methods that remain eligible for further evaluation; do not imply an unsupported ranking.\\
Main uncertainty & Distinguish reference ambiguity, sampling error, distribution shift and untested implementation.\\
Next justified action & Promote a bounded use, extend the study, repair a mechanism, or rebuild an invalid harness.\\
Limit of conclusion & State the source perimeter, task population and deployment claims not established.\\
\bottomrule
\end{tabularx}
\end{table}

A failed candidate should trigger the repair that its failure supports. Missing
source units suggest improving inventory and extraction diversity. Losing a
minority interpretation suggests changing family retention. Repeated ineffective
actions suggest a better question-selection policy. A contaminated reference
requires repairing the evaluation before comparing methods again. These are
specific next steps, not reasons to declare the entire research direction either
successful or impossible.

Every serious run should retain a manifest with source and code identities,
commands, environment, model/provider versions, configurations, seeds or request
repetitions, wall time, resource use, plan and result paths. The result note should
include the strongest alternative explanation and the evidence that would overturn
its conclusion. This makes the evaluation reproducible enough to guide product
decisions rather than merely produce an attractive benchmark table.

The immediate conclusion remains limited: the proposed method is ready to be
implemented as a bounded, testable extension, and its empirical legal-interpretation
performance remains to be established. The next chapter describes how the bank
can preserve that distinction through shadow operation, source changes and
ongoing review.

\chapter{Operating the product without losing its evidence}
\label{ch:operation}

\section{The bank remains responsible for the control}

The workbench can organize evidence, expose alternatives, verify consequences and
prevent certain unauthorized transitions. It cannot assume the bank's legal or
managerial responsibility. The operational design must identify who owns source
applicability, interpretation, factual mappings, software correctness and release.
A system that assigns every difficult question to a generic ``human review'' queue
has not completed that design.

Meaning review belongs to people authorized to assess the relevant regulatory
requirement. Data owners establish what evidence supports facts, how long it
remains valid and how corrections are represented. Engineering owners establish
that the software preserves the approved specification. Release owners verify the
package, approvals and operational readiness. Business owners determine the
workflow response to unknown or blocked results within the bank's approved policy.

These responsibilities can overlap institutionally, but conflicts and separation
requirements must be explicit. An author should not be able to approve its own
interpretation through a second display name. A model-generated reviewer identity
has no authority. The local prototype's test roles exercise the mechanics; the
bank deployment must bind them to enterprise identity and actual delegated
responsibilities.

The product's value is not that it eliminates judgment. It makes the judgment
specific and attributable. Instead of asking a reviewer to bless a long generated
summary, it asks which definition applies, why an exception attaches to a given
condition, or whether a source is current for the activity. The final control
then carries those decisions into its executable form and future amendment review.

\section{Start with a narrow operational perimeter}

A first deployment should name the activity, legal entity, products, client or
counterparty categories, source families and decision points it supports. This
perimeter is not a marketing description. It determines which circulars enter the
workbench, which factual schemas are required and which host workflows can consume
the resulting package.

For example, a paragraph-10 incentive-control pilot could focus on reviewing a
single component of a proposed offer before it is launched. It would need a
reliable definition of the component, evidence of its product relationship and a
reviewed classification process for gift and fee discount. It would not by itself
cover every marketing, disclosure or distribution requirement. The host must
continue to route the offer through other applicable controls.

Narrowing the perimeter is useful when it removes an unsupported ambiguity by an
explicit operational restriction. If the bank excludes mixed offers from the
automated pilot and routes them to manual review, the system can assess simpler
cases under a clearer subject definition. That restriction must be enforced by
input validation and workflow routing, not merely stated in a project memo.
Otherwise mixed offers will eventually arrive and receive an unintended answer.

The perimeter also defines negative tests. Inputs from another actor, product type
or jurisdiction should produce the specified out-of-scope or unsupported-profile
result. They should not be coerced into the closest known category. A controlled
failure at the boundary is preferable to a plausible answer about the wrong
activity.

Expansion requires evidence. Adding a new circular family may introduce temporal
obligations, different authority relationships or vocabulary that the original
schema does not support. The team should create a new source and method review,
update the factual mappings, and evaluate the expanded task. Success on a narrow
pilot is a reason to consider expansion, not proof that transfer is automatic.

\section{Keep regulatory meaning distinct from bank policy}

A bank can choose a more restrictive operational policy than the minimum condition
identified in a circular. It may decline an uncertain promotion, require an
additional approval or retain a manual check. These decisions should be represented
as bank policy with their own owner, rationale and effective interval. They should
not be folded into the regulatory rule and later cited as if the SFC required them.

This separation improves both accountability and maintenance. If a regulatory
interpretation changes, the team can see whether the bank policy still makes sense.
If a business owner changes an internal approval threshold, it does not appear
that the regulator amended the circular. Reports can explain a blocked offer as
blocked by bank policy pending classification, rather than falsely declaring it
prohibited by the source.

The runtime can evaluate both kinds of controls, provided the bundles identify
their authority and purpose. The host's orchestration specifies how their results
combine. A regulatory condition returning false may still be blocked by a bank
policy; a bank policy returning permissive cannot override a regulatory
prohibition. The combination rules need their own reviewed semantics and tests.

An unresolved interpretation should not be hidden in a conservative default.
Suppose the team treats every cashback as a gift until clarification. That can be
a cautious temporary policy, but the uncertainty report should still state that
the legal classification remains unresolved. Otherwise future reviewers may
mistake a risk-management choice for settled authority and propagate it into
other products.

Temporary policies need expiry or review triggers. An indefinite ``pending
clarification'' rule can become permanent through neglect. The workbench should
record the question, owner, next review condition and affected activities. A
reminder is an operational aid; it does not automatically extend the policy or
resolve the underlying source issue.

\section{Shadow operation is a study, not a hidden release}

Shadow operation runs the proposed control alongside the existing process without
letting its result directly determine the transaction. It allows the team to
observe real input quality, classification workload, latency and disagreement.
The scope and access to real data require the bank's normal authorization.
Shadow results must be clearly distinguished from approved production decisions.

A useful shadow study records the existing process's decision, the workbench's
result, the facts available to each and the timing of those facts. When they
disagree, an independent review examines the reasons. The existing process is not
automatically the truth oracle; manual procedures can be wrong. Equally, the
workbench is not correct merely because it produces a more detailed explanation.
The study needs an adjudication process consistent with Chapter~\ref{ch:evaluation}.

The team should predeclare what would justify promotion, what would veto it and
what would trigger a repair. A severe false clean result might block promotion
while motivating a targeted fix. A data-mapping defect might invalidate a subset
of observations until corrected. A provider outage may reveal operational
fragility without answering the interpretation question. These distinctions keep
the study from becoming a succession of improvised demonstrations.

Shadow operation also reveals whether the selected unit of assessment matches
business practice. If staff enter a whole campaign while the rule expects one
benefit component, the system may receive apparently valid but semantically wrong
inputs. The remedy may be a workflow redesign, not a more capable model. The
interface should help staff identify and evidence components consistently.

Reports from shadow operation remain subject to access, retention and correction
rules. They can contain sensitive facts and preliminary judgments. A later
adjudication should append a correction rather than rewrite the original
observation. This preserves the study's evidence and prevents an apparently
perfect retrospective record from being created by editing away disagreements.

\section{Source monitoring begins a new interpretation cycle}

The workbench should monitor the source families in its approved perimeter for
new publications, revisions and incorporated-material changes. Discovery creates
a candidate source record; it does not automatically establish current
applicability or overwrite the active rule. The acquisition process retains bytes,
metadata, retrieval time and the relationship to prior versions.

A textual difference is only the first signal. A changed definition may affect
many controls even when their operative paragraphs are unchanged. A formatting
change may have no substantive consequence but still change the file hash. The
impact service therefore combines byte identity, extracted text differences,
source-unit mapping and the explicit dependency graph.

The first impact report identifies which candidates, facts, rules, duties,
reports and releases depend on the changed material. It should distinguish direct
citations from indirect definition dependencies. A control can be affected through
an incorporated code provision that is not quoted in its main explanation. This
is why dependency closure was required before initial interpretation.

An unchanged structural dependency report is not a proof of semantic equivalence.
Conversely, a changed paragraph number does not prove a changed obligation.
Semantic comparison and legal review are separate actions. The ensemble can
propose interpretations of the change and construct scenarios on which old and
new rules differ, but a reviewer must establish the intended effective regime.

Every changed rule version requires the relevant review, build and verification
steps again. Reusing a previously approved JAR because its human-readable title
looks the same would break the evidence chain. Reusing a verified runtime library
can be legitimate if its identity and supported semantics remain unchanged, while
the new policy bundle and generated binding receive fresh evidence.

\section{A worked amendment scenario}

Consider a synthetic amendment used only to explain the process. The original
selected rule contains a fee-discount exception. A later source adds a condition
requiring a particular relationship between the discount and the transaction.
This is not a claim that the SFC has made such an amendment; it is a controlled
example of how a change would propagate.

The source custodian retains the new document and identifies the added condition.
The inventory maps it to the earlier exception and opens an interpretation run.
Candidate A treats the condition as applying to all offers assessed after the
new commencement date. Candidate B distinguishes offers already committed before
that date under a possible transition provision. The source packet must include
that transition provision if one exists; the system cannot infer grandfathering
from business preference.

The scenario designer creates offers before and after commencement, an offer
agreed before but performed after, and an offer with incomplete timing evidence.
These examples reveal whether the competing readings differ. The formal checker
can compare their executable conditions once the timing concepts are defined.
It cannot determine which transition interpretation the source supports without
the relevant authority.

When reviewers accept the new interpretation, the data owner confirms that the
required relationship and timing facts can be supplied. Engineering builds and
verifies the new Java package. The release manifest binds its effective interval
and applicability profile. The host's release selector is tested at the exact
boundary, including the excluded endpoint of the earlier interval.

Existing event obligations require separate treatment. If a duty was opened under
the old bundle, the current MVP preserves that binding. A legal decision may
require migration, continuation or another treatment. The workbench records the
question and any approved migration specification. No background update silently
changes the source version attached to a historical duty.

The old source, interpretation, JAR and decisions remain available for replay.
A later investigation can ask what the system would have decided under the old
rule with the evidence known then, or what a corrected assessment under the new
rule would conclude. Those are different questions and their results carry
different assessment contexts.

\section{Bitemporal replay explains rather than rewrites history}

The two-time model is particularly useful in an incident. Valid time asks when a
fact or rule applies to the underlying event. Knowledge time asks what information
was available to the system at a specified point. A replay must state both.
Without this distinction, later corrections can make earlier decisions appear
inexplicably wrong or retrospectively perfect.

Suppose a transaction was assessed on Monday using a classification record
available that morning. On Wednesday, a reviewer finds evidence that changes the
classification effective from Monday. A historical replay with Monday knowledge
reproduces the original evidence state. A corrected assessment with Wednesday
knowledge can produce a different result for the same Monday transaction. Both
results are legitimate records of different questions.

The system should label the corrected assessment and link it to the original.
It should not replace the original execution hash or pretend that Wednesday's
evidence was available on Monday. This distinction supports fair investigation
of whether the failure concerned the rule, data, process or delayed information.
It also identifies which other decisions relied on the same evidence.

Replay requires the exact historical runtime or a verified equivalent path, the
bundle, fact snapshot, event history where relevant, calendars and assessment
context. Retaining only a final status and a current source URL is insufficient.
The retention policy must preserve these dependencies for the bank's required
period; this book does not invent a statutory retention duration.

Reproducing an old decision does not endorse it. A replay can faithfully show that
an earlier rule omitted an exception. That is useful evidence about the incident.
The corrected rule and reassessment should be linked, with the reason for change
and any operational remediation. Traceability supports accountability precisely
because it can preserve an error honestly.

\section{Operational responses to unknown and conflict}

Unknown and conflict are ordinary outcomes in a system built around evidence.
They should have designed workflows rather than being treated as exceptional
software crashes. The host policy identifies the responsible queue, information
request, permitted temporary action and escalation deadline for each reason.
These choices depend on the bank's process and must be approved there.

A missing classification might route to a product specialist. Conflicting fee
records might route to the data owner. An unavailable incorporated source might
require legal review before any automated decision. An unsupported RuleIR operator
requires engineering and meaning review together. The reason code helps select
the remedy, but the workflow should expose the substantive question in ordinary
language.

A conflict should preserve both records and their provenance. Automatically taking
the newest record can hide a disagreement between systems or an unauthorized
correction. If the bank adopts a precedence rule, that rule must be explicit and
reviewed for the evidence type. It should not be embedded in a generic database
query that happens to order by timestamp.

Manual overrides need their own authority and scope. An override may authorize a
specific operational action under bank policy while leaving the regulatory
interpretation unresolved. The record should say exactly what was overridden,
by whom, for which case and period, and with what rationale. It must not change
the underlying source or make future cases automatically inherit the exception.

Queues themselves can fail. Unreviewed items can accumulate, deadlines can pass,
and staff can misunderstand a blocked status. Monitoring should therefore track
age, owner, materiality and next action for unresolved issues. A low model error
rate is little comfort if the operational process routinely ignores the issues
that the model correctly escalates.

\section{Incident analysis should locate the failed boundary}

When an incorrect control or decision is discovered, the first question is where
the evidence chain failed. Was the wrong source acquired? Was a qualification
missed? Was the accepted interpretation encoded incorrectly? Did the host supply
misclassified facts? Did a transaction race invalidate the snapshot? Did an
unauthorized or stale package reach production? These mechanisms require different
repairs and different assessments of affected decisions.

The incident record should preserve the exact failing example, source versions,
candidate, approvals, JAR, input snapshot, result and operational action. It
should also preserve the strongest competing interpretation that was available
at the time. If the objection was present but suppressed, the remedy concerns
reporting and review as well as generation quality.

Containment is a bank decision. It may involve suspending a particular offer
category, returning to an approved manual process, disabling a faulty adapter or
replacing an implementation while retaining the same accepted rule. The incident
team must distinguish these options from reverting to a legally superseded rule.
The workbench supplies impact information; it does not decide the bank's external
communications or regulatory response.

Root-cause analysis should avoid a generic conclusion that ``the AI hallucinated.''
That phrase can conceal a missing source, a contaminated oracle, an undefined
classification or a budget policy that forced premature closure. A specific
mechanism supports a specific preventive control. The discovered failure becomes
a retained regression scenario, with its source and interpretation status clear.

The team should also ask what similar decisions may be affected. Dependency and
lineage queries can identify every package using the disputed definition and
every assessment using the relevant fact mapping. A database search by final
status alone will miss decisions that happened to be correct for other reasons.
The scope of remediation follows the causal dependency, not merely the visible
outcome.

\section{Availability failures must not change the meaning of a result}

The production Java evaluator should not depend on a live model or source website.
That separation lets an approved control continue to execute during an
interpretation-provider outage, subject to the bank's release and source-monitoring
policy. An outage in the workbench still matters for new controls and amendment
review, but it need not make every existing decision nondeterministic.

The host must distinguish an unavailable interpretation service from an unknown
fact and from an expired release. These conditions can require different actions.
An approved, still-applicable local package may remain usable when a research
provider is unavailable. A release whose effective status cannot be established
presents a different problem. The fallback policy must state the conditions under
which each path is permitted.

No fallback may convert a technical error into a substantive false result. If a
Java evaluation fails because the snapshot is malformed, returning false would
imply that the selected condition was evaluated and not satisfied. The correct
result is an error with an operational disposition. Likewise, a missing
interpretation report cannot be replaced by a default approval flag.

The ensemble controller has its own recovery obligations. A restart recovers
reserved and dispatched actions, preserves their charges and determines which
results are known, late or unknown. It does not restart the investigation from
zero unless an authorized successor run is created. The supervisor monitors
absolute deadlines independently of the worker's progress messages.

Capacity planning should consider the difficult tail: large dependency graphs,
many genuine alternatives and long-running external checks. Graph and action
limits protect the service, but they create explicit incomplete outcomes. The
operational owner needs to know how often those outcomes occur and whether the
manual review process can absorb them. Increasing limits without measuring cost
and review quality is not a complete capacity strategy.

\section{Model updates are method changes}

A new model version can change which alternatives are generated, how sources are
summarized and which questions are prioritized. Even if the JSON schema remains
unchanged, the interpretation method has changed. The workbench should bind
provider, model, prompt and configuration versions to each run and evaluate
material updates before making them the default.

The update review asks whether known omission cases remain detected, minority
readings remain visible, unsupported alternatives increase, and uncertainty scores
retain their measured relationship to the reference. It also checks provider
behavior relevant to budgets, such as retries, tool use and token accounting.
A faster model that silently changes these behaviors can invalidate the
controller's operational assumptions.

Existing approved Java packages do not need to be regenerated merely because the
proposal model changes. Their authority derives from the retained interpretation,
review and verification evidence. A new model may identify a concern that triggers
review, but its existence alone does not invalidate every prior decision. This
separation limits unnecessary churn while allowing genuine discoveries to enter
an accountable amendment or incident process.

Prompt changes deserve similar treatment. Adding a leading example can improve
performance on one family while anchoring all members to the same interpretation.
Removing a request for counterarguments can reduce cost while weakening diversity.
The change record should explain the intended mechanism, evaluation evidence and
remaining uncertainty. A prompt is part of the method, not incidental prose.

Search-policy changes also require evidence. Switching from a deterministic
scheduler to MCTS, changing family quotas or increasing a no-progress threshold
alters which questions receive attention. The study should compare these choices
under equal resources and independent judgments. Passing controller invariants
is necessary for either scheduler, but does not establish which investigates law
more effectively.

\section{A release readiness review with concrete evidence}

A readiness review should inspect a concrete package, not approve a future promise
to generate one. The package contains the accepted interpretation dossier, source
coverage, factual mapping, verification report, build and release manifests,
operational dispositions and unresolved-risk record. All identities must agree.
An approval of one candidate cannot be reused for a different generated bundle.

The meaning reviewer checks that the selected source perimeter and assumptions
support the rule, that plausible alternatives have a reasoned disposition, and
that any remaining uncertainty is compatible with the proposed scope. The data
owner checks input definitions, evidence quality, freshness, corrections and
calendar or valuation policies. Engineering checks exact behavior, traces,
resource limits and integration. The release owner checks authority, deployment
mechanics, monitoring and recovery.

A useful readiness table separates evidence from a decision:
\begin{longtable}{@{}P{31mm}P{60mm}P{57mm}@{}}
\caption{Evidence required for a bounded deployment decision.}\\
\toprule
Area & Evidence to inspect & Question that remains a human responsibility\\\midrule\endfirsthead
\toprule Area & Evidence to inspect & Question that remains a human responsibility\\\midrule\endhead
Meaning & Source inventory, candidates, objections, adjudication & Is this reading justified for the selected activity and time?\\
Data & Fact dictionary, provenance, validity, conflict policy & Do the bank's records establish the facts the rule needs?\\
Execution & Conformance, witnesses, mutations, Java comparison & Is the supported implementation scope adequate for the process?\\
Authority & Exact manifests, reviewer roles, signatures & Are the required people authorized and accountable for this release?\\
Operation & Host dispositions, race tests, fallback and incident rehearsal & Can staff act correctly when the result is unknown or blocked?\\
Evidence & Retained sources, histories, reproducible reports & Can a later reviewer reconstruct the decision without relying on memory?\\
\bottomrule
\end{longtable}

A missing item should produce a specific block and next action. It should not
trigger a vague demand for more confidence. If a classification is unresolved,
obtain the authority or narrow the scope. If Java traces differ, repair the
implementation. If enterprise identity is not integrated, complete that integration
before relying on local synthetic roles. The package makes these distinctions
visible enough to assign work.

A review may approve a limited deployment with explicit exclusions and monitoring
conditions. Those conditions become enforceable configuration and expiry or
review triggers. They must not remain only in meeting minutes. The system should
be able to show which operational path enforces each exclusion and which test
would fail if it were bypassed.

\section{Three operational cases for the implementation team}

The first case is a missing definition. The source packet is intact, the candidate
formula compiles, all named scenarios pass, and the ensemble agrees on the
expression. An incorporated definition is unavailable. The correct report is
blocked for the affected classification, even though the engineering checks pass.
The next action is to obtain or adjudicate the definition, not to add more random
truth-table tests.

The second case is a wrong host mapping. Reviewers accepted a component-level
rule, but the host sends one Boolean classification for an entire mixed offer.
The Java evaluator faithfully applies its inputs and returns a reproducible
result. The error lies at the subject and fact boundary. The remedy includes a
component-aware adapter, affected-case analysis and a regression scenario that
would expose a whole-offer mapping. Reproving the Boolean evaluator alone would
not address the failure.

The third case is an unresolved exception priority. Two exception guards are
true. The current RuleIR default returns conflict, while a developer proposes to
choose the first exception in the array. That change would make the program run,
but it would introduce a priority that the specification did not authorize. The
workbench must preserve conflict and ask whether the source establishes priority,
whether explicit nesting is justified, or whether the rule lies outside the
supported fragment.

These cases demonstrate why completion criteria must follow the claimed target.
In the first, source evidence is incomplete. In the second, the application mapping
is wrong relative to the accepted interpretation. In the third, a proposed
implementation change would alter semantics. The same dashboard label ``needs
review'' is not enough; the system must name the responsible boundary and the
evidence needed to repair it.

A team rehearsal should walk through each case using the actual interface and
reports. Staff should be able to locate the source, identify the affected
candidate, understand the block and record a decision without editing hidden
configuration. Rehearsal findings become product changes and acceptance tests.
A successful technical build does not substitute for this operational exercise.

\section{Maintaining a useful literature and method record}

The paper library is a working research asset. Each retained paper has a stable
filename, source URL, version or publication information, digest and reading
status. The title-and-first-author naming convention requested for this project
helps people locate a paper, while the manifest disambiguates revised preprints
and related editions. A citation should point to the version actually inspected.

The bibliography and reading notes distinguish published results, local
derivations and proposed adaptations. Catala's treatment of defaults, Stipula's
state and resource semantics, argumentation frameworks, tree search and uncertainty
methods answer different questions. Combining them in one architecture does not
transfer every theorem or empirical result to the combination. The method record
states which mechanism is borrowed and what must be re-established locally.

Official code is valuable evidence about practical assumptions. The semantic-
clustering example in Chapter~\ref{ch:search} shows why code and paper can require
separate inspection. A pinned repository revision may implement a broader
condition than the prose suggests. The correct response is to document and test
the difference, not silently choose the version that makes the proposal look
stronger.

The literature record should also preserve inaccessible or uninspected leads.
They can guide future work, but cannot support categorical implementation claims.
The present review is a substantial bounded survey of the retained sources,
not a declaration that no relevant paper exists outside the library. New work can
change the best engineering choice, especially in a rapidly developing area.

ResearchAssistant is useful for acquisition, parsing and source structure.
MathDevMCP is useful for named formal obligations. DynareMCP contributes patterns
for explicit hypotheses, falsifiers and bounded investigations. Their reuse
should remain concrete and versioned. The bank should be able to identify which
function, data structure or proof obligation is being borrowed, rather than
relying on a broad claim that another project already solved the difficult part.

\section{A practical sequence from manuscript to product}

The first step is contract implementation with public sources and scripted
members. This establishes that the ensemble can preserve alternatives, detect
specified discrepancies, respect limits and produce complete reports. The
second-circular example provides a running integration case, while synthetic
fully resolved cases exercise the positive review route. The original blocked
example should remain blocked unless actual new evidence resolves its issues.

The second step is independent reference construction and restricted live-model
evaluation. The team freezes source inventories and judgments before seeing
model outputs, runs the baseline ladder and records uncertainty. Failures are
classified by boundary and used for planned repairs. No model is promoted because
it merely reproduces the author's original 22-case oracle.

The third step is enterprise integration: identity, approved data access, isolated
source parsing, provider controls, signed or equivalently authenticated releases,
monitoring and host transaction behavior. This work can proceed alongside method
evaluation where dependencies permit. It should not be postponed until a final
demo, because operational constraints may change which designs are feasible.

The fourth step is authorized shadow operation in a narrow activity perimeter.
The study compares the workbench with the existing process and independently
adjudicates disagreements. It measures both substantive quality and reviewer
workload. The team rehearses source amendments, unknown facts, provider outages,
worker crashes and rollback decisions.

The fifth step is a bounded deployment decision based on the concrete evidence
package. The decision may approve only a subset of controls or retain manual
review for disputed categories. Monitoring then continues the evidence process:
new sources, unexpected disagreements and incidents create successor
investigations. Deployment is not the end of interpretation assurance.

This sequence does not assign an arbitrary calendar duration or staffing number.
The bank's source perimeter, access constraints and review capacity determine
those estimates. The task dependencies and acceptance evidence are concrete enough
to build a schedule once those inputs are known. Inventing a confident delivery
date without them would add apparent precision without resolving the underlying
work.

\section{What the completed prototype will and will not establish}

A completed deterministic prototype will establish that the specified records,
controller, budgets, reports and release checks behave as tested. It will produce
actual Java packages from accepted RuleIR and preserve the existing deterministic
execution contract. It will demonstrate that known omissions and persistent
ambiguities lead to the intended different outcomes in scripted scenarios.

A successful independent interpretation study can establish empirical performance
on its declared task population, with uncertainty and limitations. It may support
a bounded use or a choice among methods. It cannot prove that every future
circular will be interpreted correctly, especially under new vocabulary,
authority or distribution shift. A formal proof of the controller's release
invariant similarly depends on the recorded evidence and does not prove that no
source obligation was ever missed.

These limits do not make the product pointless. The present manual process also
contains interpretation, transcription, coding and review errors, often without a
complete record of competing readings. The proposed workbench makes those stages
explicit, introduces complementary checks and preserves the questions that remain.
Its value should be demonstrated against that real process with a fair reference
and comparable resources.

The high stakes justify both ambition and precision. The system should actively
search for ways its preferred reading could be wrong, retain those challenges,
and stop automatic work honestly when the evidence or budget is insufficient.
It should deliver code whose relationship to an approved specification can be
checked, while keeping the approval of meaning visible as a separate act.
That combination is a concrete engineering objective and a defensible basis for
building the next version of LegalMath.

\backmatter
\begingroup
\raggedright
\bibliographystyle{plainnat}
\bibliography{references}
\endgroup
\end{document}
